% !TEX program = lualatex
% ---------------------------------------------------------------------------------------------
%  Readback note 40: Elementary divisors of a Frobenius-twisted Hecke operator
%
%  Build:  python3 tools/rb.py note <slug>
%
%  One unit of a Lean artifact, rewritten as ordinary mathematics. The prose is hand-written. Lean names,
%  line ranges and counts come from data.tex; quoted Lean statements are copied from the artifact by the build.
%
%  Rules of the text: standard terminology; every symbol is explained where it first appears and listed again
%  in the notation at the end; no letter has two meanings; a proof is complete or says what it takes for granted;
%  Lean names appear only in the LEAN lines and in the concordance.
%
%  The signature rule (the build enforces it): a function, relation or predicate that the artifact defines is
%  written only through its notation macro, which the build makes from the spec (`notation`) with one slot for
%  each explicit argument of the Lean definition. No argument is fixed and dropped, also not "for brevity" and
%  not for one section. Every argument and every hypothesis of a stated declaration stands inside its
%  statement, or in a Setting that the statement cites by number; only implicit and instance arguments (the
%  ambient objects) may stand in the conventions.
% ---------------------------------------------------------------------------------------------
\documentclass[11pt,a4paper]{article}
\ifdefined\pdfvariable\pdfvariable suppressoptionalinfo \numexpr 1+2+4+8\relax\fi
\usepackage[a4paper,hmargin=35mm,top=22mm,bottom=27mm,footskip=12mm,nohead]{geometry}

\usepackage{fontspec}
\directlua{
  luaotfload.add_fallback("gfallback", {"STIXTwoMath-Regular.otf:mode=node;", "DejaVuSans.ttf:mode=node;"})
  luaotfload.add_fallback("gmonofallback", {"DejaVuSansMono.ttf:mode=node;", "STIXTwoMath-Regular.otf:mode=node;"})
}
\setmainfont{NewCM10}[Extension=.otf,UprightFont=*-Book,ItalicFont=*-BookItalic,
  BoldFont=*-Bold,BoldItalicFont=*-BoldItalic,SmallCapsFont=*-Book,SmallCapsFeatures={Letters=SmallCaps},
  RawFeature={fallback=gfallback}]
\setsansfont{NewCMSans10}[Extension=.otf,UprightFont=*-Book,ItalicFont=*-BookOblique,
  BoldFont=*-Bold,BoldItalicFont=*-BoldOblique,RawFeature={fallback=gfallback}]
\setmonofont{JetBrainsMono}[Extension=.otf,UprightFont=*-Regular,ItalicFont=*-Italic,
  BoldFont=*-SemiBold,BoldItalicFont=*-SemiBoldItalic,Scale=MatchLowercase,
  RawFeature={-calt;-liga;fallback=gmonofallback}]
\usepackage{amsmath}
\usepackage{amsthm}
\usepackage[math-style=TeX]{unicode-math}
\setmathfont{NewCMMath-Book.otf}
\emergencystretch=2.5em\relax

\usepackage{xcolor}
\usepackage{tikz}
\usetikzlibrary{calc,arrows.meta,positioning,backgrounds}
\usepackage[protrusion=true,expansion=true]{microtype}
\usepackage{enumitem}
\usepackage{needspace}
\usepackage{booktabs}
\usepackage{array}
\usepackage{longtable}
\usepackage{fancyhdr}
\usepackage{lastpage}
\usepackage{multicol}
\usepackage[hidelinks,pdfusetitle]{hyperref}

\definecolor{ink}{HTML}{0F172A}
\definecolor{inksoft}{HTML}{475569}
\definecolor{inkfaint}{HTML}{94A3B8}
\definecolor{rulegrey}{HTML}{CBD5E1}
\definecolor{arith}{HTML}{C2410C}
\definecolor{est}{HTML}{7C3AED}
\definecolor{ana}{HTML}{0891B2}
\definecolor{ok}{HTML}{047857}
\definecolor{leancol}{HTML}{3730A3}

\title{Elementary divisors of a Frobenius-twisted Hecke operator}
\author{grh-compact}
\date{}

% ------------------------------------------------------------------------------------- macros
\newcommand{\leanfont}{\ttfamily\fontsize{8.6}{10.5}\selectfont\color{leancol}}
\newcommand{\lean}[1]{\csname LN@#1\endcsname}
\newcommand{\fillin}[1]{{\color{arith}\itshape[#1]}}  % a place of the template that is not written yet; the build reports every one that is left
\newcommand{\leanquote}[1]{\csname LQ@#1\endcsname}
% generated by rb.py note; do not edit
\newcommand{\ConcDecls}{86}
\newcommand{\Pin}{72fead10}
\newcommand{\PinText}{that of commit {\leanfont 72fead10}}
\newcommand{\ArtifactName}{grh-compact}
\expandafter\newcommand\csname LN@conjSubst\endcsname{\href{https://github.com/lsernau/grh-lean/blob/72fead102f663d5c7a407964a13dc0cf9d977207/grh-compact/src/New/TwistedHecke/ClassMeans.lean\#L23-L24}{{\leanfont ClassMean.\allowbreak{}conjSubst}}}
\expandafter\newcommand\csname LN@mapC\endcsname{\href{https://github.com/lsernau/grh-lean/blob/72fead102f663d5c7a407964a13dc0cf9d977207/grh-compact/src/New/TwistedHecke/ClassMeans.lean\#L26-L28}{{\leanfont ClassMean.\allowbreak{}map\_C\_mul\_mul\_C}}}
\expandafter\newcommand\csname LN@zariski\endcsname{\href{https://github.com/lsernau/grh-lean/blob/72fead102f663d5c7a407964a13dc0cf9d977207/grh-compact/src/New/TwistedHecke/ClassMeans.lean\#L36-L50}{{\leanfont ClassMean.\allowbreak{}eq\_top\_of\_eval\_monomial}}}
\expandafter\newcommand\csname LN@diagform\endcsname{\href{https://github.com/lsernau/grh-lean/blob/72fead102f663d5c7a407964a13dc0cf9d977207/grh-compact/src/New/TwistedHecke/ClassMeans.lean\#L57-L71}{{\leanfont ClassMean.\allowbreak{}exists\_invariant\_diagonal\_form}}}
\expandafter\newcommand\csname LN@mean\endcsname{\href{https://github.com/lsernau/grh-lean/blob/72fead102f663d5c7a407964a13dc0cf9d977207/grh-compact/src/New/TwistedHecke/ClassMeans.lean\#L80-L136}{{\leanfont ClassMean.\allowbreak{}exists\_mean}}}
\expandafter\newcommand\csname LN@borel\endcsname{\href{https://github.com/lsernau/grh-lean/blob/72fead102f663d5c7a407964a13dc0cf9d977207/grh-compact/src/New/TwistedHecke/Counting.lean\#L29-L38}{{\leanfont index\_borel\_le}}}
\expandafter\newcommand\csname LN@exCone\endcsname{\href{https://github.com/lsernau/grh-lean/blob/72fead102f663d5c7a407964a13dc0cf9d977207/grh-compact/src/New/TwistedHecke/Counting.lean\#L43-L44}{{\leanfont TwistingData.\allowbreak{}excessCone}}}
\expandafter\newcommand\csname LN@dimvol\endcsname{\href{https://github.com/lsernau/grh-lean/blob/72fead102f663d5c7a407964a13dc0cf9d977207/grh-compact/src/New/TwistedHecke/Counting.lean\#L49-L115}{{\leanfont TwistingData.\allowbreak{}finrank\_excessCond\_le\_volume}}}
\expandafter\newcommand\csname LN@exdot\endcsname{\href{https://github.com/lsernau/grh-lean/blob/72fead102f663d5c7a407964a13dc0cf9d977207/grh-compact/src/New/TwistedHecke/Counting.lean\#L120-L127}{{\leanfont TwistingData.\allowbreak{}excessWeight\_dotProduct\_permStep}}}
\expandafter\newcommand\csname LN@conepoint\endcsname{\href{https://github.com/lsernau/grh-lean/blob/72fead102f663d5c7a407964a13dc0cf9d977207/grh-compact/src/New/TwistedHecke/Counting.lean\#L131-L145}{{\leanfont TwistingData.\allowbreak{}exists\_mem\_excessCone\_mem\_Icc}}}
\expandafter\newcommand\csname LN@simplex\endcsname{\href{https://github.com/lsernau/grh-lean/blob/72fead102f663d5c7a407964a13dc0cf9d977207/grh-compact/src/New/TwistedHecke/Counting.lean\#L149-L159}{{\leanfont volume\_sum\_mem\_Icc\_le}}}
\expandafter\newcommand\csname LN@HeckeSizes\endcsname{\href{https://github.com/lsernau/grh-lean/blob/72fead102f663d5c7a407964a13dc0cf9d977207/grh-compact/src/New/TwistedHecke/Counting.lean\#L163-L166}{{\leanfont HeckeSizes}}}
\expandafter\newcommand\csname LN@dimex\endcsname{\href{https://github.com/lsernau/grh-lean/blob/72fead102f663d5c7a407964a13dc0cf9d977207/grh-compact/src/New/TwistedHecke/Counting.lean\#L173-L203}{{\leanfont TwistingData.\allowbreak{}index\_mul\_finrank\_excessCond\_le}}}
\expandafter\newcommand\csname LQ@dimex\endcsname{\leanquoteline{0}{lemma TwistingData.index\_mul\_finrank\_excessCond\_le [CharP K T.p] (ha : a.1 + 1 = a') \{n ℓ h κ z\}}\leanquoteline{2}{(big : HeckeSizes T.P T.o₀ N k n a' ℓ κ h z) : (parahoric W (twoBlock k a')).index * finrank K}\leanquoteline{2}{(T.excessCond n a a' ⊥ h) ≤ (1 + κ) * Fintype.card (Monomials (Fin k) n) \textasciicircum{} N}}
\expandafter\newcommand\csname LN@dimdeep\endcsname{\href{https://github.com/lsernau/grh-lean/blob/72fead102f663d5c7a407964a13dc0cf9d977207/grh-compact/src/New/TwistedHecke/Counting.lean\#L210-L233}{{\leanfont TwistingData.\allowbreak{}sum\_finrank\_deepForms\_le}}}
\expandafter\newcommand\csname LQ@dimdeep\endcsname{\leanquoteline{0}{lemma TwistingData.sum\_finrank\_deepForms\_le [CharP K T.p] \{n r ℓ κ h z\} (big :}\leanquoteline{2}{HeckeSizes T.P T.o₀ N k n r ℓ κ h z) M : (∑ d ∈ range M, finrank K (T.deepForms k n r d) : ℝ) ≤}\leanquoteline{2}{(1 + z) * Fintype.card (Monomials (Fin k) n) \textasciicircum{} N}}
\expandafter\newcommand\csname LN@sumdeep\endcsname{\href{https://github.com/lsernau/grh-lean/blob/72fead102f663d5c7a407964a13dc0cf9d977207/grh-compact/src/New/TwistedHecke/DeepFormCost.lean\#L13-L16}{{\leanfont TwistingData.\allowbreak{}sum\_deepIndicator}}}
\expandafter\newcommand\csname LN@deepweight\endcsname{\href{https://github.com/lsernau/grh-lean/blob/72fead102f663d5c7a407964a13dc0cf9d977207/grh-compact/src/New/TwistedHecke/DeepFormCost.lean\#L21-L92}{{\leanfont TwistingData.\allowbreak{}exists\_le\_weight\_of\_mem\_deepForms}}}
\expandafter\newcommand\csname LN@cardle\endcsname{\href{https://github.com/lsernau/grh-lean/blob/72fead102f663d5c7a407964a13dc0cf9d977207/grh-compact/src/New/TwistedHecke/ElementaryDivisors.lean\#L34-L40}{{\leanfont TwistingData.\allowbreak{}card\_le\_finrank\_translatesIn}}}
\expandafter\newcommand\csname LN@pairing\endcsname{\href{https://github.com/lsernau/grh-lean/blob/72fead102f663d5c7a407964a13dc0cf9d977207/grh-compact/src/New/TwistedHecke/ElementaryDivisors.lean\#L45-L69}{{\leanfont TwistingData.\allowbreak{}of\_pairing}}}
\expandafter\newcommand\csname LN@depthfact\endcsname{\href{https://github.com/lsernau/grh-lean/blob/72fead102f663d5c7a407964a13dc0cf9d977207/grh-compact/src/New/TwistedHecke/ElementaryDivisors.lean\#L77-L137}{{\leanfont TwistingData.\allowbreak{}exists\_depth\_factorization}}}
\expandafter\newcommand\csname LN@tensordual\endcsname{\href{https://github.com/lsernau/grh-lean/blob/72fead102f663d5c7a407964a13dc0cf9d977207/grh-compact/src/New/TwistedHecke/ElementaryDivisors.lean\#L142-L146}{{\leanfont exists\_tensor\_dual}}}
\expandafter\newcommand\csname LN@projector\endcsname{\href{https://github.com/lsernau/grh-lean/blob/72fead102f663d5c7a407964a13dc0cf9d977207/grh-compact/src/New/TwistedHecke/ElementaryDivisors.lean\#L153-L208}{{\leanfont TwistingData.\allowbreak{}exists\_projector}}}
\expandafter\newcommand\csname LQ@projector\endcsname{\leanquoteline{0}{lemma exists\_projector (hp : ¬IsUnit (T.p : O)) (ha : a + 1 = r.1) \{ℓ\} θ (e : ℓ + ℓ = N)}\leanquoteline{4}{(h : N * (n * (k * (ℓ + 2))) + n + k < T.P \textasciicircum{} ℓ) : ∃ b Z, (∀ F, Z *ᵥ F ∈ T.excessCond n a r}\leanquoteline{4}{(span \{T.w \textasciicircum{} c\}) θ) ∧ (∀ y ∈ T.excessCond n a r (span \{T.w \textasciicircum{} c\}) (θ - (4 * c + 2)), ∀ F,}\leanquoteline{4}{T.w \textasciicircum{} c ∣ y ⬝ᵥ (b * Z *ᵥ F) - T.p \textasciicircum{} ((k + k \textasciicircum{} 2) * (T.o₀ * N) * N) * y ⬝ᵥ F) ∧}\leanquoteline{4}{∀ γ α β, b β * T.letterRep n (transposeGL γ) β α = b α * T.letterRep n γ α β}}
\expandafter\newcommand\csname LN@hecke\endcsname{\href{https://github.com/lsernau/grh-lean/blob/72fead102f663d5c7a407964a13dc0cf9d977207/grh-compact/src/New/TwistedHecke/ElementaryDivisors.lean\#L217-L261}{{\leanfont TwistingData.\allowbreak{}hasType\_twisted\_hecke}}}
\expandafter\newcommand\csname LQ@hecke\endcsname{\leanquoteline{0}{lemma hasType\_twisted\_hecke (hp : ¬IsUnit (T.p : O)) \{r ℓ κ y z\} (big : HeckeSizes T.P T.o₀ N k n r}\leanquoteline{2}{ℓ κ y z) \{I M : Type\} [Fintype H] [Fintype (\_ ⧸ parahoric W (twoBlock k r))] [Fintype I]}\leanquoteline{2}{[AddCommGroup M] [Module O M] (f : (X → I → O) →ₗ[O] M) \{Y : End O M\} (src :}\leanquoteline{2}{X × \_ ⧸ parahoric W (twoBlock k r) → X) (g : \_ → GL (Fin k) W) (u : \_ → GL H O) (J : I ≃ \_ × H)}\leanquoteline{2}{(hs : Injective fun e ↦ ((src e, ↑(transposeGL (g e))⁻¹) : X × \_ ⧸ parahoric W (twoBlock k r)))}\leanquoteline{2}{(hY : ∀ F, Y (f F) = f fun c i ↦ ∑ b, ∑ j, u (c, b) (J i).2 (J j).2 * (T.letterRep n (g (c, b)) *}\leanquoteline{2}{diagonal ((T.p : O) \textasciicircum{} T.depth k n r ·) * T.letterRep n b.out⁻¹ :) (J j).1 (J i).1 * F (src (c, b))}\leanquoteline{2}{j) m : HasType ↑T.p f.range m Y ((z + κ * m + y) * (Nat.card X * Nat.card I))}}
\expandafter\newcommand\csname LN@frobX\endcsname{\href{https://github.com/lsernau/grh-lean/blob/72fead102f663d5c7a407964a13dc0cf9d977207/grh-compact/src/New/TwistedHecke/FrobeniusFactors.lean\#L20-L20}{{\leanfont TwistingData.\allowbreak{}frobGenericMatrix}}}
\expandafter\newcommand\csname LN@frobOperator\endcsname{\href{https://github.com/lsernau/grh-lean/blob/72fead102f663d5c7a407964a13dc0cf9d977207/grh-compact/src/New/TwistedHecke/FrobeniusFactors.lean\#L25-L26}{{\leanfont TwistingData.\allowbreak{}frobOperator}}}
\expandafter\newcommand\csname LN@letterConj\endcsname{\href{https://github.com/lsernau/grh-lean/blob/72fead102f663d5c7a407964a13dc0cf9d977207/grh-compact/src/New/TwistedHecke/FrobeniusFactors.lean\#L30-L31}{{\leanfont TwistingData.\allowbreak{}letterConj}}}
\expandafter\newcommand\csname LN@idxw\endcsname{\href{https://github.com/lsernau/grh-lean/blob/72fead102f663d5c7a407964a13dc0cf9d977207/grh-compact/src/New/TwistedHecke/FrobeniusFactors.lean\#L35-L35}{{\leanfont TwistingData.\allowbreak{}indexWeight}}}
\expandafter\newcommand\csname LN@frobmap\endcsname{\href{https://github.com/lsernau/grh-lean/blob/72fead102f663d5c7a407964a13dc0cf9d977207/grh-compact/src/New/TwistedHecke/FrobeniusFactors.lean\#L44-L56}{{\leanfont TwistingData.\allowbreak{}frobOperator\_map}}}
\expandafter\newcommand\csname LN@frobmem\endcsname{\href{https://github.com/lsernau/grh-lean/blob/72fead102f663d5c7a407964a13dc0cf9d977207/grh-compact/src/New/TwistedHecke/FrobeniusFactors.lean\#L62-L83}{{\leanfont TwistingData.\allowbreak{}frobOperator\_mem}}}
\expandafter\newcommand\csname LN@Xmem\endcsname{\href{https://github.com/lsernau/grh-lean/blob/72fead102f663d5c7a407964a13dc0cf9d977207/grh-compact/src/New/TwistedHecke/FrobeniusFactors.lean\#L85-L86}{{\leanfont TwistingData.\allowbreak{}X\_mem\_gaussFiltration}}}
\expandafter\newcommand\csname LN@cjfrob\endcsname{\href{https://github.com/lsernau/grh-lean/blob/72fead102f663d5c7a407964a13dc0cf9d977207/grh-compact/src/New/TwistedHecke/FrobeniusFactors.lean\#L89-L96}{{\leanfont TwistingData.\allowbreak{}p\_dvd\_cj\_frobenius\_sub}}}
\expandafter\newcommand\csname LN@powdvd\endcsname{\href{https://github.com/lsernau/grh-lean/blob/72fead102f663d5c7a407964a13dc0cf9d977207/grh-compact/src/New/TwistedHecke/FrobeniusFactors.lean\#L102-L131}{{\leanfont TwistingData.\allowbreak{}pow\_dvd\_coeff\_conj\_frobOperator}}}
\expandafter\newcommand\csname LN@degle\endcsname{\href{https://github.com/lsernau/grh-lean/blob/72fead102f663d5c7a407964a13dc0cf9d977207/grh-compact/src/New/TwistedHecke/FrobeniusFactors.lean\#L133-L138}{{\leanfont TwistingData.\allowbreak{}degree\_le\_of\_coeff\_frobOperator}}}
\expandafter\newcommand\csname LN@permStep\endcsname{\href{https://github.com/lsernau/grh-lean/blob/72fead102f663d5c7a407964a13dc0cf9d977207/grh-compact/src/New/TwistedHecke/InterlacingCone.lean\#L14-L15}{{\leanfont permStep}}}
\expandafter\newcommand\csname LN@slab\endcsname{\href{https://github.com/lsernau/grh-lean/blob/72fead102f663d5c7a407964a13dc0cf9d977207/grh-compact/src/New/TwistedHecke/InterlacingCone.lean\#L18-L18}{{\leanfont slab}}}
\expandafter\newcommand\csname LN@pieri\endcsname{\href{https://github.com/lsernau/grh-lean/blob/72fead102f663d5c7a407964a13dc0cf9d977207/grh-compact/src/New/TwistedHecke/InterlacingCone.lean\#L22-L23}{{\leanfont pieriCone}}}
\expandafter\newcommand\csname LN@detinter\endcsname{\href{https://github.com/lsernau/grh-lean/blob/72fead102f663d5c7a407964a13dc0cf9d977207/grh-compact/src/New/TwistedHecke/InterlacingCone.lean\#L27-L36}{{\leanfont det\_interlacing\_indicator}}}
\expandafter\newcommand\csname LN@isBounded\endcsname{\href{https://github.com/lsernau/grh-lean/blob/72fead102f663d5c7a407964a13dc0cf9d977207/grh-compact/src/New/TwistedHecke/InterlacingCone.lean\#L40-L42}{{\leanfont isBounded\_inter\_slab}}}
\expandafter\newcommand\csname LN@signedsum\endcsname{\href{https://github.com/lsernau/grh-lean/blob/72fead102f663d5c7a407964a13dc0cf9d977207/grh-compact/src/New/TwistedHecke/InterlacingCone.lean\#L47-L106}{{\leanfont volume\_inter\_slab\_eq\_signed\_sum}}}
\expandafter\newcommand\csname LN@prodvol\endcsname{\href{https://github.com/lsernau/grh-lean/blob/72fead102f663d5c7a407964a13dc0cf9d977207/grh-compact/src/New/TwistedHecke/InterlacingCone.lean\#L111-L138}{{\leanfont prod\_mul\_volume\_inter\_slab}}}
\expandafter\newcommand\csname LQ@prodvol\endcsname{\leanquoteline{0}{lemma prod\_mul\_volume\_inter\_slab \{N\} [NeZero N] \{v P\} (hP : 1 < P) (D : ZMod N → ℝ) \{Ω\} (hΩ : ∀ ν,}\leanquoteline{2}{permStep v P 1 ν ∈ Ω ↔ ν ∈ pieriCone v P) : (∏ l : Fin k, ((P \textasciicircum{} N) \textasciicircum{} (l.1 + 1) - 1) /}\leanquoteline{2}{(P \textasciicircum{} N - 1)) * volume.real (Ω ∩ slab k D) = volume.real (slab k D)}}
\expandafter\newcommand\csname LN@stepanov\endcsname{\href{https://github.com/lsernau/grh-lean/blob/72fead102f663d5c7a407964a13dc0cf9d977207/grh-compact/src/New/TwistedHecke/RootLemma.lean\#L27-L34}{{\leanfont stepanov\_count}}}
\expandafter\newcommand\csname LN@IsRF\endcsname{\href{https://github.com/lsernau/grh-lean/blob/72fead102f663d5c7a407964a13dc0cf9d977207/grh-compact/src/New/TwistedHecke/RootLemma.lean\#L38-L38}{{\leanfont IsRealFiltration}}}
\expandafter\newcommand\csname LN@adic\endcsname{\href{https://github.com/lsernau/grh-lean/blob/72fead102f663d5c7a407964a13dc0cf9d977207/grh-compact/src/New/TwistedHecke/RootLemma.lean\#L44-L44}{{\leanfont adicFiltration}}}
\expandafter\newcommand\csname LN@gauss\endcsname{\href{https://github.com/lsernau/grh-lean/blob/72fead102f663d5c7a407964a13dc0cf9d977207/grh-compact/src/New/TwistedHecke/RootLemma.lean\#L53-L54}{{\leanfont gaussFiltration}}}
\expandafter\newcommand\csname LN@memadic\endcsname{\href{https://github.com/lsernau/grh-lean/blob/72fead102f663d5c7a407964a13dc0cf9d977207/grh-compact/src/New/TwistedHecke/RootLemma.lean\#L58-L59}{{\leanfont mem\_adicFiltration\_of\_nonpos}}}
\expandafter\newcommand\csname LN@monomialmem\endcsname{\href{https://github.com/lsernau/grh-lean/blob/72fead102f663d5c7a407964a13dc0cf9d977207/grh-compact/src/New/TwistedHecke/RootLemma.lean\#L61-L62}{{\leanfont monomial\_mem\_gaussFiltration}}}
\expandafter\newcommand\csname LN@aevalsub\endcsname{\href{https://github.com/lsernau/grh-lean/blob/72fead102f663d5c7a407964a13dc0cf9d977207/grh-compact/src/New/TwistedHecke/RootLemma.lean\#L70-L81}{{\leanfont aeval\_sub\_mem\_of\_gaussFiltration}}}
\expandafter\newcommand\csname LN@rootlemma\endcsname{\href{https://github.com/lsernau/grh-lean/blob/72fead102f663d5c7a407964a13dc0cf9d977207/grh-compact/src/New/TwistedHecke/RootLemma.lean\#L88-L154}{{\leanfont TwistingData.\allowbreak{}root\_lemma}}}
\expandafter\newcommand\csname LQ@rootlemma\endcsname{\leanquoteline{0}{lemma TwistingData.root\_lemma i (ha : a ≠ a') τ c \{f α\} (l : ℕ) (h : ∀ g ∈ levi W τ,}\leanquoteline{2}{T.polyAction g⁻¹ f ∈ gaussFiltration (adicFiltration T.w)}\leanquoteline{2}{c (l + weight c α)) (h0 : ¬T.w \textasciicircum{} (l + 1) ∣ f.coeff α) (hτ : τ a = τ a') (hc : ∀ s, c (T.rootVar a}\leanquoteline{2}{a' s) - c (T.rootVar a' a s) ∈ Set.Icc 0.4 1.6) : weight (T.excessWeight i a a') α ≤ 0}}
\expandafter\newcommand\csname LN@twoBlock\endcsname{\href{https://github.com/lsernau/grh-lean/blob/72fead102f663d5c7a407964a13dc0cf9d977207/grh-compact/src/New/TwistedHecke/TwistingData.lean\#L21-L21}{{\leanfont twoBlock}}}
\expandafter\newcommand\csname LN@TwD\endcsname{\href{https://github.com/lsernau/grh-lean/blob/72fead102f663d5c7a407964a13dc0cf9d977207/grh-compact/src/New/TwistedHecke/TwistingData.lean\#L25-L34}{{\leanfont TwistingData}}}
\expandafter\newcommand\csname LQ@TwD\endcsname{\leanquoteline{0}{structure TwistingData where}\leanquoteline{2}{(p o₀ : ℕ) ι : ZMod N → W →+* R}\leanquoteline{2}{isEven : ZMod N → \_}\leanquoteline{2}{w : R}\leanquoteline{2}{prime : p.Prime}\leanquoteline{2}{five\_le\_P : 5 ≤ p \textasciicircum{} o₀}\leanquoteline{2}{w\_sq : w \textasciicircum{} 2 = p}\leanquoteline{2}{frobenius\_congr s x : ↑p ∣ ι (s + 1) x - ι s x \textasciicircum{} p \textasciicircum{} o₀}\leanquoteline{2}{isEven\_succ s : isEven (s + 1) = !isEven s}\leanquoteline{2}{card\_residueField : .card (ResidueField W) = (p \textasciicircum{} o₀) \textasciicircum{} N}}
\expandafter\newcommand\csname LN@PT\endcsname{\href{https://github.com/lsernau/grh-lean/blob/72fead102f663d5c7a407964a13dc0cf9d977207/grh-compact/src/New/TwistedHecke/TwistingData.lean\#L42-L42}{{\leanfont TwistingData.\allowbreak{}P}}}
\expandafter\newcommand\csname LN@qT\endcsname{\href{https://github.com/lsernau/grh-lean/blob/72fead102f663d5c7a407964a13dc0cf9d977207/grh-compact/src/New/TwistedHecke/TwistingData.lean\#L45-L45}{{\leanfont TwistingData.\allowbreak{}q}}}
\expandafter\newcommand\csname LN@Tmap\endcsname{\href{https://github.com/lsernau/grh-lean/blob/72fead102f663d5c7a407964a13dc0cf9d977207/grh-compact/src/New/TwistedHecke/TwistingData.lean\#L49-L51}{{\leanfont TwistingData.\allowbreak{}map}}}
\expandafter\newcommand\csname LN@letterMatrix\endcsname{\href{https://github.com/lsernau/grh-lean/blob/72fead102f663d5c7a407964a13dc0cf9d977207/grh-compact/src/New/TwistedHecke/TwistingData.lean\#L55-L56}{{\leanfont TwistingData.\allowbreak{}letterMatrix}}}
\expandafter\newcommand\csname LN@polyAction\endcsname{\href{https://github.com/lsernau/grh-lean/blob/72fead102f663d5c7a407964a13dc0cf9d977207/grh-compact/src/New/TwistedHecke/TwistingData.lean\#L60-L61}{{\leanfont TwistingData.\allowbreak{}polyAction}}}
\expandafter\newcommand\csname LN@letterRep\endcsname{\href{https://github.com/lsernau/grh-lean/blob/72fead102f663d5c7a407964a13dc0cf9d977207/grh-compact/src/New/TwistedHecke/TwistingData.lean\#L65-L66}{{\leanfont TwistingData.\allowbreak{}letterRep}}}
\expandafter\newcommand\csname LN@dpi\endcsname{\href{https://github.com/lsernau/grh-lean/blob/72fead102f663d5c7a407964a13dc0cf9d977207/grh-compact/src/New/TwistedHecke/TwistingData.lean\#L70-L70}{{\leanfont TwistingData.\allowbreak{}deepIndicator}}}
\expandafter\newcommand\csname LN@depth\endcsname{\href{https://github.com/lsernau/grh-lean/blob/72fead102f663d5c7a407964a13dc0cf9d977207/grh-compact/src/New/TwistedHecke/TwistingData.lean\#L74-L74}{{\leanfont TwistingData.\allowbreak{}depth}}}
\expandafter\newcommand\csname LN@rootVar\endcsname{\href{https://github.com/lsernau/grh-lean/blob/72fead102f663d5c7a407964a13dc0cf9d977207/grh-compact/src/New/TwistedHecke/TwistingData.lean\#L77-L77}{{\leanfont TwistingData.\allowbreak{}rootVar}}}
\expandafter\newcommand\csname LN@exw\endcsname{\href{https://github.com/lsernau/grh-lean/blob/72fead102f663d5c7a407964a13dc0cf9d977207/grh-compact/src/New/TwistedHecke/TwistingData.lean\#L81-L82}{{\leanfont TwistingData.\allowbreak{}excessWeight}}}
\expandafter\newcommand\csname LN@translatesIn\endcsname{\href{https://github.com/lsernau/grh-lean/blob/72fead102f663d5c7a407964a13dc0cf9d977207/grh-compact/src/New/TwistedHecke/TwistingData.lean\#L86-L87}{{\leanfont TwistingData.\allowbreak{}translatesIn}}}
\expandafter\newcommand\csname LN@excessCond\endcsname{\href{https://github.com/lsernau/grh-lean/blob/72fead102f663d5c7a407964a13dc0cf9d977207/grh-compact/src/New/TwistedHecke/TwistingData.lean\#L91-L92}{{\leanfont TwistingData.\allowbreak{}excessCond}}}
\expandafter\newcommand\csname LN@deepForms\endcsname{\href{https://github.com/lsernau/grh-lean/blob/72fead102f663d5c7a407964a13dc0cf9d977207/grh-compact/src/New/TwistedHecke/TwistingData.lean\#L96-L96}{{\leanfont TwistingData.\allowbreak{}deepForms}}}
\expandafter\newcommand\csname LN@oneltP\endcsname{\href{https://github.com/lsernau/grh-lean/blob/72fead102f663d5c7a407964a13dc0cf9d977207/grh-compact/src/New/TwistedHecke/TwistingData.lean\#L103-L103}{{\leanfont TwistingData.\allowbreak{}one\_lt\_P}}}
\expandafter\newcommand\csname LN@oneltq\endcsname{\href{https://github.com/lsernau/grh-lean/blob/72fead102f663d5c7a407964a13dc0cf9d977207/grh-compact/src/New/TwistedHecke/TwistingData.lean\#L105-L105}{{\leanfont TwistingData.\allowbreak{}one\_lt\_q}}}
\expandafter\newcommand\csname LN@frobpow\endcsname{\href{https://github.com/lsernau/grh-lean/blob/72fead102f663d5c7a407964a13dc0cf9d977207/grh-compact/src/New/TwistedHecke/TwistingData.lean\#L108-L109}{{\leanfont TwistingData.\allowbreak{}frobenius\_pow\_congr}}}
\expandafter\newcommand\csname LN@wnezero\endcsname{\href{https://github.com/lsernau/grh-lean/blob/72fead102f663d5c7a407964a13dc0cf9d977207/grh-compact/src/New/TwistedHecke/TwistingData.lean\#L114-L114}{{\leanfont TwistingData.\allowbreak{}w\_ne\_zero}}}
\expandafter\newcommand\csname LN@finres\endcsname{\href{https://github.com/lsernau/grh-lean/blob/72fead102f663d5c7a407964a13dc0cf9d977207/grh-compact/src/New/TwistedHecke/TwistingData.lean\#L117-L117}{{\leanfont TwistingData.\allowbreak{}finite\_residueField}}}
\expandafter\newcommand\csname LN@pdvdpar\endcsname{\href{https://github.com/lsernau/grh-lean/blob/72fead102f663d5c7a407964a13dc0cf9d977207/grh-compact/src/New/TwistedHecke/TwistingData.lean\#L119-L124}{{\leanfont TwistingData.\allowbreak{}p\_dvd\_letterMatrix\_of\_mem\_parahoric}}}
\expandafter\newcommand\csname LN@polyActionX\endcsname{\href{https://github.com/lsernau/grh-lean/blob/72fead102f663d5c7a407964a13dc0cf9d977207/grh-compact/src/New/TwistedHecke/TwistingData.lean\#L128-L129}{{\leanfont TwistingData.\allowbreak{}polyAction\_X}}}
\expandafter\newcommand\csname LN@diagdepth\endcsname{\href{https://github.com/lsernau/grh-lean/blob/72fead102f663d5c7a407964a13dc0cf9d977207/grh-compact/src/New/TwistedHecke/TwistingData.lean\#L131-L133}{{\leanfont TwistingData.\allowbreak{}diagonal\_pow\_depth}}}
\expandafter\newcommand\csname LN@polyOfCoeffs\endcsname{\href{https://github.com/lsernau/grh-lean/blob/72fead102f663d5c7a407964a13dc0cf9d977207/grh-compact/src/New/TwistedHecke/TwistingData.lean\#L137-L140}{{\leanfont TwistingData.\allowbreak{}polyAction\_polyOfCoeffs}}}
\expandafter\newcommand\csname LN@letterRepmap\endcsname{\href{https://github.com/lsernau/grh-lean/blob/72fead102f663d5c7a407964a13dc0cf9d977207/grh-compact/src/New/TwistedHecke/TwistingData.lean\#L142-L142}{{\leanfont TwistingData.\allowbreak{}letterRep\_map}}}
\expandafter\newcommand\csname LN@oneaddval\endcsname{\href{https://github.com/lsernau/grh-lean/blob/72fead102f663d5c7a407964a13dc0cf9d977207/grh-compact/src/New/TwistedHecke/TwistingData.lean\#L144-L145}{{\leanfont TwistingData.\allowbreak{}one\_add\_val\_sub\_one\_sub}}}
\expandafter\newcommand\csname LN@fivemul\endcsname{\href{https://github.com/lsernau/grh-lean/blob/72fead102f663d5c7a407964a13dc0cf9d977207/grh-compact/src/New/TwistedHecke/TwistingData.lean\#L147-L148}{{\leanfont TwistingData.\allowbreak{}five\_mul\_P\_pow\_le\_q}}}
\expandafter\newcommand\csname LN@exwIcc\endcsname{\href{https://github.com/lsernau/grh-lean/blob/72fead102f663d5c7a407964a13dc0cf9d977207/grh-compact/src/New/TwistedHecke/TwistingData.lean\#L150-L152}{{\leanfont TwistingData.\allowbreak{}excessWeight\_mem\_Icc}}}
\expandafter\newcommand\csname LN@IndexDvd\endcsname{\href{https://github.com/lsernau/grh-lean/blob/72fead102f663d5c7a407964a13dc0cf9d977207/grh-compact/src/New/Core/ModuleIndex.lean\#L20-L20}{{\leanfont IndexDvd}}}
\expandafter\newcommand\csname LN@HasType\endcsname{\href{https://github.com/lsernau/grh-lean/blob/72fead102f663d5c7a407964a13dc0cf9d977207/grh-compact/src/New/Core/ModuleIndex.lean\#L55-L56}{{\leanfont HasType}}}
\expandafter\newcommand\csname LN@ofdetdvd\endcsname{\href{https://github.com/lsernau/grh-lean/blob/72fead102f663d5c7a407964a13dc0cf9d977207/grh-compact/src/New/Core/ModuleIndex.lean\#L68-L73}{{\leanfont IndexDvd.\allowbreak{}of\_det\_dvd}}}
\expandafter\newcommand\csname LN@idtrans\endcsname{\href{https://github.com/lsernau/grh-lean/blob/72fead102f663d5c7a407964a13dc0cf9d977207/grh-compact/src/New/Core/ModuleIndex.lean\#L31-L33}{{\leanfont IndexDvd.\allowbreak{}trans}}}
\expandafter\newcommand\csname LN@idmap\endcsname{\href{https://github.com/lsernau/grh-lean/blob/72fead102f663d5c7a407964a13dc0cf9d977207/grh-compact/src/New/Core/ModuleIndex.lean\#L37-L41}{{\leanfont IndexDvd.\allowbreak{}map}}}
\expandafter\newcommand\csname LN@idcomp\endcsname{\href{https://github.com/lsernau/grh-lean/blob/72fead102f663d5c7a407964a13dc0cf9d977207/grh-compact/src/New/Core/ModuleIndex.lean\#L50-L51}{{\leanfont IndexDvd.\allowbreak{}comp}}}
\expandafter\newcommand\csname LN@hermite\endcsname{\href{https://github.com/lsernau/grh-lean/blob/72fead102f663d5c7a407964a13dc0cf9d977207/grh-compact/src/Lib/Adelic/Witt/FrameAction.lean\#L58-L69}{{\leanfont IsValuationPair.\allowbreak{}hermite}}}
\expandafter\newcommand\csname LN@smith\endcsname{\href{https://github.com/lsernau/grh-lean/blob/72fead102f663d5c7a407964a13dc0cf9d977207/grh-compact/src/Lib/General/LinearAlgebra/ColumnLattices.lean\#L23-L44}{{\leanfont IsValuationPair.\allowbreak{}smith\_normal\_form}}}
\expandafter\newcommand\csname LN@valtop\endcsname{\href{https://github.com/lsernau/grh-lean/blob/72fead102f663d5c7a407964a13dc0cf9d977207/grh-compact/src/Lib/Adelic/Witt/FrameAction.lean\#L50-L51}{{\leanfont isValuationPair\_top}}}
\expandafter\newcommand\csname LN@nullst\endcsname{\href{https://github.com/lsernau/grh-lean/blob/72fead102f663d5c7a407964a13dc0cf9d977207/grh-compact/src/Lib/General/Algebra/ZerosCostDegree.lean\#L271-L313}{{\leanfont exists\_le\_cost\_of\_vanishing}}}
\expandafter\newcommand\csname LN@cardGL\endcsname{\href{https://github.com/lsernau/grh-lean/blob/72fead102f663d5c7a407964a13dc0cf9d977207/grh-compact/src/Lib/LocalGL/Green/CuspidalAtScalars.lean\#L17-L21}{{\leanfont card\_GL\_eq\_prod}}}
\newcommand{\ConcordanceRows}{\crow{conjSubst}{\href{https://github.com/lsernau/grh-lean/blob/72fead102f663d5c7a407964a13dc0cf9d977207/grh-compact/src/New/TwistedHecke/ClassMeans.lean\#L23-L24}{{\leanfont ClassMean.\allowbreak{}conjSubst}}}{ClassMeans.lean}{23--24}
\crow{mapC}{\href{https://github.com/lsernau/grh-lean/blob/72fead102f663d5c7a407964a13dc0cf9d977207/grh-compact/src/New/TwistedHecke/ClassMeans.lean\#L26-L28}{{\leanfont ClassMean.\allowbreak{}map\_C\_mul\_mul\_C}}}{ClassMeans.lean}{26--28}
\crow{zariski}{\href{https://github.com/lsernau/grh-lean/blob/72fead102f663d5c7a407964a13dc0cf9d977207/grh-compact/src/New/TwistedHecke/ClassMeans.lean\#L36-L50}{{\leanfont ClassMean.\allowbreak{}eq\_top\_of\_eval\_monomial}}}{ClassMeans.lean}{36--50}
\crow{diagform}{\href{https://github.com/lsernau/grh-lean/blob/72fead102f663d5c7a407964a13dc0cf9d977207/grh-compact/src/New/TwistedHecke/ClassMeans.lean\#L57-L71}{{\leanfont ClassMean.\allowbreak{}exists\_invariant\_diagonal\_form}}}{ClassMeans.lean}{57--71}
\crow{mean}{\href{https://github.com/lsernau/grh-lean/blob/72fead102f663d5c7a407964a13dc0cf9d977207/grh-compact/src/New/TwistedHecke/ClassMeans.lean\#L80-L136}{{\leanfont ClassMean.\allowbreak{}exists\_mean}}}{ClassMeans.lean}{80--136}
\crow{borel}{\href{https://github.com/lsernau/grh-lean/blob/72fead102f663d5c7a407964a13dc0cf9d977207/grh-compact/src/New/TwistedHecke/Counting.lean\#L29-L38}{{\leanfont index\_borel\_le}}}{Counting.lean}{29--38}
\crow{exCone}{\href{https://github.com/lsernau/grh-lean/blob/72fead102f663d5c7a407964a13dc0cf9d977207/grh-compact/src/New/TwistedHecke/Counting.lean\#L43-L44}{{\leanfont TwistingData.\allowbreak{}excessCone}}}{Counting.lean}{43--44}
\crow{dimvol}{\href{https://github.com/lsernau/grh-lean/blob/72fead102f663d5c7a407964a13dc0cf9d977207/grh-compact/src/New/TwistedHecke/Counting.lean\#L49-L115}{{\leanfont TwistingData.\allowbreak{}finrank\_excessCond\_le\_volume}}}{Counting.lean}{49--115}
\crow{exdot}{\href{https://github.com/lsernau/grh-lean/blob/72fead102f663d5c7a407964a13dc0cf9d977207/grh-compact/src/New/TwistedHecke/Counting.lean\#L120-L127}{{\leanfont TwistingData.\allowbreak{}excessWeight\_dotProduct\_permStep}}}{Counting.lean}{120--127}
\crow{conepoint}{\href{https://github.com/lsernau/grh-lean/blob/72fead102f663d5c7a407964a13dc0cf9d977207/grh-compact/src/New/TwistedHecke/Counting.lean\#L131-L145}{{\leanfont TwistingData.\allowbreak{}exists\_mem\_excessCone\_mem\_Icc}}}{Counting.lean}{131--145}
\crow{simplex}{\href{https://github.com/lsernau/grh-lean/blob/72fead102f663d5c7a407964a13dc0cf9d977207/grh-compact/src/New/TwistedHecke/Counting.lean\#L149-L159}{{\leanfont volume\_sum\_mem\_Icc\_le}}}{Counting.lean}{149--159}
\crow{HeckeSizes}{\href{https://github.com/lsernau/grh-lean/blob/72fead102f663d5c7a407964a13dc0cf9d977207/grh-compact/src/New/TwistedHecke/Counting.lean\#L163-L166}{{\leanfont HeckeSizes}}}{Counting.lean}{163--166}
\crow{dimex}{\href{https://github.com/lsernau/grh-lean/blob/72fead102f663d5c7a407964a13dc0cf9d977207/grh-compact/src/New/TwistedHecke/Counting.lean\#L173-L203}{{\leanfont TwistingData.\allowbreak{}index\_mul\_finrank\_excessCond\_le}}}{Counting.lean}{173--203}
\crow{dimdeep}{\href{https://github.com/lsernau/grh-lean/blob/72fead102f663d5c7a407964a13dc0cf9d977207/grh-compact/src/New/TwistedHecke/Counting.lean\#L210-L233}{{\leanfont TwistingData.\allowbreak{}sum\_finrank\_deepForms\_le}}}{Counting.lean}{210--233}
\crow{sumdeep}{\href{https://github.com/lsernau/grh-lean/blob/72fead102f663d5c7a407964a13dc0cf9d977207/grh-compact/src/New/TwistedHecke/DeepFormCost.lean\#L13-L16}{{\leanfont TwistingData.\allowbreak{}sum\_deepIndicator}}}{DeepFormCost.lean}{13--16}
\crow{deepweight}{\href{https://github.com/lsernau/grh-lean/blob/72fead102f663d5c7a407964a13dc0cf9d977207/grh-compact/src/New/TwistedHecke/DeepFormCost.lean\#L21-L92}{{\leanfont TwistingData.\allowbreak{}exists\_le\_weight\_of\_mem\_deepForms}}}{DeepFormCost.lean}{21--92}
\crow{cardle}{\href{https://github.com/lsernau/grh-lean/blob/72fead102f663d5c7a407964a13dc0cf9d977207/grh-compact/src/New/TwistedHecke/ElementaryDivisors.lean\#L34-L40}{{\leanfont TwistingData.\allowbreak{}card\_le\_finrank\_translatesIn}}}{ElementaryDivisors.lean}{34--40}
\crow{pairing}{\href{https://github.com/lsernau/grh-lean/blob/72fead102f663d5c7a407964a13dc0cf9d977207/grh-compact/src/New/TwistedHecke/ElementaryDivisors.lean\#L45-L69}{{\leanfont TwistingData.\allowbreak{}of\_pairing}}}{ElementaryDivisors.lean}{45--69}
\crow{depthfact}{\href{https://github.com/lsernau/grh-lean/blob/72fead102f663d5c7a407964a13dc0cf9d977207/grh-compact/src/New/TwistedHecke/ElementaryDivisors.lean\#L77-L137}{{\leanfont TwistingData.\allowbreak{}exists\_depth\_factorization}}}{ElementaryDivisors.lean}{77--137}
\crow{tensordual}{\href{https://github.com/lsernau/grh-lean/blob/72fead102f663d5c7a407964a13dc0cf9d977207/grh-compact/src/New/TwistedHecke/ElementaryDivisors.lean\#L142-L146}{{\leanfont exists\_tensor\_dual}}}{ElementaryDivisors.lean}{142--146}
\crow{projector}{\href{https://github.com/lsernau/grh-lean/blob/72fead102f663d5c7a407964a13dc0cf9d977207/grh-compact/src/New/TwistedHecke/ElementaryDivisors.lean\#L153-L208}{{\leanfont TwistingData.\allowbreak{}exists\_projector}}}{ElementaryDivisors.lean}{153--208}
\crow{hecke}{\href{https://github.com/lsernau/grh-lean/blob/72fead102f663d5c7a407964a13dc0cf9d977207/grh-compact/src/New/TwistedHecke/ElementaryDivisors.lean\#L217-L261}{{\leanfont TwistingData.\allowbreak{}hasType\_twisted\_hecke}}}{ElementaryDivisors.lean}{217--261}
\crow{frobX}{\href{https://github.com/lsernau/grh-lean/blob/72fead102f663d5c7a407964a13dc0cf9d977207/grh-compact/src/New/TwistedHecke/FrobeniusFactors.lean\#L20-L20}{{\leanfont TwistingData.\allowbreak{}frobGenericMatrix}}}{FrobeniusFactors.lean}{20--20}
\crow{frobOperator}{\href{https://github.com/lsernau/grh-lean/blob/72fead102f663d5c7a407964a13dc0cf9d977207/grh-compact/src/New/TwistedHecke/FrobeniusFactors.lean\#L25-L26}{{\leanfont TwistingData.\allowbreak{}frobOperator}}}{FrobeniusFactors.lean}{25--26}
\crow{letterConj}{\href{https://github.com/lsernau/grh-lean/blob/72fead102f663d5c7a407964a13dc0cf9d977207/grh-compact/src/New/TwistedHecke/FrobeniusFactors.lean\#L30-L31}{{\leanfont TwistingData.\allowbreak{}letterConj}}}{FrobeniusFactors.lean}{30--31}
\crow{idxw}{\href{https://github.com/lsernau/grh-lean/blob/72fead102f663d5c7a407964a13dc0cf9d977207/grh-compact/src/New/TwistedHecke/FrobeniusFactors.lean\#L35-L35}{{\leanfont TwistingData.\allowbreak{}indexWeight}}}{FrobeniusFactors.lean}{35--35}
\crow{frobmap}{\href{https://github.com/lsernau/grh-lean/blob/72fead102f663d5c7a407964a13dc0cf9d977207/grh-compact/src/New/TwistedHecke/FrobeniusFactors.lean\#L44-L56}{{\leanfont TwistingData.\allowbreak{}frobOperator\_map}}}{FrobeniusFactors.lean}{44--56}
\crow{frobmem}{\href{https://github.com/lsernau/grh-lean/blob/72fead102f663d5c7a407964a13dc0cf9d977207/grh-compact/src/New/TwistedHecke/FrobeniusFactors.lean\#L62-L83}{{\leanfont TwistingData.\allowbreak{}frobOperator\_mem}}}{FrobeniusFactors.lean}{62--83}
\crow{Xmem}{\href{https://github.com/lsernau/grh-lean/blob/72fead102f663d5c7a407964a13dc0cf9d977207/grh-compact/src/New/TwistedHecke/FrobeniusFactors.lean\#L85-L86}{{\leanfont TwistingData.\allowbreak{}X\_mem\_gaussFiltration}}}{FrobeniusFactors.lean}{85--86}
\crow{cjfrob}{\href{https://github.com/lsernau/grh-lean/blob/72fead102f663d5c7a407964a13dc0cf9d977207/grh-compact/src/New/TwistedHecke/FrobeniusFactors.lean\#L89-L96}{{\leanfont TwistingData.\allowbreak{}p\_dvd\_cj\_frobenius\_sub}}}{FrobeniusFactors.lean}{89--96}
\crow{powdvd}{\href{https://github.com/lsernau/grh-lean/blob/72fead102f663d5c7a407964a13dc0cf9d977207/grh-compact/src/New/TwistedHecke/FrobeniusFactors.lean\#L102-L131}{{\leanfont TwistingData.\allowbreak{}pow\_dvd\_coeff\_conj\_frobOperator}}}{FrobeniusFactors.lean}{102--131}
\crow{degle}{\href{https://github.com/lsernau/grh-lean/blob/72fead102f663d5c7a407964a13dc0cf9d977207/grh-compact/src/New/TwistedHecke/FrobeniusFactors.lean\#L133-L138}{{\leanfont TwistingData.\allowbreak{}degree\_le\_of\_coeff\_frobOperator}}}{FrobeniusFactors.lean}{133--138}
\crow{permStep}{\href{https://github.com/lsernau/grh-lean/blob/72fead102f663d5c7a407964a13dc0cf9d977207/grh-compact/src/New/TwistedHecke/InterlacingCone.lean\#L14-L15}{{\leanfont permStep}}}{InterlacingCone.lean}{14--15}
\crow{slab}{\href{https://github.com/lsernau/grh-lean/blob/72fead102f663d5c7a407964a13dc0cf9d977207/grh-compact/src/New/TwistedHecke/InterlacingCone.lean\#L18-L18}{{\leanfont slab}}}{InterlacingCone.lean}{18--18}
\crow{pieri}{\href{https://github.com/lsernau/grh-lean/blob/72fead102f663d5c7a407964a13dc0cf9d977207/grh-compact/src/New/TwistedHecke/InterlacingCone.lean\#L22-L23}{{\leanfont pieriCone}}}{InterlacingCone.lean}{22--23}
\crow{detinter}{\href{https://github.com/lsernau/grh-lean/blob/72fead102f663d5c7a407964a13dc0cf9d977207/grh-compact/src/New/TwistedHecke/InterlacingCone.lean\#L27-L36}{{\leanfont det\_interlacing\_indicator}}}{InterlacingCone.lean}{27--36}
\crow{isBounded}{\href{https://github.com/lsernau/grh-lean/blob/72fead102f663d5c7a407964a13dc0cf9d977207/grh-compact/src/New/TwistedHecke/InterlacingCone.lean\#L40-L42}{{\leanfont isBounded\_inter\_slab}}}{InterlacingCone.lean}{40--42}
\crow{signedsum}{\href{https://github.com/lsernau/grh-lean/blob/72fead102f663d5c7a407964a13dc0cf9d977207/grh-compact/src/New/TwistedHecke/InterlacingCone.lean\#L47-L106}{{\leanfont volume\_inter\_slab\_eq\_signed\_sum}}}{InterlacingCone.lean}{47--106}
\crow{prodvol}{\href{https://github.com/lsernau/grh-lean/blob/72fead102f663d5c7a407964a13dc0cf9d977207/grh-compact/src/New/TwistedHecke/InterlacingCone.lean\#L111-L138}{{\leanfont prod\_mul\_volume\_inter\_slab}}}{InterlacingCone.lean}{111--138}
\crow{stepanov}{\href{https://github.com/lsernau/grh-lean/blob/72fead102f663d5c7a407964a13dc0cf9d977207/grh-compact/src/New/TwistedHecke/RootLemma.lean\#L27-L34}{{\leanfont stepanov\_count}}}{RootLemma.lean}{27--34}
\crow{IsRF}{\href{https://github.com/lsernau/grh-lean/blob/72fead102f663d5c7a407964a13dc0cf9d977207/grh-compact/src/New/TwistedHecke/RootLemma.lean\#L38-L38}{{\leanfont IsRealFiltration}}}{RootLemma.lean}{38--38}
\crow{adic}{\href{https://github.com/lsernau/grh-lean/blob/72fead102f663d5c7a407964a13dc0cf9d977207/grh-compact/src/New/TwistedHecke/RootLemma.lean\#L44-L44}{{\leanfont adicFiltration}}}{RootLemma.lean}{44--44}
\crow{gauss}{\href{https://github.com/lsernau/grh-lean/blob/72fead102f663d5c7a407964a13dc0cf9d977207/grh-compact/src/New/TwistedHecke/RootLemma.lean\#L53-L54}{{\leanfont gaussFiltration}}}{RootLemma.lean}{53--54}
\crow{memadic}{\href{https://github.com/lsernau/grh-lean/blob/72fead102f663d5c7a407964a13dc0cf9d977207/grh-compact/src/New/TwistedHecke/RootLemma.lean\#L58-L59}{{\leanfont mem\_adicFiltration\_of\_nonpos}}}{RootLemma.lean}{58--59}
\crow{monomialmem}{\href{https://github.com/lsernau/grh-lean/blob/72fead102f663d5c7a407964a13dc0cf9d977207/grh-compact/src/New/TwistedHecke/RootLemma.lean\#L61-L62}{{\leanfont monomial\_mem\_gaussFiltration}}}{RootLemma.lean}{61--62}
\crow{aevalsub}{\href{https://github.com/lsernau/grh-lean/blob/72fead102f663d5c7a407964a13dc0cf9d977207/grh-compact/src/New/TwistedHecke/RootLemma.lean\#L70-L81}{{\leanfont aeval\_sub\_mem\_of\_gaussFiltration}}}{RootLemma.lean}{70--81}
\crow{rootlemma}{\href{https://github.com/lsernau/grh-lean/blob/72fead102f663d5c7a407964a13dc0cf9d977207/grh-compact/src/New/TwistedHecke/RootLemma.lean\#L88-L154}{{\leanfont TwistingData.\allowbreak{}root\_lemma}}}{RootLemma.lean}{88--154}
\crow{twoBlock}{\href{https://github.com/lsernau/grh-lean/blob/72fead102f663d5c7a407964a13dc0cf9d977207/grh-compact/src/New/TwistedHecke/TwistingData.lean\#L21-L21}{{\leanfont twoBlock}}}{TwistingData.lean}{21--21}
\crow{TwD}{\href{https://github.com/lsernau/grh-lean/blob/72fead102f663d5c7a407964a13dc0cf9d977207/grh-compact/src/New/TwistedHecke/TwistingData.lean\#L25-L34}{{\leanfont TwistingData}}}{TwistingData.lean}{25--34}
\crow{PT}{\href{https://github.com/lsernau/grh-lean/blob/72fead102f663d5c7a407964a13dc0cf9d977207/grh-compact/src/New/TwistedHecke/TwistingData.lean\#L42-L42}{{\leanfont TwistingData.\allowbreak{}P}}}{TwistingData.lean}{42--42}
\crow{qT}{\href{https://github.com/lsernau/grh-lean/blob/72fead102f663d5c7a407964a13dc0cf9d977207/grh-compact/src/New/TwistedHecke/TwistingData.lean\#L45-L45}{{\leanfont TwistingData.\allowbreak{}q}}}{TwistingData.lean}{45--45}
\crow{Tmap}{\href{https://github.com/lsernau/grh-lean/blob/72fead102f663d5c7a407964a13dc0cf9d977207/grh-compact/src/New/TwistedHecke/TwistingData.lean\#L49-L51}{{\leanfont TwistingData.\allowbreak{}map}}}{TwistingData.lean}{49--51}
\crow{letterMatrix}{\href{https://github.com/lsernau/grh-lean/blob/72fead102f663d5c7a407964a13dc0cf9d977207/grh-compact/src/New/TwistedHecke/TwistingData.lean\#L55-L56}{{\leanfont TwistingData.\allowbreak{}letterMatrix}}}{TwistingData.lean}{55--56}
\crow{polyAction}{\href{https://github.com/lsernau/grh-lean/blob/72fead102f663d5c7a407964a13dc0cf9d977207/grh-compact/src/New/TwistedHecke/TwistingData.lean\#L60-L61}{{\leanfont TwistingData.\allowbreak{}polyAction}}}{TwistingData.lean}{60--61}
\crow{letterRep}{\href{https://github.com/lsernau/grh-lean/blob/72fead102f663d5c7a407964a13dc0cf9d977207/grh-compact/src/New/TwistedHecke/TwistingData.lean\#L65-L66}{{\leanfont TwistingData.\allowbreak{}letterRep}}}{TwistingData.lean}{65--66}
\crow{dpi}{\href{https://github.com/lsernau/grh-lean/blob/72fead102f663d5c7a407964a13dc0cf9d977207/grh-compact/src/New/TwistedHecke/TwistingData.lean\#L70-L70}{{\leanfont TwistingData.\allowbreak{}deepIndicator}}}{TwistingData.lean}{70--70}
\crow{depth}{\href{https://github.com/lsernau/grh-lean/blob/72fead102f663d5c7a407964a13dc0cf9d977207/grh-compact/src/New/TwistedHecke/TwistingData.lean\#L74-L74}{{\leanfont TwistingData.\allowbreak{}depth}}}{TwistingData.lean}{74--74}
\crow{rootVar}{\href{https://github.com/lsernau/grh-lean/blob/72fead102f663d5c7a407964a13dc0cf9d977207/grh-compact/src/New/TwistedHecke/TwistingData.lean\#L77-L77}{{\leanfont TwistingData.\allowbreak{}rootVar}}}{TwistingData.lean}{77--77}
\crow{exw}{\href{https://github.com/lsernau/grh-lean/blob/72fead102f663d5c7a407964a13dc0cf9d977207/grh-compact/src/New/TwistedHecke/TwistingData.lean\#L81-L82}{{\leanfont TwistingData.\allowbreak{}excessWeight}}}{TwistingData.lean}{81--82}
\crow{translatesIn}{\href{https://github.com/lsernau/grh-lean/blob/72fead102f663d5c7a407964a13dc0cf9d977207/grh-compact/src/New/TwistedHecke/TwistingData.lean\#L86-L87}{{\leanfont TwistingData.\allowbreak{}translatesIn}}}{TwistingData.lean}{86--87}
\crow{excessCond}{\href{https://github.com/lsernau/grh-lean/blob/72fead102f663d5c7a407964a13dc0cf9d977207/grh-compact/src/New/TwistedHecke/TwistingData.lean\#L91-L92}{{\leanfont TwistingData.\allowbreak{}excessCond}}}{TwistingData.lean}{91--92}
\crow{deepForms}{\href{https://github.com/lsernau/grh-lean/blob/72fead102f663d5c7a407964a13dc0cf9d977207/grh-compact/src/New/TwistedHecke/TwistingData.lean\#L96-L96}{{\leanfont TwistingData.\allowbreak{}deepForms}}}{TwistingData.lean}{96--96}
\crow{oneltP}{\href{https://github.com/lsernau/grh-lean/blob/72fead102f663d5c7a407964a13dc0cf9d977207/grh-compact/src/New/TwistedHecke/TwistingData.lean\#L103-L103}{{\leanfont TwistingData.\allowbreak{}one\_lt\_P}}}{TwistingData.lean}{103--103}
\crow{oneltq}{\href{https://github.com/lsernau/grh-lean/blob/72fead102f663d5c7a407964a13dc0cf9d977207/grh-compact/src/New/TwistedHecke/TwistingData.lean\#L105-L105}{{\leanfont TwistingData.\allowbreak{}one\_lt\_q}}}{TwistingData.lean}{105--105}
\crow{frobpow}{\href{https://github.com/lsernau/grh-lean/blob/72fead102f663d5c7a407964a13dc0cf9d977207/grh-compact/src/New/TwistedHecke/TwistingData.lean\#L108-L109}{{\leanfont TwistingData.\allowbreak{}frobenius\_pow\_congr}}}{TwistingData.lean}{108--109}
\crow{wnezero}{\href{https://github.com/lsernau/grh-lean/blob/72fead102f663d5c7a407964a13dc0cf9d977207/grh-compact/src/New/TwistedHecke/TwistingData.lean\#L114-L114}{{\leanfont TwistingData.\allowbreak{}w\_ne\_zero}}}{TwistingData.lean}{114--114}
\crow{finres}{\href{https://github.com/lsernau/grh-lean/blob/72fead102f663d5c7a407964a13dc0cf9d977207/grh-compact/src/New/TwistedHecke/TwistingData.lean\#L117-L117}{{\leanfont TwistingData.\allowbreak{}finite\_residueField}}}{TwistingData.lean}{117--117}
\crow{pdvdpar}{\href{https://github.com/lsernau/grh-lean/blob/72fead102f663d5c7a407964a13dc0cf9d977207/grh-compact/src/New/TwistedHecke/TwistingData.lean\#L119-L124}{{\leanfont TwistingData.\allowbreak{}p\_dvd\_letterMatrix\_of\_mem\_parahoric}}}{TwistingData.lean}{119--124}
\crow{polyActionX}{\href{https://github.com/lsernau/grh-lean/blob/72fead102f663d5c7a407964a13dc0cf9d977207/grh-compact/src/New/TwistedHecke/TwistingData.lean\#L128-L129}{{\leanfont TwistingData.\allowbreak{}polyAction\_X}}}{TwistingData.lean}{128--129}
\crow{diagdepth}{\href{https://github.com/lsernau/grh-lean/blob/72fead102f663d5c7a407964a13dc0cf9d977207/grh-compact/src/New/TwistedHecke/TwistingData.lean\#L131-L133}{{\leanfont TwistingData.\allowbreak{}diagonal\_pow\_depth}}}{TwistingData.lean}{131--133}
\crow{polyOfCoeffs}{\href{https://github.com/lsernau/grh-lean/blob/72fead102f663d5c7a407964a13dc0cf9d977207/grh-compact/src/New/TwistedHecke/TwistingData.lean\#L137-L140}{{\leanfont TwistingData.\allowbreak{}polyAction\_polyOfCoeffs}}}{TwistingData.lean}{137--140}
\crow{letterRepmap}{\href{https://github.com/lsernau/grh-lean/blob/72fead102f663d5c7a407964a13dc0cf9d977207/grh-compact/src/New/TwistedHecke/TwistingData.lean\#L142-L142}{{\leanfont TwistingData.\allowbreak{}letterRep\_map}}}{TwistingData.lean}{142--142}
\crow{oneaddval}{\href{https://github.com/lsernau/grh-lean/blob/72fead102f663d5c7a407964a13dc0cf9d977207/grh-compact/src/New/TwistedHecke/TwistingData.lean\#L144-L145}{{\leanfont TwistingData.\allowbreak{}one\_add\_val\_sub\_one\_sub}}}{TwistingData.lean}{144--145}
\crow{fivemul}{\href{https://github.com/lsernau/grh-lean/blob/72fead102f663d5c7a407964a13dc0cf9d977207/grh-compact/src/New/TwistedHecke/TwistingData.lean\#L147-L148}{{\leanfont TwistingData.\allowbreak{}five\_mul\_P\_pow\_le\_q}}}{TwistingData.lean}{147--148}
\crow{exwIcc}{\href{https://github.com/lsernau/grh-lean/blob/72fead102f663d5c7a407964a13dc0cf9d977207/grh-compact/src/New/TwistedHecke/TwistingData.lean\#L150-L152}{{\leanfont TwistingData.\allowbreak{}excessWeight\_mem\_Icc}}}{TwistingData.lean}{150--152}
\crow{IndexDvd}{\href{https://github.com/lsernau/grh-lean/blob/72fead102f663d5c7a407964a13dc0cf9d977207/grh-compact/src/New/Core/ModuleIndex.lean\#L20-L20}{{\leanfont IndexDvd}}}{ModuleIndex.lean}{20--20}
\crow{HasType}{\href{https://github.com/lsernau/grh-lean/blob/72fead102f663d5c7a407964a13dc0cf9d977207/grh-compact/src/New/Core/ModuleIndex.lean\#L55-L56}{{\leanfont HasType}}}{ModuleIndex.lean}{55--56}
\crow{ofdetdvd}{\href{https://github.com/lsernau/grh-lean/blob/72fead102f663d5c7a407964a13dc0cf9d977207/grh-compact/src/New/Core/ModuleIndex.lean\#L68-L73}{{\leanfont IndexDvd.\allowbreak{}of\_det\_dvd}}}{ModuleIndex.lean}{68--73}
\crow{idtrans}{\href{https://github.com/lsernau/grh-lean/blob/72fead102f663d5c7a407964a13dc0cf9d977207/grh-compact/src/New/Core/ModuleIndex.lean\#L31-L33}{{\leanfont IndexDvd.\allowbreak{}trans}}}{ModuleIndex.lean}{31--33}
\crow{idmap}{\href{https://github.com/lsernau/grh-lean/blob/72fead102f663d5c7a407964a13dc0cf9d977207/grh-compact/src/New/Core/ModuleIndex.lean\#L37-L41}{{\leanfont IndexDvd.\allowbreak{}map}}}{ModuleIndex.lean}{37--41}
\crow{idcomp}{\href{https://github.com/lsernau/grh-lean/blob/72fead102f663d5c7a407964a13dc0cf9d977207/grh-compact/src/New/Core/ModuleIndex.lean\#L50-L51}{{\leanfont IndexDvd.\allowbreak{}comp}}}{ModuleIndex.lean}{50--51}
\crow{hermite}{\href{https://github.com/lsernau/grh-lean/blob/72fead102f663d5c7a407964a13dc0cf9d977207/grh-compact/src/Lib/Adelic/Witt/FrameAction.lean\#L58-L69}{{\leanfont IsValuationPair.\allowbreak{}hermite}}}{FrameAction.lean}{58--69}
\crow{smith}{\href{https://github.com/lsernau/grh-lean/blob/72fead102f663d5c7a407964a13dc0cf9d977207/grh-compact/src/Lib/General/LinearAlgebra/ColumnLattices.lean\#L23-L44}{{\leanfont IsValuationPair.\allowbreak{}smith\_normal\_form}}}{ColumnLattices.lean}{23--44}
\crow{valtop}{\href{https://github.com/lsernau/grh-lean/blob/72fead102f663d5c7a407964a13dc0cf9d977207/grh-compact/src/Lib/Adelic/Witt/FrameAction.lean\#L50-L51}{{\leanfont isValuationPair\_top}}}{FrameAction.lean}{50--51}
\crow{nullst}{\href{https://github.com/lsernau/grh-lean/blob/72fead102f663d5c7a407964a13dc0cf9d977207/grh-compact/src/Lib/General/Algebra/ZerosCostDegree.lean\#L271-L313}{{\leanfont exists\_le\_cost\_of\_vanishing}}}{ZerosCostDegree.lean}{271--313}
\crow{cardGL}{\href{https://github.com/lsernau/grh-lean/blob/72fead102f663d5c7a407964a13dc0cf9d977207/grh-compact/src/Lib/LocalGL/Green/CuspidalAtScalars.lean\#L17-L21}{{\leanfont card\_GL\_eq\_prod}}}{CuspidalAtScalars.lean}{17--21}}
% the notations of the defined functions: one slot for each explicit argument of the Lean declaration
\newcommand{\conjSubst}[2]{\mathrm{cs}_{#1}(#2)}
\newcommand{\exCone}[3]{\mathrm{Ec}_{#1}(#2;#3)}
\newcommand{\frobX}[3]{X^{\mathrm{Fr}}_{#1}(#2,#3)}
\newcommand{\letterConj}[3]{\mathrm{cj}_{#1,#2,#3}}
\newcommand{\idxw}[5]{\mathrm{iw}_{#1,#2}^{#3,#4}(#5)}
\newcommand{\permStep}[3]{\mathrm{St}_{#1,#2,#3}}
\newcommand{\slab}[2]{\mathrm{Sl}_{#1}(#2)}
\newcommand{\pieri}[2]{\mathrm{Pc}_{#1,#2}}
\newcommand{\IsRF}[1]{$#1$ is a real filtration}
\newcommand{\adic}[2]{\mathrm{ad}_{#1}(#2)}
\newcommand{\gauss}[3]{\mathrm{Gs}_{#1,#2}(#3)}
\newcommand{\twoBlock}[3]{\tau_{#1,#2}(#3)}
\newcommand{\TwD}[3]{\mathrm{TD}_{#1}(#2,#3)}
\newcommand{\PT}[1]{P_{#1}}
\newcommand{\qT}[1]{q_{#1}}
\newcommand{\Tmap}[2]{#1_{*}#2}
\newcommand{\letterMatrix}[3]{#2_{#1,#3}}
\newcommand{\polyAction}[2]{[#2]_{#1}}
\newcommand{\letterRep}[3]{\rho_{#1,#2}(#3)}
\newcommand{\dpi}[3]{\mathrm{dp}_{#1,#2}(#3)}
\newcommand{\rootVar}[4]{\mathrm{rt}_{#1,#4}(#2,#3)}
\newcommand{\exw}[5]{\mathrm{ex}_{#1,#2}^{#3,#4}(#5)}
\newcommand{\deepForms}[5]{\mathrm{Dp}_{#1,#4}^{#2,#3}(#5)}
\newcommand{\IndexDvd}[3]{[#1\,{:}\,#2]\mathrel{\big|}#3}
\newcommand{\HasType}[5]{$#4$ has type at most~$#5$ at~$#1^{#3}$ on~$#2$}
\newcommand{\UnitMembers}{78}
\newcommand{\UnitInNote}{75}
\newcommand{\UnitProofs}{32}
\newcommand{\UnitSteps}{320}
\newcommand{\UnitAllSteps}{357}


\linespread{1.08}                       % 11 pt on 14.7 pt
\setlength{\parindent}{0pt}
\setlength{\parskip}{5.5pt plus 1.5pt minus 0.5pt}
\setlength{\abovedisplayskip}{7.5pt plus 2pt minus 3pt}
\setlength{\belowdisplayskip}{7.5pt plus 2pt minus 3pt}
\setlength{\abovedisplayshortskip}{3pt plus 2pt}
\setlength{\belowdisplayshortskip}{6pt plus 2pt minus 2pt}
\emergencystretch=1.5em
\clubpenalty=10000 \widowpenalty=10000 \displaywidowpenalty=10000
\relpenalty=7000 \binoppenalty=9000
\raggedbottom
\pagestyle{fancy}\fancyhf{}\renewcommand{\headrulewidth}{0pt}
\fancyfoot[C]{\sffamily\fontsize{7.6}{9}\selectfont\color{inkfaint}\thepage\,/\,\pageref*{LastPage}}
\setlist{topsep=2pt,partopsep=0pt,parsep=0pt,itemsep=4pt,leftmargin=2.1em,labelsep=0.55em}

\newcommand{\Q}{\mathbb{Q}}
\newcommand{\Z}{\mathbb{Z}}
\newcommand{\F}{\mathbb{F}}
\newcommand{\len}{\operatorname{len}}
\newcommand{\cok}{\operatorname{cok}}
\newcommand{\End}{\operatorname{End}}
\newcommand{\Hom}{\operatorname{Hom}}
\newcommand{\Gal}{\operatorname{Gal}}
\newcommand{\diag}{\operatorname{diag}}
\newcommand{\kw}[1]{\emph{#1}}
\newcommand{\m}[1]{\mbox{$#1$}}        % a short formula that must not be broken across lines
\newcommand{\head}[1]{\par\addvspace{17pt plus 4pt minus 2pt}\Needspace{6\baselineskip}%
  \noindent{\fontsize{13.2}{16}\selectfont\bfseries\color{ink}#1}\par\nopagebreak\vspace{-1pt}}
\newcommand{\sect}[2]{\refstepcounter{section}\label{#2}\par\addvspace{19pt plus 4pt minus 2pt}\Needspace{7\baselineskip}%
  \noindent\llap{{\fontsize{13.2}{16}\selectfont\bfseries\color{inkfaint}\thesection}\hspace{9pt}}%
  {\fontsize{13.2}{16}\selectfont\bfseries\color{ink}#1}\par\nopagebreak\vspace{-1pt}}
\newcommand{\subhead}[1]{\par\addvspace{9pt plus 2pt}\Needspace{4\baselineskip}\noindent\textbf{#1}\enspace\ignorespaces}
% where a statement lives in the Lean archive
\newcommand{\leantag}{{\sffamily\bfseries\fontsize{6.3}{8}\selectfont\color{leancol}\addfontfeatures{LetterSpace=7}LEAN}}
\newcommand{\sep}{\unskip\nobreak\hspace{5pt}{\color{inkfaint}\textperiodcentered}\hspace{5pt plus 2pt}\ignorespaces}
\newcommand{\leanline}[1]{\par\nopagebreak\vspace{-2.5pt}{\fontsize{8.9}{12.4}\selectfont\color{inksoft}\raggedright
  \noindent\llap{\leantag\hspace{9pt}}#1\par}\vspace{1pt}}
% a displayed statement in a tinted box: #1 colour, #2 text
\newcommand{\statement}[2]{\par\addvspace{7pt}\begingroup\setlength{\fboxsep}{8.5pt}\setlength{\fboxrule}{0.6pt}%
  \noindent\fcolorbox{#1!30!white}{#1!3!white}{\parbox{\dimexpr\linewidth-2\fboxsep-2\fboxrule\relax}{#2}}%
  \endgroup\par\addvspace{7pt}}
\newcommand{\nt}[2]{\par\noindent\hangindent=1.1em\hangafter=1 #1\enspace{\color{inksoft}#2}\par\vspace{1pt}}
\newcommand{\legdot}[1]{\tikz[baseline=-0.62ex]\fill[#1] (0,0) circle (2.5pt);\hspace{3pt}}
\newcommand{\crow}[4]{\csname CS@#1\endcsname & #2 & \textcolor{inksoft}{#3} & \textcolor{inksoft}{#4}\\}
\newcommand{\leanquoteline}[2]{\par\noindent\hangindent=\dimexpr#1\fontdimen2\font+2.2em\relax\hangafter=1
  \hspace*{#1\fontdimen2\font}#2}
\newcommand{\quoted}[1]{\par\vspace{3pt}{\ttfamily\fontsize{7.5}{10.4}\selectfont\color{leancol}\raggedright
  \setlength{\parskip}{0pt}\leanquote{#1}\par}\vspace{4pt}}
\newcommand{\bx}[3]{{\sffamily\bfseries\fontsize{7.3}{9.2}\selectfont\color{#1!82!black}#2}\\[1.5pt]#3}
\newcommand{\cs}[2]{\expandafter\newcommand\csname CS@#1\endcsname{#2}}

% theorems
\newtheoremstyle{plainx}{7pt plus 2pt}{7pt plus 2pt}{\itshape}{}{\bfseries}{.}{0.5em}{}
\newtheoremstyle{uprightx}{7pt plus 2pt}{7pt plus 2pt}{}{}{\bfseries}{.}{0.5em}{}
\theoremstyle{plainx}
\newtheorem{theorem}{Theorem}[section]
\newtheorem{lemma}[theorem]{Lemma}
\newtheorem{proposition}[theorem]{Proposition}
\newtheorem{corollary}[theorem]{Corollary}
\theoremstyle{uprightx}
\newtheorem{definition}[theorem]{Definition}
\newtheorem{setting}[theorem]{Setting}          % hypotheses shared by several statements; each of them cites it by number
\newtheorem{remark}[theorem]{Remark}
\newtheorem{example}[theorem]{Example}
\renewcommand{\qedsymbol}{{\setlength{\fboxsep}{0pt}\setlength{\fboxrule}{0.5pt}\fbox{\rule{0pt}{0.46em}\hspace{0.46em}}}}
% what holds throughout the note: the ambient objects only (implicit and instance arguments), never a hypothesis
\newenvironment{conventions}{\par\addvspace{5pt}\noindent\textbf{Conventions.}\enspace\ignorespaces}{\par\addvspace{3pt}}
\makeatletter
\renewenvironment{proof}[1][Proof]{\par\pushQED{\qed}\normalfont\topsep0pt\partopsep0pt
  \trivlist\item[\hskip\labelsep\itshape #1.]\ignorespaces}{\popQED\endtrivlist\@endpefalse}
\makeatother

% where each declaration of the concordance is treated: \cs{key}{place}, the place naming a \label or a lettered item
\cs{rootlemma}{Theorem 1}
\cs{prodvol}{Theorem 2}
\cs{dimex}{Theorem 3}
\cs{dimdeep}{Theorem 3}
\cs{projector}{Theorem 4}
\cs{hecke}{Theorem 5}
\cs{TwD}{Definition~\ref{d:T}}
\cs{PT}{Definition~\ref{d:T}}
\cs{qT}{Definition~\ref{d:T}}
\cs{Tmap}{Definition~\ref{d:T}}
\cs{twoBlock}{Definition~\ref{d:action}}
\cs{letterMatrix}{Definition~\ref{d:action}}
\cs{polyAction}{Definition~\ref{d:action}}
\cs{letterRep}{Definition~\ref{d:action}}
\cs{dpi}{Definition~\ref{d:depth}}
\cs{depth}{Definition~\ref{d:depth}}
\cs{rootVar}{Definition~\ref{d:excess}}
\cs{exw}{Definition~\ref{d:excess}}
\cs{translatesIn}{Definition~\ref{d:spaces}}
\cs{excessCond}{Definition~\ref{d:spaces}}
\cs{deepForms}{Definition~\ref{d:spaces}}
\cs{IsRF}{Definition~\ref{d:filt}}
\cs{adic}{Definition~\ref{d:filt}}
\cs{gauss}{Definition~\ref{d:filt}}
\cs{permStep}{Definition~\ref{d:cones}}
\cs{slab}{Definition~\ref{d:cones}}
\cs{pieri}{Definition~\ref{d:cones}}
\cs{exCone}{Definition~\ref{d:cones}}
\cs{HeckeSizes}{Definition~\ref{d:sizes}}
\cs{frobX}{Definition~\ref{d:frob}}
\cs{frobOperator}{Definition~\ref{d:frob}}
\cs{letterConj}{Definition~\ref{d:frob}}
\cs{idxw}{Definition~\ref{d:frob}}
\cs{conjSubst}{Definition~\ref{d:frob}}
\cs{IndexDvd}{Input 1}
\cs{HasType}{Input 1}
\cs{ofdetdvd}{Input 1}
\cs{idtrans}{Input 1}
\cs{idmap}{Input 1}
\cs{idcomp}{Input 1}
\cs{hermite}{Input 2}
\cs{smith}{Input 2}
\cs{valtop}{Input 2}
\cs{nullst}{Input 3}
\cs{cardGL}{Input 4}
\cs{frobpow}{Lemma~\ref{l:basic}}
\cs{polyActionX}{Lemma~\ref{l:basic}}
\cs{polyOfCoeffs}{Lemma~\ref{l:basic}}
\cs{letterRepmap}{Lemma~\ref{l:basic}}
\cs{diagdepth}{Lemma~\ref{l:basic}}
\cs{oneltP}{Lemma~\ref{l:basic}}
\cs{oneltq}{Lemma~\ref{l:basic}}
\cs{wnezero}{Lemma~\ref{l:basic}}
\cs{finres}{Lemma~\ref{l:basic}}
\cs{pdvdpar}{Lemma~\ref{l:par}}
\cs{exwIcc}{Lemma~\ref{l:exIcc}}
\cs{oneaddval}{Lemma~\ref{l:exIcc}}
\cs{fivemul}{Lemma~\ref{l:exIcc}}
\cs{sumdeep}{Lemma~\ref{l:sumdeep}}
\cs{stepanov}{Lemma~\ref{l:stepanov}}
\cs{aevalsub}{Lemma~\ref{l:subst}}
\cs{memadic}{Lemma~\ref{l:subst}}
\cs{monomialmem}{Lemma~\ref{l:subst}}
\cs{Xmem}{Lemma~\ref{l:subst}}
\cs{detinter}{Lemma~\ref{l:interlace}}
\cs{signedsum}{Proposition~\ref{p:signed}}
\cs{isBounded}{Proposition~\ref{p:signed}}
\cs{borel}{Lemma~\ref{l:borel}}
\cs{exdot}{Lemma~\ref{l:exdot}}
\cs{conepoint}{Lemma~\ref{l:conepoint}}
\cs{simplex}{Lemma~\ref{l:simplex}}
\cs{dimvol}{Proposition~\ref{p:dimvol}}
\cs{deepweight}{Proposition~\ref{p:deepweight}}
\cs{zariski}{Lemma~\ref{l:zariski}}
\cs{diagform}{Lemma~\ref{l:diagform}}
\cs{mean}{Proposition~\ref{p:mean}}
\cs{mapC}{Theorem 4}
\cs{frobmap}{Lemma~\ref{l:frobmap}}
\cs{frobmem}{Lemma~\ref{l:frobmem}}
\cs{cjfrob}{Lemma~\ref{l:cjfrob}}
\cs{powdvd}{Proposition~\ref{p:powdvd}}
\cs{degle}{Lemma~\ref{l:degle}}
\cs{tensordual}{Lemma~\ref{l:tensor}}
\cs{cardle}{Lemma~\ref{l:cardle}}
\cs{pairing}{Proposition~\ref{p:pairing}}
\cs{depthfact}{Proposition~\ref{p:depthfact}}
\newcommand{\cO}{\mathcal{O}}
\newcommand{\FF}{\mathbb{F}}
\newcommand{\RR}{\mathbb{R}}
\newcommand{\cF}{\mathcal{F}}
\newcommand{\cG}{\mathcal{G}}
\newcommand{\cA}{\mathcal{A}}
\newcommand{\cR}{\mathcal{R}}
\newcommand{\cX}{\mathcal{X}}
\newcommand{\cE}{\mathcal{E}}
\newcommand{\cH}{\mathcal{H}}
\newcommand{\cI}{\mathcal{I}}
\newcommand{\cV}{\mathcal{V}}
\newcommand{\cC}{\mathcal{C}}
\newcommand{\cL}{\mathcal{L}}
\newcommand{\cM}{\mathcal{M}}
\newcommand{\cQ}{\mathcal{Q}}
\newcommand{\cS}{\mathcal{S}}
\newcommand{\cK}{\mathcal{K}}
\newcommand{\cN}{\mathcal{N}}
\newcommand{\fm}{\mathfrak{m}}
\newcommand{\fx}{\mathfrak{x}}
\newcommand{\fq}{\mathfrak{q}}
\newcommand{\fh}{\mathfrak{h}}
\newcommand{\fii}{\mathfrak{i}}
\newcommand{\fe}{\mathfrak{e}}
\newcommand{\fs}{\mathfrak{s}}
\newcommand{\GLk}{\mathrm{GL}_{k}}
\newcommand{\Mk}{\mathrm{M}_{k}}
\newcommand{\Lev}{\mathrm{Lev}}
\newcommand{\Par}{\mathrm{Par}}
\newcommand{\Mon}{\mathrm{Mon}}
\newcommand{\Fam}{\mathrm{Mon}_{k,n}^{L}}
\newcommand{\ZL}{\mathbb{Z}/L}
\newcommand{\vol}{\operatorname{vol}}
\newcommand{\sgn}{\operatorname{sgn}}
\newcommand{\pr}[2]{\langle#1,#2\rangle}
\newcommand{\tp}{\mathsf{T}}
\newcommand{\thr}{\mathfrak{t}}        % the threshold of the depth (Lean's r); the write-up keeps \mu for the exponent of the data
\newcommand{\pa}{\mathrm{par}}          % the parity map of the letters (Lean's isEven); the write-up keeps \varsigma for a root of unity
% the four definitions whose reading by program lists a bound variable as an argument, and the predicate of ten
% arguments: written with the arguments that Lean takes (see the spec, `notation: none`)
\newcommand{\dep}[5]{\mathrm{depth}_{#1,#4}^{#2,#3}(#5)}
\newcommand{\trl}[5]{\mathrm{Tr}_{#1,#2}(#3;#4,#5)}
\newcommand{\exc}[6]{\mathrm{Ex}_{#1,#2}^{#3,#4}(#5;#6)}
\newcommand{\frob}[5]{\mathrm{Fr}_{#1,#2,#3,#4}(#5)}
\newcommand{\Sz}[2]{\mathrm{Sz}(#1;#2)}

\begin{document}
\color{ink}

{\raggedright
{\sffamily\bfseries\fontsize{7.2}{9}\selectfont\color{leancol}\addfontfeatures{LetterSpace=8}READBACK NOTE 40\par}
\vspace{5pt}
{\fontsize{21.5}{25.5}\selectfont\bfseries Elementary divisors of a Frobenius-twisted Hecke operator\par}
\vspace{5pt}
{\fontsize{12.4}{16}\selectfont\itshape\color{inksoft} One unit of the Lean artifact \ArtifactName, read back into
mathematics\par}}
\vspace{8pt}
{\color{rulegrey}\hrule height 0.6pt}
\vspace{7pt}
\thispagestyle{empty}

\textbf{Summary.} A group \m{\GLk(W)} over a local ring with finite residue field acts on forms in $L$ sets of
$k$ variables, through $L$ ring homomorphisms that are linked in a circle by a power of Frobenius modulo a
prime~$p$, alternately directly and by the inverse transpose. A Hecke operator of this action, scaled by powers
of~$p$ so as to be integral, is shown to have small elementary divisors: modulo~\m{p^{e}} its cokernel has length
at most a small share of $e$ times the rank, plus a constant times the rank. Four ideas carry the proof. A \emph{root lemma}: a
term of a form that is least, for valuation plus a weight, among all translates under a Levi subgroup has
exponents that obey linear inequalities, by a count of zeros against degree. Two \emph{dimension bounds}
modulo~$p$: forms whose exponents nearly obey these inequalities in all frames fill about one part in the degree
of the operator, by a count of flags against the volume of a cone of cyclically interlacing vectors; and forms
that the scaling makes divisible in all frames are rare. An \emph{approximate projector} onto the first kind of
forms, made from conjugation-invariant means on polynomials in generic matrices. And normal forms of matrices
over a valuation ring, which turn the three into a statement about indices.

\statement{leancol}{\fontsize{9.6}{13.2}\selectfont\textbf{What this note is.} The artifact \ArtifactName{} is a
body of Lean code: the compact edition of a formal proof, due to L.~Alpöge with the AI system Claude, of the
Riemann hypothesis and of its generalisation to Dirichlet $L$-functions. That headline is a claim. The authors
report that every file compiles with Lean v4.35.0 against a pinned version of the library Mathlib and that Lean's
kernel accepted the proofs; an independent build is recorded in the handout's record of checks when it reports.
No mathematician has read the argument through, and nothing here says that the Riemann hypothesis is proved. This
note treats one unit of the artifact, the folder on the twisted Hecke operator: \UnitInNote{} of its \UnitMembers{}
declarations, which are all that carry a name (the others are three instances, described in §\ref{s:notes}).
Proofs are given in full, \emph{given} four inputs from other units: the calculus of index divisibility
(unit~149), the Hermite and Smith normal forms over a valuation ring, a combinatorial Nullstellensatz with
multiplicity, and the order of the general linear group of a finite field. Standard facts of algebra and of
measure theory are used without proof: Nakayama's lemma, the freeness of a finitely
generated torsion-free module over a valuation domain, the vanishing of the trace of a nilpotent endomorphism of
a free module of finite rank over a domain, Legendre's formula for the power of a prime in a factorial, the change of
variables for Lebesgue measure under a linear map, and the convergence of normalised counts of lattice points to
the volume of a bounded convex set. Every statement of the unit and of
the Inputs was compared with its Lean text, argument by argument. The Lean text read is \PinText; Appendix~A is
the concordance of \ConcDecls{} declarations.}

\sect{The statements}{s:intro}
The notions are defined in §\ref{s:setting}; here is what is needed to read the five theorems. \emph{Twisting
data}~$T$ of length~$L$ over a local ring~$W$ and a ring~$\cO$ are a prime~$p$, a power \m{\PT{T}=p^{o}\ge5}, ring
homomorphisms \m{\iota_s\colon W\to\cO} for \m{s\in\ZL} (the \emph{letters}) with
\m{\iota_{s+1}(\upsilon)\equiv\iota_s(\upsilon)^{\PT{T}}} modulo~$p$ for $\upsilon$ in~$W$, a parity of the letters that alternates, and an element
$\varpi$ of~$\cO$ with \m{\varpi^{2}=p}; the residue field of~$W$ has \m{\qT{T}=\PT{T}^{L}} elements. An element
$\gamma$ of \m{\GLk(W)} acts on the variables \m{x_{i,s}} (\m{0\le i<k}, \m{s\in\ZL}) of letter~$s$ through the
matrix \m{\letterMatrix{T}{\gamma}{s}}, which is \m{\iota_s(\gamma)} for even~$s$ and
\m{\iota_s((\gamma^{-1})^{\tp})} for odd~$s$; \m{\polyAction{T}{\gamma}} is this action on polynomials, and
\m{\letterRep{T}{n}{\gamma}} its matrix on the forms of degree~$n$ in each letter, in the basis of the monomials
\m{x^{\beta}}, \m{\beta\in\Fam}. For a pair of indices \m{(i_0,i_1)} and a letter~$s_0$,
\m{\exw{T}{s_0}{i_0}{i_1}{\cdot}} is a real weight on the variables, the \emph{excess weight}; for a
threshold~$\thr$, \m{\dep{T}{k}{n}{\thr}{\beta}} is the degree of \m{x^{\beta}} in the variables that are
\emph{deep} for~$\thr$. \m{\exc{T}{n}{i_0}{i_1}{I}{h}} is the module of the forms~$F$ such that every translate
of~$F$ by a block diagonal element has, modulo the ideal~$I$, only monomials of excess at most~$h$ at every
letter; \m{\deepForms{T}{k}{n}{\thr}{j}} is the module of the forms all of whose translates have only monomials
of depth at least~$j$. \m{\Sz{\PT{T},o,L,k,n,\thr}{\lambda,\kappa,y,z}} is a list of ten conditions on
sizes (Definition~\ref{d:sizes}).

\statement{arith}{\phantomsection\label{t:1}\textbf{Theorem 1} (the root lemma).\enspace
\emph{In Setting~\ref{set:T}, let $s_0$ be a letter, let \m{i_0\ne i_1} be indices, let $\tau$ be a labelling of
the indices with \m{\tau(i_0)=\tau(i_1)}, and let $\omega$ be a real weight on the variables with}
\[0.4\ \le\ \omega\bigl(\rootVar{T}{i_0}{i_1}{s}\bigr)-\omega\bigl(\rootVar{T}{i_1}{i_0}{s}\bigr)\ \le\ 1.6
\qquad\text{\emph{for every letter} } s .\]
\emph{Let $\Phi$ be a polynomial in the variables \m{x_v} over~$\cO$, let $\alpha$ be an exponent and $e$ a
natural number. Suppose that \m{\polyAction{T}{\gamma^{-1}}\Phi} lies in
\m{\gauss{\adic{\varpi}{\cdot}}{\omega}{e+\pr{\omega}{\alpha}}} for every $\gamma$ in \m{\Lev_{\tau}(W)}: the
coefficient of \m{x^{\alpha'}} in every such translate is divisible by~$\varpi$ to the power
\m{\max(\lceil e+\pr{\omega}{\alpha}-\pr{\omega}{\alpha'}\rceil,0)}. Suppose that \m{\varpi^{e+1}} does not
divide the coefficient of \m{x^{\alpha}} in~$\Phi$. Then}
\[\pr{\exw{T}{s_0}{i_0}{i_1}{\cdot}}{\alpha}\ \le\ 0 .\]}
\leanline{\lean{rootlemma}}

\statement{arith}{\phantomsection\label{t:2}\textbf{Theorem 2} (the cyclic interlacing cone fills the share of
one flag).\enspace
\emph{Let \m{k\ge0} and \m{L\ge1} be integers, $\pa$ a map from \m{\ZL} to the signs \m{\{+,-\}}, and
$\varrho$ a real number with \m{1<\varrho}. Let \m{\delta=(\delta_s)_{s\in\ZL}} be real numbers and $\cC$ a subset
of \m{\RR^{k\times L}} such that, for every $\xi$ in \m{\RR^{k\times L}}: \m{\permStep{\pa}{\varrho}{1}(\xi)}
lies in~$\cC$ if and only if $\xi$ lies in \m{\pieri{\pa}{\varrho}}. Then}
\[\prod_{l=1}^{k}\frac{\varrho^{Ll}-1}{\varrho^{L}-1}\ \cdot\ \vol\bigl(\cC\cap\slab{k}{\delta}\bigr)\ =\
\vol\bigl(\slab{k}{\delta}\bigr).\]}
\leanline{\lean{prodvol}}

The identity holds for every real \m{\varrho>1}. When \m{\varrho^{L}} is the number~$q$ of elements of a finite
field, the product is the number of full flags in a vector space of dimension~$k$ over that field.

\statement{arith}{\phantomsection\label{t:3}\textbf{Theorem 3} (two dimension bounds in
characteristic~$p$).\enspace
\emph{In Setting~\ref{set:T}, let the ring of~$T$ be a field~$\FF$ of characteristic~$p$, the prime of~$T$. Let
$n$, $\lambda$ be natural numbers and $\kappa$, $z$ real numbers.}

\emph{\textup{(a)} Let \m{i_0}, \m{i_1} be indices with \m{i_0+1=i_1}, and $h$ a real number such that
\m{\Sz{\PT{T},o,L,k,n,i_1}{\lambda,\kappa,h,z}}. Then}
\[\bigl[\GLk(W):\Par_{\twoBlock{k}{i_1}{\cdot}}(W)\bigr]\cdot\dim_{\FF}\exc{T}{n}{i_0}{i_1}{0}{h}\ \le\
(1+\kappa)\cdot\bigl(\#\Mon_{k,n}\bigr)^{L}.\]

\emph{\textup{(b)} Let $\thr$ be a natural number and $y$ a real number such that
\m{\Sz{\PT{T},o,L,k,n,\thr}{\lambda,\kappa,y,z}}, and let $j_0$ be a natural number. Then}
\[\sum_{j<j_0}\dim_{\FF}\deepForms{T}{k}{n}{\thr}{j}\ \le\ (1+z)\cdot\bigl(\#\Mon_{k,n}\bigr)^{L}.\]}
\leanline{\lean{dimex}\sep \lean{dimdeep}}

Here \m{(\#\Mon_{k,n})^{L}} is the number of the monomials \m{x^{\beta}}, the rank of the module of forms. In~(a)
the index of the parahoric subgroup is the number of cosets over which a Hecke operator sums, its degree; so the
forms obeying the excess condition fill about one part in the degree.

\statement{arith}{\phantomsection\label{t:4}\textbf{Theorem 4} (an approximate projector).\enspace
\emph{In Setting~\ref{set:T}, let $\cO$ be a valuation domain of characteristic zero in which the prime~$p$
of~$T$ is not a unit. Let $c$, $n$, $\lambda$ be natural numbers, \m{i_0}, \m{i_1} indices with \m{i_0+1=i_1},
and $h$ a real number. Suppose that \m{\lambda+\lambda=L} and
\m{L\,n\,k\,(\lambda+2)+n+k<\PT{T}^{\lambda}}. Then there are a function $w$ from \m{\Fam} to~$\cO$ and a square
matrix~$Z$ over~$\cO$ indexed by \m{\Fam} such that:}
\begin{enumerate}[label=\textup{(\roman*)},leftmargin=2.3em,itemsep=1pt,topsep=2pt]
\item \emph{\m{ZF} lies in \m{\exc{T}{n}{i_0}{i_1}{\varpi^{c}\cO}{h}} for every $F$ in \m{\cO^{\Fam}};}
\item \emph{for every $F'$ in \m{\exc{T}{n}{i_0}{i_1}{\varpi^{c}\cO}{h-(4c+2)}} and every $F$ in
\m{\cO^{\Fam}}, \m{\varpi^{c}} divides}
\[\sum_{\beta}F'(\beta)\,w(\beta)\,(ZF)(\beta)\ -\ p^{(k+k^{2})\,oL\cdot L}\sum_{\beta}F'(\beta)\,F(\beta);\]
\item \emph{\m{w(\beta)\,\letterRep{T}{n}{\gamma^{\tp}}_{\beta,\beta'}=w(\beta')\,\letterRep{T}{n}{\gamma}_{\beta',\beta}}
for every $\gamma$ in \m{\GLk(W)} and all \m{\beta,\beta'} in \m{\Fam}.}
\end{enumerate}}
\leanline{\lean{projector}}

\statement{arith}{\phantomsection\label{t:5}\textbf{Theorem 5} (elementary divisors of the twisted Hecke
operator).\enspace
\emph{In Setting~\ref{set:T}, let $\cO$ be a valuation domain of characteristic zero in which the prime~$p$
of~$T$ is not a unit. Let $n$, $\thr$, $\lambda$ be natural numbers and $\kappa$, $y$, $z$ real numbers with
\m{\Sz{\PT{T},o,L,k,n,\thr}{\lambda,\kappa,y,z}}. Let $\cX$, $\cH$, $\cI$ be finite sets, write
\m{\cQ=\GLk(W)/\Par_{\twoBlock{k}{\thr}{\cdot}}(W)} for the finite set of left cosets, and \m{\cE=\cX\times\cQ}. Let $V$
be an $\cO$-module, \m{\Psi\colon\cO^{\cX\times\cI}\to V} an $\cO$-linear map, and $Y$ an $\cO$-linear endomorphism
of~$V$. Let \m{\mathrm{src}\colon\cE\to\cX} be a map, \m{(\gamma_{\fe})_{\fe\in\cE}} elements of
\m{\GLk(W)}, \m{(u_{\fe})_{\fe\in\cE}} elements of \m{\mathrm{GL}_{\cH}(\cO)} \textup{(}invertible matrices
over~$\cO$ indexed by~$\cH$\textup{)}, and
\m{\jmath\colon\cI\to\Fam\times\cH}, \m{\fii\mapsto(\beta_{\fii},\fh_{\fii})}, a bijection. Suppose that the map
\m{\fe\mapsto\bigl(\mathrm{src}(\fe),\,(\gamma_{\fe}^{\tp})^{-1}\Par_{\twoBlock{k}{\thr}{\cdot}}(W)\bigr)}
from $\cE$ to~$\cE$ is injective, and that for every $F$ in \m{\cO^{\cX\times\cI}}}
\begin{gather*}
Y(\Psi(F))=\Psi(F^{\flat}),\\
F^{\flat}(\fx,\fii)=\sum_{\fq\in\cQ}\,\sum_{\fii'\in\cI}(u_{(\fx,\fq)})_{\fh_{\fii},\fh_{\fii'}}\,
\bigl(H_{(\fx,\fq)}\bigr)_{\beta_{\fii'},\beta_{\fii}}\;F\bigl(\mathrm{src}(\fx,\fq),\fii'\bigr),
\end{gather*}
\emph{where
\m{H_{(\fx,\fq)}=\letterRep{T}{n}{\gamma_{(\fx,\fq)}}\cdot\operatorname{diag}\bigl(p^{\dep{T}{k}{n}{\thr}{\beta}}\bigr)_{\beta}\cdot\letterRep{T}{n}{\tilde{\fq}^{-1}}}
and \m{\tilde{\fq}} is the chosen representative of the coset~$\fq$. Then, for every natural number~$e$:}

\vspace{2pt}\centerline{\HasType{p}{\Psi(\cO^{\cX\times\cI})}{e}{Y}{(z+\kappa e+y)\cdot\#\cX\cdot\#\cI}.}}
\leanline{\lean{hecke}}

The last line says (Input~1): $Y$ maps the image of~$\Psi$ into itself, and the index of
\m{Y(\operatorname{im}\Psi)+p^{e}\operatorname{im}\Psi} in \m{\operatorname{im}\Psi} divides \m{p^{l}} for a natural number
\m{l\le(z+\kappa e+y)\,\#\cX\,\#\cI}. The rank of \m{\cO^{\cX\times\cI}} is \m{\#\cX\cdot\#\cI}; so of the
\m{e\cdot\#\cX\cdot\#\cI} units of length that \m{p^{e}} could cost, the operator loses the share~$\kappa$, and a
constant \m{z+y} for each basis vector. The size conditions make $\kappa$ small when~$y$ is large and $n$ larger
still. The operator \m{F\mapsto F^{\flat}} is a sum over the cosets of the parahoric subgroup, as a Hecke operator
is, of the scaling \m{p^{\mathrm{depth}}} between two translations, twisted by the matrices~$u_{\fe}$ on an
auxiliary space and by the map~\m{\mathrm{src}} on the points.

The note on the unit of the columns (unit~145) quotes this theorem as its Input~I10, with other letters: it writes
$n$, $k'$, $a'$, $q'$, $r'$, $l'$, $c_1$, $c_2$, $c_3$, \m{\phi_s} and~$s$ for what is here $L$, $k$, $n$, \m{\PT{T}},
$\thr$, $\lambda$, $\kappa$, $y$, $z$, \m{\iota_s} and~$e$, and it states the theorem for a valuation ring~$W$.

\begin{center}
\begin{tikzpicture}[x=1mm,y=1mm,>={Stealth[length=1.6mm]},
  bx/.style={draw=rulegrey,line width=0.5pt,rounded corners=1.2pt,fill=arith!3!white,inner xsep=4pt,inner ysep=3.5pt,
    align=center,font=\fontsize{7.9}{9.8}\selectfont},
  gx/.style={bx,fill=inksoft!5!white}]
\node[gx] (st) at (0,0) {count of zeros\\(Lemma~\ref{l:stepanov})};
\node[bx] (rl) at (36,0) {Theorem 1\\root lemma};
\node[gx] (df) at (76,0) {depth factorisation\\(Proposition~\ref{p:depthfact})};
\node[bx] (vol) at (0,-15) {Theorem 2\\cone volume};
\node[bx] (ta) at (36,-15) {Theorem 3(a)\\excess forms};
\node[bx] (t5) at (116,-22) {Theorem 5};
\node[gx] (ns) at (0,-30) {Nullstellensatz\\(Input~3)};
\node[bx] (tb) at (36,-30) {Theorem 3(b)\\deep forms};
\node[gx] (mean) at (0,-45) {invariant means\\(Proposition~\ref{p:mean})};
\node[bx] (pj) at (36,-45) {Theorem 4\\projector};
\node[gx] (pa) at (76,-45) {pairing\\(Proposition~\ref{p:pairing})};
\draw[->,inksoft] (st) -- (rl);
\draw[->,inksoft] (rl) -- (ta);
\draw[->,inksoft] (vol) -- (ta);
\draw[->,inksoft] (ns) -- (tb);
\draw[->,inksoft] (mean) -- (pj);
\draw[->,inksoft] (rl) -- (df);
\draw[->,inksoft] (df.east) -- (t5.north west);
\draw[->,inksoft] (ta.east) -- (t5.west);
\draw[->,inksoft] (tb.east) -- (t5.west);
\draw[->,inksoft] (pj.east) -- (t5.south west);
\draw[->,inksoft] (pa.east) -- (t5.south);
\end{tikzpicture}
\end{center}
\vspace{-4pt}
{\fontsize{8.9}{12}\selectfont\color{inksoft}\noindent How the five theorems depend on each other and on the main
intermediate results.\par}

\sect{The setting}{s:setting}
\begin{conventions}\phantomsection\label{c:conventions}
The natural numbers are \m{0,1,2,\dots}. Every ring is a commutative ring with~$1$. A \kw{local ring}~$W$ has
one maximal ideal~$\fm_W$ and the residue field \m{\FF_W=W/\fm_W}. A \kw{valuation domain} is an integral domain
in which, of any two elements, one divides the other; it is a local ring, and its maximal ideal is written~$\fm$.
The number of elements of a finite set~$\cS$ is \m{\#\cS}; for the residue field of a local ring the formal
text uses a count that is~$0$ for an infinite field (§\ref{s:notes}, note~5).

\emph{Indices, letters, variables.} $k$ and $L$ are natural numbers, \m{L\ge1} unless said otherwise. The
\kw{indices} are \m{0,\dots,k-1}, written $i$, $i'$, $l$. The \kw{letters} are the elements of \m{\ZL}, written
$s$, $s'$, $s_0$; \m{s+1} is the next letter. A \kw{variable} is a pair \m{v=(i,s)} of an index and a letter; the
set of variables is finite. A \kw{labelling} is a map~$\tau$ from the indices to the natural numbers. For a
labelling~$\tau$ and a local ring~$W$, \m{\Lev_{\tau}(W)} is the \kw{Levi subgroup}, of the $\gamma$ in
\m{\GLk(W)} with \m{\gamma_{ii'}=0} whenever \m{\tau(i)\ne\tau(i')}, and \m{\Par_{\tau}(W)} is the \kw{parahoric
subgroup}, of the $\gamma$ with \m{\gamma_{ii'}\in\fm_W} whenever \m{\tau(i)<\tau(i')}. For the constant labelling
\m{\Lev_{\tau}(W)=\GLk(W)}; for the labelling \m{i\mapsto i}, \m{\Par_{\tau}(W)} consists of the matrices that are
lower triangular modulo~$\fm_W$ (the lower Iwahori subgroup).

\emph{Polynomials.} For a set $\cV$ of variables and a ring~$\cO$, \m{\cO[x_{\cV}]} is the polynomial ring in
the variables \m{x_v}, \m{v\in\cV}. An \kw{exponent} is a function~$\alpha$ from $\cV$ to the natural numbers
with finite support; \m{x^{\alpha}=\prod_vx_v^{\alpha(v)}}, \m{|\alpha|=\sum_v\alpha(v)} is its degree, and for a
real function~$\omega$ on~$\cV$, a \kw{weight}, \m{\pr{\omega}{\alpha}=\sum_v\omega(v)\alpha(v)}. The same
notation is used for vectors~$\alpha$ with real coordinates. For a square matrix~$G$ over~$\cO$ indexed by a
finite set~$\cV$, the \kw{linear substitution} of~$G$ is the endomorphism of the $\cO$-algebra \m{\cO[x_{\cV}]}
with \m{x_{v'}\mapsto\sum_vG_{vv'}x_v}; it is multiplicative in~$G$. \m{\Mon_{k,n}} is the set of the exponents of
degree~$n$ on the indices, and \m{\mathrm{Sym}^{n}(G)}, for $G$ in \m{\Mk(\cO)}, is the matrix indexed by
\m{\Mon_{k,n}} whose \m{(\beta,\beta')} entry is the coefficient of \m{x^{\beta}} in the linear substitution of~$G$
applied to \m{x^{\beta'}}. A \kw{family} is an element \m{\beta=(\beta_s)_{s}} of
\m{\Fam=(\Mon_{k,n})^{\ZL}}; it is also read as the exponent \m{(i,s)\mapsto\beta_s(i)} on the variables, of
degree~$n$ in each letter. A \kw{form} is a vector \m{F=(F(\beta))_{\beta}} in \m{\cO^{\Fam}}, also read as the
polynomial \m{\sum_{\beta}F(\beta)x^{\beta}}. For matrices \m{G_s} indexed by \m{\Mon_{k,n}}, one for each letter,
\m{\bigotimes_sG_s} is the matrix indexed by \m{\Fam} with entries \m{\prod_s(G_s)_{\beta_s,\beta'_s}}. Matrices
act on column vectors; \m{G^{\tp}} is the transpose.

\emph{Volumes.} \m{\RR^{k\times L}} is the space of the real functions on the variables, with Lebesgue
measure~\m{\vol}; a volume is a real number, finite in all cases that occur.
\end{conventions}

\begin{definition}[twisting data]\label{d:T}
Let $L$ be a natural number, $W$ a local ring and $\cO$ a ring. \kw{Twisting data of length~$L$ over $W$
and~$\cO$} are a tuple \m{T=(p,o,(\iota_s)_{s\in\ZL},\pa_T,\varpi)} of natural numbers $p$ and~$o$, ring
homomorphisms \m{\iota_s\colon W\to\cO}, a map $\pa_T$ from \m{\ZL} to the signs \m{\{+,-\}}, the \kw{parity}
(the letter~$s$ is \kw{even} if \m{\pa_T(s)=+} and \kw{odd} otherwise), and an element~$\varpi$ of~$\cO$, such
that: $p$ is prime; \m{5\le p^{o}}; \m{\varpi^{2}=p} in~$\cO$; $p$ divides \m{\iota_{s+1}(\upsilon)-\iota_s(\upsilon)^{p^{o}}}
in~$\cO$ for every letter~$s$ and every $\upsilon$ in~$W$; \m{\pa_T(s+1)\ne\pa_T(s)} for every letter~$s$; and
\m{\#\FF_W=(p^{o})^{L}}. The set of all twisting data is \m{\TwD{L}{W}{\cO}}. For $T$ in it we put
\[\PT{T}=p^{o},\qquad \qT{T}=\PT{T}^{L},\]
and for a ring homomorphism \m{\chi\colon\cO\to\cO'} we write
\[\Tmap{\chi}{T}=\bigl(p,o,(\chi\circ\iota_s)_s,\pa_T,\chi(\varpi)\bigr),\]
an element of \m{\TwD{L}{W}{\cO'}}.
\end{definition}
\leanline{\lean{TwD}\sep \lean{PT}\sep \lean{qT}\sep \lean{Tmap}}

For \m{L=0} there are no twisting data, as a field has at least two elements. For \m{L\ge1} the alternation
forces $L$ to be even. The components of~$T$ keep their names: $p$, $o$, $\iota_s$, $\pa_T$ and~$\varpi$
always belong to the twisting data called~$T$ in the statement at hand, and those of \m{\Tmap{\chi}{T}} are as
just said.

\begin{setting}\label{set:T}
\m{L\ge1} and \m{k\ge0} are integers, $W$ is a local ring, $\cO$ is a ring, and $T$ is an element of
\m{\TwD{L}{W}{\cO}}, with components \m{p,o,(\iota_s)_s,\pa_T,\varpi}. Indices, letters and variables are
those of the conventions, for this $k$ and~$L$.
\end{setting}

\emph{Letters.} The letters of this note are local to it. The threshold of the depth is written~$\thr$ and the
parity~$\pa$, because the write-up keeps $\mu$ and $\varsigma$ for the exponent of the data and for a root of
unity. The letters $\lambda$, $\kappa$, $\tau$, $\sigma$, $m$, $q$, $\theta$ and the polynomials~$\Phi$ have no relation
to the objects that the overview of the handout writes with them (an arithmetic function and a constant of
Chapter~I, the infinite places, the automorphisms of the complex numbers, the number of lifts, a rational number
of Theorem~A, the family of twists, a function of Chapter~I), and \m{\PT{T}} none to the parameters. In the
application (§\ref{s:use}) $k$ and~$n$ become the rank of
the columns and the number \m{a(D)} of the data~$D$.

\begin{definition}[the action on forms]\label{d:action}
In Setting~\ref{set:T}, let $\gamma$ be an element of \m{\GLk(W)}, $s$ a letter and $n$ a natural number.
\begin{enumerate}[label=\textup{(\alph*)},leftmargin=2.3em,itemsep=2pt]
\item \m{\letterMatrix{T}{\gamma}{s}} is the matrix in \m{\Mk(\cO)} obtained by applying $\iota_s$ to the entries
of~$\gamma$ if $s$ is even, and of \m{(\gamma^{-1})^{\tp}} if $s$ is odd. The map
\m{\gamma\mapsto\letterMatrix{T}{\gamma}{s}} is multiplicative.
\item \m{\polyAction{T}{\gamma}} is the endomorphism of the $\cO$-algebra of polynomials in the variables
\m{x_{i,s}} that acts on the variables of each letter~$s$ by the linear substitution of
\m{\letterMatrix{T}{\gamma}{s}}. The map \m{\gamma\mapsto\polyAction{T}{\gamma}} is multiplicative: \m{\GLk(W)}
acts on the left.
\item \m{\letterRep{T}{n}{\gamma}=\bigotimes_{s}\mathrm{Sym}^{n}\bigl(\letterMatrix{T}{\gamma}{s}\bigr)}, a square
matrix over~$\cO$ indexed by \m{\Fam}, multiplicative in~$\gamma$.
\item For natural numbers $k$, $\thr$ and an index~$i$, \m{\twoBlock{k}{\thr}{i}} is $0$ if \m{i<\thr} and $1$
otherwise. The labelling \m{\twoBlock{k}{\thr}{\cdot}} has the two blocks \m{\{i<\thr\}} and \m{\{i\ge\thr\}}.
\end{enumerate}
\end{definition}
\leanline{\lean{letterMatrix}\sep \lean{polyAction}\sep \lean{letterRep}\sep \lean{twoBlock}}

\begin{definition}[deep variables and depth]\label{d:depth}
In Setting~\ref{set:T}, let $\thr$ and $n$ be natural numbers. For a variable \m{v=(i,s)} put
\m{\dpi{T}{\thr}{v}=1} if $s$ is even and \m{i\ge\thr}, or $s$ is odd and \m{i<\thr}, and \m{\dpi{T}{\thr}{v}=0}
otherwise; in the first case $v$ is \kw{deep} for~$\thr$. For a family~$\beta$ in \m{\Fam}, its \kw{depth} is
\[\dep{T}{k}{n}{\thr}{\beta}=\sum_{v=(i,s)}\beta_s(i)\,\dpi{T}{\thr}{v}=\pr{\dpi{T}{\thr}{\cdot}}{\beta},\]
the degree of \m{x^{\beta}} in the deep variables. (A reading of the formal text by program gives the depth a
sixth argument: it takes the name under which the degree is passed on for a further one.)
\end{definition}
\leanline{\lean{dpi}\sep \lean{depth}}

The scaling of the Hecke operator is the diagonal matrix \m{\operatorname{diag}(p^{\dep{T}{k}{n}{\thr}{\beta}})_{\beta}}.
It is the action on forms of the diagonal matrix with entries $1$ at the indices \m{i<\thr} and $p$ at the
others, read through the letters: directly at the even ones, and by the inverse transpose, up to the scalar~$p$,
at the odd ones.

\begin{definition}[root variables and excess]\label{d:excess}
In Setting~\ref{set:T}, let \m{i_0}, \m{i_1} be indices and $s$, $s_0$ letters. The \kw{root variable} is
\[\rootVar{T}{i_0}{i_1}{s}=(i_1,s)\ \text{ if $s$ is even},\qquad\rootVar{T}{i_0}{i_1}{s}=(i_0,s)\ \text{ if $s$ is odd}.\]
For a letter~$s$ let \m{[s-s_0]} be the representative of \m{s-s_0} in \m{\{0,\dots,L-1\}}. The \kw{excess
weight} at~$s_0$ is the real function on the variables \m{v=(i,s)} given by
\[\exw{T}{s_0}{i_0}{i_1}{v}=\frac{1}{\qT{T}}\Bigl(\bigl[v=\rootVar{T}{i_0}{i_1}{s}\bigr]\cdot\PT{T}^{\,1+[s-1-s_0]}
-\bigl[v=\rootVar{T}{i_1}{i_0}{s}\bigr]\cdot\PT{T}^{\,[s-s_0]}\Bigr),\]
where a bracket around a statement is $1$ if it holds and $0$ if not. The \kw{excess} of an
exponent~$\alpha$ at~$s_0$ is \m{\pr{\exw{T}{s_0}{i_0}{i_1}{\cdot}}{\alpha}}.
\end{definition}
\leanline{\lean{rootVar}\sep \lean{exw}}

Since \m{1+[s-1-s_0]} is $L$ for \m{s=s_0} and \m{[s-s_0]} otherwise, the excess of~$\alpha$ at~$s_0$ is, for
\m{i_0\ne i_1} and with \m{\alpha^{+}_s=\alpha(\rootVar{T}{i_0}{i_1}{s})} and
\m{\alpha^{-}_s=\alpha(\rootVar{T}{i_1}{i_0}{s})},
\begin{equation}\label{e:excess}
\pr{\exw{T}{s_0}{i_0}{i_1}{\cdot}}{\alpha}=\alpha^{+}_{s_0}-\frac{\alpha^{-}_{s_0}}{\qT{T}}
+\sum_{s\ne s_0}\frac{\PT{T}^{[s-s_0]}}{\qT{T}}\,\bigl(\alpha^{+}_s-\alpha^{-}_s\bigr).
\end{equation}
The two root variables of a letter are the two variables that the elementary matrix
\m{1+\upsilon E_{i_0i_1}}, with $\upsilon$ in~$W$, couples at that letter; here \m{E_{ii'}} is the matrix with the
entry~$1$ at \m{(i,i')} and $0$ elsewhere. In words: at an even letter it adds a multiple of \m{x_{i_0,s}} to \m{x_{i_1,s}}, at an odd letter a multiple
of \m{x_{i_1,s}} to \m{x_{i_0,s}}.

\begin{definition}[conditions on forms]\label{d:spaces}
In Setting~\ref{set:T}, let $n$ be a natural number.
\begin{enumerate}[label=\textup{(\alph*)},leftmargin=2.3em,itemsep=2pt]
\item For a labelling~$\tau$, an ideal~$I$ of~$\cO$ and a subset $\mathcal{B}$ of \m{\Fam},
\m{\trl{T}{n}{\tau}{I}{\mathcal{B}}} is
the $\cO$-submodule of the forms $F$ in \m{\cO^{\Fam}} such that
\m{(\letterRep{T}{n}{\gamma}F)(\beta)\in I} for every $\gamma$ in \m{\Lev_{\tau}(W)} and every $\beta$ in~$\mathcal{B}$.
\item For indices \m{i_0}, \m{i_1}, an ideal~$I$ of~$\cO$ and a real number~$h$,
\begin{gather*}
\exc{T}{n}{i_0}{i_1}{I}{h}=\trl{T}{n}{\twoBlock{k}{i_1}{\cdot}}{I}{\mathcal{B}},\qquad\text{where}\\
\mathcal{B}=\bigl\{\beta:\ h<\pr{\exw{T}{s_0}{i_0}{i_1}{\cdot}}{\beta}\ \text{for some letter } s_0\bigr\}.
\end{gather*}
\item For natural numbers $k$, $\thr$ and $j$,
\m{\deepForms{T}{k}{n}{\thr}{j}=\trl{T}{n}{0}{0}{\{\beta:\dep{T}{k}{n}{\thr}{\beta}<j\}}}, with the constant
labelling~$0$ and the zero ideal.
\end{enumerate}
\end{definition}
\leanline{\lean{translatesIn}\sep \lean{excessCond}\sep \lean{deepForms}}

So $F$ lies in \m{\exc{T}{n}{i_0}{i_1}{I}{h}} when, for every $\gamma$ that is block diagonal with the blocks
\m{\{i<i_1\}} and \m{\{i\ge i_1\}}, the form \m{\letterRep{T}{n}{\gamma}F} has modulo~$I$ only monomials whose
excess is at most~$h$ at every letter; and $F$ lies in \m{\deepForms{T}{k}{n}{\thr}{j}} when every translate of~$F$
under \m{\GLk(W)} has only monomials of depth at least~$j$.

\begin{definition}[filtrations]\label{d:filt}
Let $\cO$ be a ring.
\begin{enumerate}[label=\textup{(\alph*)},leftmargin=2.3em,itemsep=2pt]
\item Let \m{\cF=(\cF(\theta))_{\theta\in\RR}} be a family of additive subgroups of~$\cO$. We say that
\IsRF{\cF} if \m{1\in\cF(0)}, if \m{\cF(\theta)\cdot\cF(\theta')\subseteq\cF(\theta+\theta')} for all
\m{\theta,\theta'}, and if \m{\cF(\theta')\subseteq\cF(\theta)} whenever \m{\theta\le\theta'}.
\item For $\varpi$ in~$\cO$ and a real number~$\theta$, \m{\adic{\varpi}{\theta}} is the additive group of the
elements of~$\cO$ divisible by \m{\varpi^{\max(\lceil\theta\rceil,0)}}.
\item For a family $\cF$ of additive subgroups of~$\cO$, a set $\cV$ of variables, a weight~$\omega$ on~$\cV$ and
a real number~$\theta$, \m{\gauss{\cF}{\omega}{\theta}} is the additive group of the polynomials in
\m{\cO[x_{\cV}]} whose coefficient at each \m{x^{\alpha}} lies in \m{\cF(\theta-\pr{\omega}{\alpha})}.
\end{enumerate}
\end{definition}
\leanline{\lean{IsRF}\sep \lean{adic}\sep \lean{gauss}}

In the sense of~(a), \IsRF{\adic{\varpi}{\cdot}}, because
\m{\lceil\theta+\theta'\rceil\le\lceil\theta\rceil+\lceil\theta'\rceil}; and if \IsRF{\cF}, then
\IsRF{\gauss{\cF}{\omega}{\cdot}}, on the polynomial ring, because the coefficient of \m{x^{\alpha}} in a product is a
sum of products of coefficients at exponents with sum~$\alpha$. These are the two instances of the unit
(§\ref{s:notes}, note~1). In \m{\gauss{\adic{\varpi}{\cdot}}{\omega}{\theta}} a variable \m{x_v} sits at
level~\m{\omega(v)}: the level of a polynomial is the least value of \emph{valuation plus weight} over its terms.
For \m{\varpi=0} the group \m{\adic{0}{\theta}} is~$\cO$ for \m{\theta\le0} and $0$ for \m{\theta>0}.

\begin{definition}[cones and slabs]\label{d:cones}
Let \m{k\ge0}, \m{L\ge1}, let $\pa$ be a map from \m{\ZL} to the signs, $\varrho$ a real number, and
\m{\sigma=(\sigma_s)_{s\in\ZL}} permutations of the indices.
\begin{enumerate}[label=\textup{(\alph*)},leftmargin=2.3em,itemsep=2pt]
\item \m{\permStep{\pa}{\varrho}{\sigma}} is the linear endomorphism of \m{\RR^{k\times L}} sending $\xi$ to
the vector with the coordinates
\m{\pa(s)\cdot\bigl(\varrho\,\xi(i,s+1)-\xi(\sigma_s(i),s)\bigr)} at \m{(i,s)}. For the identity
permutations we write \m{\permStep{\pa}{\varrho}{1}}.
\item For reals \m{\delta=(\delta_s)_s}, \m{\slab{k}{\delta}} is the set of the $\eta$ in \m{\RR^{k\times L}} with
all coordinates \m{\ge0} and \m{\delta_s-1\le\sum_i\eta(i,s)\le\delta_s} for every letter~$s$.
\item \m{\pieri{\pa}{\varrho}} is the set of the $\xi$ in \m{\RR^{k\times L}} such that for every
letter~$s$ and all indices \m{i}, \m{i+1}: \m{\varrho\,\xi(i+1,s+1)\le\xi(i,s)} if \m{\pa(s)=+}, and
\m{\xi(i+1,s)\le\varrho\,\xi(i,s+1)} if \m{\pa(s)=-}.
\item In Setting~\ref{set:T} with $\cO$ a field, for an index~$i_1$ and a real number~$h$,
\m{\exCone{T}{i_1}{h}} is the set of the $\eta$ in \m{\RR^{k\times L}} such that for every letter~$s_0$ and all
indices \m{i}, \m{i+1}: \m{\pr{\exw{T}{s_0}{i}{i+1}{\cdot}}{\eta}\le h} if \m{i+1=i_1}, and
\m{\pr{\exw{T}{s_0}{i}{i+1}{\cdot}}{\eta}\le0} otherwise.
\end{enumerate}
\end{definition}
\leanline{\lean{permStep}\sep \lean{slab}\sep \lean{pieri}\sep \lean{exCone}}

With \m{\eta=\permStep{\pa}{\varrho}{1}(\xi)}, the condition of~(c) at an even letter, together with
\m{\eta(i,s)>0}, reads \m{\varrho\,\xi(i,s+1)>\xi(i,s)\ge\varrho\,\xi(i+1,s+1)}: the column \m{\xi(\cdot,s)}
interlaces the column \m{\varrho\,\xi(\cdot,s+1)} from below. At an odd letter the roles are exchanged. Hence
the name: a cone of cyclically interlacing columns.

\begin{definition}[the size conditions]\label{d:sizes}
For natural numbers $\mathbf{P}$, $o$, $L$, $k$, $n$, $\thr$, $\lambda$ and real numbers $\kappa$, $y$, $z$ we write
\m{\Sz{\mathbf{P},o,L,k,n,\thr}{\lambda,\kappa,y,z}} for the conjunction of the ten conditions
\begin{gather*}
0<\thr,\quad\ 2\thr\le k,\quad\ \kappa\le1,\quad\ 1\le y,\quad\ \lambda+\lambda=L,\quad\
L\,n\,k\,(\lambda+2)+n+k<\mathbf{P}^{\lambda},\\
6\bigl(2(k+k^{2})\,oL\cdot L+2\bigr)+2\le y,\qquad \frac{2(L+2)k}{y}\le\kappa,\\
\bigl(y+L(k+1)\bigr)\,k\,\mathbf{P}^{3\thr}+k+2\le\frac{n}{y},\qquad
L\,n\,\Bigl(\frac{1}{\mathbf{P}^{2\thr}}+\frac{\thr^{2}}{kP}\Bigr)\le z .
\end{gather*}
\end{definition}
\leanline{\lean{HeckeSizes}\sep in the formal text the ten arguments are called $P$, $d$, $N$, $k$, $n$, $r$, $\ell$,
$\kappa$, $y$, $z$}

\begin{definition}[Frobenius operators and conjugations]\label{d:frob}
In Setting~\ref{set:T}, let $\thr$, $\lambda$, $n$, $j$ be natural numbers, $s$ a letter, $\gamma$ an element of
\m{\GLk(W)}, and \m{i_0}, \m{i_1}, $i$ indices.
\begin{enumerate}[label=\textup{(\alph*)},leftmargin=2.3em,itemsep=2pt]
\item The \kw{generic variables} are \m{x_{s,i,i'}}, for a letter and two indices; \m{X_s=(x_{s,i,i'})_{i,i'}} is
the generic matrix of the letter~$s$. \m{\frobX{T}{s}{j}} is the matrix over the polynomial ring in the generic
variables over~$\cO$ whose \m{(i,i')} entry is \m{x_{s-j,i,i'}^{\,\PT{T}^{j}}}: the generic matrix of the letter
\m{s-j} with its entries raised to the power \m{\PT{T}^{j}}.
\item Let $\cR$ be a ring and \m{\psi=(\psi(s,j))} matrices in \m{\Mk(\cR)}, for $s$ a letter and $j$ a natural
number. For a letter~$s$ and a natural number~$m$ put \m{\Pi_s^{m}=\psi(s,m-1)\cdots\psi(s,1)\,\psi(s,0)}, and
let $B_s$ be the matrix in \m{\Mk(\cR)} whose row~$i$ is row~$i$ of \m{\Pi_s^{m}} if $s$ is even, and of the
transpose of the adjugate of \m{\Pi_s^{m}} if $s$ is odd, where \m{m=\lambda+\dpi{T}{\thr}{(i,s)}}. Then
\[\frob{T}{\thr}{\lambda}{n}{\psi}=\bigotimes_{s}\mathrm{Sym}^{n}(B_s),\]
a square matrix over~$\cR$ indexed by \m{\Fam}.
\item \m{\letterConj{T}{\gamma}{s}} is the multiplicative map
\m{G\mapsto\iota_s(\gamma)\,G\,\iota_s(\gamma^{-1})} of \m{\Mk(\cO)}.
\item \m{\idxw{T}{\lambda}{i_0}{i_1}{i}=\bigl([i=i_1]-[i=i_0]\bigr)/\PT{T}^{\lambda}}, a real number.
\item For a ring homomorphism \m{\chi\colon W\to\cO}, \m{\conjSubst{\chi}{\gamma}} is the endomorphism of the
$\cO$-algebra of polynomials in the entries \m{x_{i,i'}} of one generic matrix~$X$ that substitutes the entries
of \m{\chi(\gamma^{-1})\,X\,\chi(\gamma)} for those of~$X$.
\end{enumerate}
\end{definition}
\leanline{\lean{frobX}\sep \lean{frobOperator}\sep \lean{letterConj}\sep \lean{idxw}\sep \lean{conjSubst}}

If the generic matrix of every letter~$s$ is specialised to \m{\iota_s(\gamma)}, for one $\gamma$ in
\m{\GLk(W)}, then \m{\frobX{T}{s}{j}} becomes congruent to \m{\iota_s(\gamma)} modulo~$p$, and the matrices $B_s$ of~(b),
for \m{\psi=\frobX{T}{\cdot}{\cdot}}, become congruent to matrices whose rows are rows of
\m{\iota_s(\gamma)^{\lambda}} or of \m{\iota_s(\gamma)^{\lambda+1}} (for odd letters, of their transposed
adjugates). The operator on generic matrices keeps more: the exponents of the generic variables of each letter.

\statement{inksoft}{\textbf{Inputs from other units.} The results below are taken for granted in this note.

\vspace{4pt}\phantomsection\label{i:1}\textbf{Input 1} (index divisibility; unit 149).\enspace
\emph{Let $\cO$ be a commutative ring, $V$ and $V'$ modules over~$\cO$, let $M_1$, $M_2$, $M_3$ be submodules
of~$V$ and $M'_1$, $M'_2$ submodules of~$V'$, and let $\pi$, $\pi'$, $\pi_0$ be elements of~$\cO$. One writes
\m{\IndexDvd{M_1}{M_2}{\pi}}, the index of $M_2$ in $M_1$ divides~$\pi$, if for every natural number~$n'$ and every
family of $n'$ elements of~$M_1$ there is a square matrix of size~$n'$ over~$\cO$ whose determinant
divides~$\pi$ and whose rows, used as coefficients of combinations of the family, give $n'$ elements of~$M_2$.
For an element~$p$ of~$\cO$, a natural number~$e$, an $\cO$-linear endomorphism~$Y$ of~$V$ and a real
number~$\theta$ one says that \HasType{p}{M_1}{e}{Y}{\theta} if \m{Y(M_1)\subseteq M_1} and
\m{\IndexDvd{M_1}{Y(M_1)+p^{e}M_1}{p^{l}}} for some natural number \m{l\le\theta}. Then:}
\begin{enumerate}[label=\textup{(\alph*)},leftmargin=2.3em,itemsep=1pt,topsep=2pt]
\item \emph{if $\cO$ is a valuation domain, $\mathcal{J}$ a finite set, $G$ a square matrix over~$\cO$ indexed
by~$\mathcal{J}$ whose determinant divides~$\pi$, and $M_2$ a submodule of \m{\cO^{\mathcal{J}}} that contains
\m{G(\cO^{\mathcal{J}})}, then \m{\IndexDvd{\cO^{\mathcal{J}}}{M_2}{\pi}};}
\item \emph{if \m{\IndexDvd{M_1}{M_2}{\pi}} and \m{\IndexDvd{M_2}{M_3}{\pi'}}, then
\m{\IndexDvd{M_1}{M_3}{\pi\pi'}};}
\item \emph{if \m{\Psi\colon V\to V'} is an $\cO$-linear map, \m{\IndexDvd{M_1}{M_2}{\pi}},
\m{M'_1\subseteq\Psi(M_1)+M'_2} and \m{\Psi(M_2)\subseteq M'_2}, then \m{\IndexDvd{M'_1}{M'_2}{\pi}};}
\item \emph{if \m{Y\colon V\to V} is an $\cO$-linear map with \m{Y(M_1)\subseteq M_1}, if
\m{\IndexDvd{M_1}{Y(M_1)+\pi_0M_1}{\pi}} and \m{\IndexDvd{M_1}{M_2+\pi_0M_1}{\pi'}}, then
\m{\IndexDvd{M_1}{Y(M_2)+\pi_0M_1}{\pi\pi'}}.}
\end{enumerate}

\vspace{3pt}\phantomsection\label{i:2}\textbf{Input 2} (normal forms over a valuation ring; units 120 and~112,
background).\enspace
\emph{Let $\cA$ be a commutative ring and $\cR$ a subring such that \m{(\cA,\cR)} is a valuation pair: for any
two elements of~$\cA$ there is an idempotent of~$\cR$ such that the $\cR$-module spanned by the two elements
lies in the $\cR$-multiples of the first element minus the idempotent times the difference of the first and the
second. If $\cA$ is a valuation domain, then \m{(\cA,\cA)} is a valuation pair. Let $\cN$ be a finite set.}
\begin{enumerate}[label=\textup{(\alph*)},leftmargin=2.3em,itemsep=1pt,topsep=2pt]
\item \emph{\textup{(Hermite)} Let $\cK$ be a finite set with an injective map to~$\cN$, by which $\cK$ is
regarded as a subset of~$\cN$, let $X$ be a matrix over~$\cA$ with rows indexed by~$\cN$ and columns by~$\cK$, and
let weights, natural numbers, be given on~$\cK$. Then there are an invertible matrix~$U$ over~$\cA$ indexed
by~$\cN$, with the entries of $U$ and of \m{U^{-1}} in~$\cR$, and a square matrix~$\Delta$ over~$\cA$ indexed
by~$\cK$, whose entry at \m{(l,l')} vanishes whenever the weight of~$l'$ is less than that of~$l$, such that
\m{X=U_{\cdot,\cK}\,\Delta} and \m{(U^{-1})_{\cK,\cdot}\,X=\Delta}; here \m{U_{\cdot,\cK}} consists of the columns
of~$U$ in~$\cK$ and \m{(U^{-1})_{\cK,\cdot}} of the rows of \m{U^{-1}} in~$\cK$.}
\item \emph{\textup{(Smith)} Every square matrix $X$ over~$\cA$ indexed by~$\cN$ is a product
\m{U\operatorname{diag}(\varepsilon_l)_{l\in\cN}\,U'} with $U$, $U'$ invertible, the entries of $U$, \m{U^{-1}},
$U'$, \m{U'^{-1}} in~$\cR$, and \m{\varepsilon_l} in~$\cA$.}
\end{enumerate}
(The reading of the Hermite form by program also lists a ring with involution, from an earlier part of its
file; the statement does not involve it.)}
\leanline{\lean{IndexDvd}\sep \lean{HasType}\sep \lean{ofdetdvd}\sep \lean{idtrans}\sep \lean{idmap}\sep
\lean{idcomp}\sep \lean{hermite}\sep \lean{smith}\sep \lean{valtop}}

\statement{inksoft}{\phantomsection\label{i:3}\textbf{Input 3} (a combinatorial Nullstellensatz with
multiplicity; unit 57, background).\enspace
\emph{Let $\FF$ be a field, $\cV$ and $\cS$ finite sets, and $\varkappa$ a map from $\cV$ to the natural numbers
with \m{0<\varkappa(v)} for every~$v$; the \kw{cost} of an exponent~$\alpha'$ on~$\cV$ is
\m{\sum_v\alpha'(v)/\varkappa(v)}. Let \m{(\upsilon_{\fs})_{\fs\in\cS}} be pairwise distinct points of
\m{\FF^{\cV}} with \m{\prod_v\varkappa(v)\le\#\cS}. Suppose that for every $v$ in~$\cV$ there is a polynomial
\m{\Phi_v} in \m{\FF[x_{\cV}]}, all of whose monomials \m{x^{\alpha'}} have cost less than~$1$, with
\m{\Phi_v(\upsilon_{\fs})=\upsilon_{\fs}(v)^{\varkappa(v)}} for every~$\fs$. Let $j$ be a natural number and
$\Phi$ a nonzero polynomial in \m{\FF[x_{\cV}]} that lies, for every~$\fs$, in the $j$-th power of the ideal of
the polynomials vanishing at~\m{\upsilon_{\fs}}. Then $\Phi$ has a monomial \m{x^{\alpha}}, with nonzero
coefficient, of cost at least~$j$.} (The reading by program also lists a polynomial and an exponent of an earlier
part of the file; they are the bound objects \m{\Phi_v} and \m{x^{\alpha}} of the statement.)

\vspace{3pt}\phantomsection\label{i:4}\textbf{Input 4} (the order of the general linear group; unit 24,
background).\enspace
\emph{For a finite field~$\FF$ and a finite set of \m{k'} elements, the group of the invertible matrices
over~$\FF$ indexed by the set has \m{\prod_{l=0}^{k'-1}(\#\FF)^{l}\bigl((\#\FF)^{l+1}-1\bigr)} elements; the formal
text states this as an equation in an arbitrary field.}}
\leanline{\lean{nullst}\sep \lean{cardGL}}

Input~1 is the subject of the note on unit~149 of this write-up; there the definitions are stated with their
arguments in the formal order and the four properties are proved. Input~2(a) is also an input of that note (its
Input~4), where it is not proved either, and the fact that a valuation domain gives a valuation pair has a proof
of two lines there; Input~2(b), the Smith form, comes from unit~112. The background library also supplies the
vocabulary of the conventions (linear substitutions, symmetric powers, Kronecker products, Levi and parahoric
subgroups; units 57 and~112).

\emph{References for the Inputs.} The catalogue of quoted results of this write-up lists the Hermite form as an
elementary lemma with no confirmed place in the literature; it has no entry yet for the Smith form, for Input~3
or for Input~4. What can be said of each:
\begin{itemize}[leftmargin=1.6em,itemsep=2pt,topsep=2pt]
\item \emph{Input~2.} This note uses both forms only for a valuation domain~$\cO$ and the pair \m{(\cO,\cO)}
(Propositions \ref{p:pairing} and~\ref{p:depthfact}). There they are classical: among the entries of a nonzero
matrix over a valuation domain one divides all the others, and operations on rows and columns move it to a corner
and clear its row and its column. The rings over which every matrix has a triangular or a diagonal reduction are
the subject of I.~Kaplansky, \emph{Elementary divisors and modules}, Trans.\ Amer.\ Math.\ Soc.\ 66 (1949),
464--491, and valuation rings are among them; the place in that paper was not confirmed. For the form with a
valuation pair, which the artifact needs elsewhere: \emph{reference to be confirmed}.
\item \emph{Input~3.} The artifact calls it a known result, a combinatorial Nullstellensatz with multiplicity (a
nonzero polynomial that vanishes to high order on a large grid-like set of points has a term of large cost), and
the header of its file names as sources N.~Alon, \emph{Combinatorial Nullstellensatz}, Combin.\ Probab.\ Comput.\
8 (1999), 7--29, and S.~Ball and O.~Serra, \emph{Punctured combinatorial Nullstellens\"atze}, Combinatorica 29
(2009), 511--522. Those papers concern the zeros of a polynomial on a product of finite subsets of the field, the
second with multiplicity. Input~3 replaces the product by any set of points on which each power
\m{x_v^{\varkappa(v)}} agrees with a polynomial of cost less than~$1$; this note did not compare the statements,
and for the form stated here: \emph{reference to be confirmed}.
\item \emph{Input~4} is classical, and its proof is one line: the column of index~$l$ of an invertible matrix,
\m{l=0,\dots,k'-1}, is any vector outside the span of the $l$ columns before it, so there are
\m{\prod_{l=0}^{k'-1}\bigl((\#\FF)^{k'}-(\#\FF)^{l}\bigr)} invertible matrices, which is the product of Input~4. A
place in a textbook: \emph{reference to be confirmed}.
\end{itemize}

\sect{First properties}{s:first}
\begin{lemma}[elementary properties]\label{l:basic}
In Setting~\ref{set:T}:
\begin{enumerate}[label=\textup{(\alph*)},leftmargin=2.3em,itemsep=1pt,topsep=2pt]
\item for every natural number~$j$, every letter~$s$ and every $\upsilon$ in~$W$, $p$ divides
\m{\iota_{s+j}(\upsilon)-\iota_s(\upsilon)^{\PT{T}^{j}}};
\item for $\gamma$ in \m{\GLk(W)} and a variable \m{v=(i',s)},
\m{\polyAction{T}{\gamma}x_{v}=\sum_{i}(\letterMatrix{T}{\gamma}{s})_{ii'}\,x_{i,s}};
\item for a natural number~$n$, $\gamma$ in \m{\GLk(W)} and a form $F$ in \m{\cO^{\Fam}},
\m{\polyAction{T}{\gamma}} sends the polynomial \m{\sum_{\beta}F(\beta)x^{\beta}} to the polynomial of the form
\m{\letterRep{T}{n}{\gamma}F};
\item for a ring homomorphism \m{\chi\colon\cO\to\cO'}, a natural number~$n$ and $\gamma$ in \m{\GLk(W)}, the
matrix \m{\letterRep{\Tmap{\chi}{T}}{n}{\gamma}} is obtained by applying $\chi$ to the entries of
\m{\letterRep{T}{n}{\gamma}};
\item for natural numbers $n$, $\thr$ and an element $\pi$ of~$\cO$, the diagonal matrix with the entries
\m{\pi^{\dep{T}{k}{n}{\thr}{\beta}}}, \m{\beta\in\Fam}, equals
\m{\bigotimes_s\mathrm{Sym}^{n}\bigl(\operatorname{diag}(\pi^{\dpi{T}{\thr}{(i,s)}})_{i}\bigr)}.
\end{enumerate}
Moreover \m{1<\PT{T}} and \m{1<\qT{T}}; the residue field of~$W$ is finite; and \m{\varpi\ne0} if $\cO$ has
characteristic zero.
\end{lemma}
\begin{proof}
(a) By induction on~$j$: if \m{\iota_{s+j}(\upsilon)\equiv\iota_s(\upsilon)^{\PT{T}^{j}}} modulo~$p$, then, modulo~$p$,
\[\iota_{s+j+1}(\upsilon)\equiv\iota_{s+j}(\upsilon)^{\PT{T}}\equiv\iota_s(\upsilon)^{\PT{T}^{j+1}},\]
because congruent elements have congruent powers. (b) is the definition of the linear substitution. (c) By~(b)
and the definition of \m{\mathrm{Sym}^{n}}, \m{\polyAction{T}{\gamma}} sends the monomial \m{\prod_sx^{\beta_s}}, with
\m{x^{\beta_s}} in the variables of the letter~$s$, to
\[\prod_s\sum_{\beta'_s}\mathrm{Sym}^{n}(\letterMatrix{T}{\gamma}{s})_{\beta'_s,\beta_s}x^{\beta'_s}
=\sum_{\beta'}\letterRep{T}{n}{\gamma}_{\beta',\beta}\,x^{\beta'}.\]
(d) The entries of \m{\mathrm{Sym}^{n}(G)} are
polynomials with integer coefficients in the entries of~$G$. (e) The linear substitution of a diagonal matrix
multiplies each monomial by the product of the diagonal entries to the powers of the exponent. The last
assertions: \m{\PT{T}\ge5}; \m{\#\FF_W=\qT{T}} is a positive natural number; and \m{\varpi^{2}=p\ne0}.
\end{proof}
\leanline{\lean{frobpow}\sep \lean{polyActionX}\sep \lean{polyOfCoeffs}\sep \lean{letterRepmap}\sep
\lean{diagdepth}\sep \lean{oneltP}\sep \lean{oneltq}\sep \lean{finres}\sep \lean{wnezero}}

\begin{lemma}[entries of parahoric elements]\label{l:par}
In Setting~\ref{set:T}, let $\tau$ be a labelling, $\gamma$ an element of \m{\Par_{\tau}(W)}, $s$ a letter, and
$i$, $i'$ indices such that \m{\tau(i)<\tau(i')} if $s$ is even and \m{\tau(i')<\tau(i)} if $s$ is odd. Then $p$
divides the entry \m{(\letterMatrix{T}{\gamma}{s})_{ii'}} in~$\cO$.
\end{lemma}
\begin{proof}
Let $\upsilon$ be the \m{(i,i')} entry of $\gamma$ if $s$ is even and of \m{(\gamma^{-1})^{\tp}} if $s$ is odd, so
that the entry in question is \m{\iota_s(\upsilon)}. In both cases $\upsilon$ lies in~$\fm_W$: for even~$s$ by
the definition of the parahoric subgroup, for odd~$s$ because \m{\gamma^{-1}} lies in it and
\m{\upsilon=(\gamma^{-1})_{i'i}}. Hence \m{1-\upsilon^{\qT{T}-1}} is a unit of~$W$ and its image under~$\iota_s$ a
unit of~$\cO$. By Lemma~\ref{l:basic}(a) with \m{j=L}, where \m{s+L=s} and \m{\PT{T}^{L}=\qT{T}}, the prime $p$
divides \m{\iota_s(\upsilon)-\iota_s(\upsilon)^{\qT{T}}=\iota_s(\upsilon)\cdot\iota_s(1-\upsilon^{\qT{T}-1})}, hence
\m{\iota_s(\upsilon)}.
\end{proof}
\leanline{\lean{pdvdpar}}

\begin{lemma}[size of the excess weights]\label{l:exIcc}
In Setting~\ref{set:T}, for a letter~$s_0$, indices \m{i_0}, \m{i_1} and a variable~$v$:
\m{-\tfrac15\le\exw{T}{s_0}{i_0}{i_1}{v}\le1}.
\end{lemma}
\begin{proof}
Let \m{v=(i,s)}. As \m{[s-s_0]\le L-1} and \m{\PT{T}\ge5},
\[5\,\PT{T}^{[s-s_0]}\le\qT{T};\]
and \m{1+[s-1-s_0]} is $L$ if \m{s=s_0} and \m{[s-s_0]} otherwise. So the first term of the definition is $0$, or
$1$, or lies in \m{(0,\tfrac15]}, and the second is $0$ or lies in \m{[-\tfrac15,0)}.
\end{proof}
\leanline{\lean{exwIcc}\sep \lean{oneaddval}\sep \lean{fivemul}}

\begin{lemma}[half of the letters are deep]\label{l:sumdeep}
In Setting~\ref{set:T}, let $\thr$ be a natural number and $i$ an index. Then
\m{\sum_{s}\dpi{T}{\thr}{(i,s)}=L/2}. (The formal statement asks that $\cO$ be a field and does not use it; it
writes the quotient of natural numbers.)
\end{lemma}
\begin{proof}
The parity alternates, so \m{\dpi{T}{\thr}{(i,s)}+\dpi{T}{\thr}{(i,s+1)}=1} for every letter; summing over the
letters gives twice the sum on the left and $L$ on the right.
\end{proof}
\leanline{\lean{sumdeep}}

\sect{The root lemma}{s:root}
\begin{lemma}[a count of zeros against degree]\label{l:stepanov}
Let $W$ be a local ring, $\cO$ a ring, \m{\chi\colon W\to\cO} a ring homomorphism, $\Phi$ a nonzero polynomial in
one variable~$X$ over~$\cO$, and $e$ a natural number. Suppose that \m{(X-\chi(\upsilon))^{e}} divides~$\Phi$ for
every $\upsilon$ in~$W$; equivalently, \m{X^{e}} divides \m{\Phi(X+\chi(\upsilon))}. Then
\m{e\cdot\#\FF_W\le\deg\Phi}.
\end{lemma}
\begin{proof}
If \m{\FF_W} is infinite the count on the left is~$0$. Otherwise choose one element of~$W$ in each residue class.
The difference of two of them is a unit of~$W$, so its image is a unit of~$\cO$, and the polynomials
\m{(X-\chi(\upsilon))^{e}} for the chosen~$\upsilon$ are pairwise coprime. Each divides~$\Phi$; so does their
product, which is monic of degree \m{e\cdot\#\FF_W}. A nonzero polynomial is not divisible by a monic polynomial
of larger degree.
\end{proof}
\leanline{\lean{stepanov}}

\begin{lemma}[filtrations under substitution]\label{l:subst}
Let $\cO$ be a ring, $\cA$ a commutative $\cO$-algebra, $\cV$ a set of variables, \m{\cF=(\cF(\theta))_{\theta}} a
family of additive subgroups of~$\cO$, $\omega$ a weight on~$\cV$ and $\theta$ a real number. Let $\chi$ and
$\chi'$ be homomorphisms of $\cO$-algebras from \m{\cO[x_{\cV}]} to~$\cA$, and let \m{\cG=(\cG(\theta))_{\theta}} be
a family of additive subgroups of~$\cA$ such that \IsRF{\cG}. Let $\Phi$ be a polynomial in
\m{\gauss{\cF}{\omega}{\theta}}. Suppose that
\begin{enumerate}[label=\textup{(\roman*)},leftmargin=2.3em,itemsep=1pt,topsep=2pt]
\item \m{\chi(x_v)\in\cG(\omega(v))} for every $v$ in~$\cV$;
\item \m{\chi'(x_v)-\chi(x_v)\in\cG(\omega(v)+\tfrac1{20})} for every $v$ in~$\cV$;
\item the image in~$\cA$ of every element of \m{\cF(\theta')} lies in \m{\cG(\theta')}, for every real~$\theta'$.
\end{enumerate}
Then \m{\chi'(\Phi)\in\cG(\theta)} and \m{\chi'(\Phi)-\chi(\Phi)\in\cG(\theta+\tfrac1{20})}.
\end{lemma}
\begin{proof}
It suffices to treat one term \m{\pi x^{\alpha}} with $\pi$ in \m{\cF(\theta-\pr{\omega}{\alpha})}. By (i) and~(ii)
and since $\cG$ is decreasing, \m{\chi'(x_v)} and \m{\chi(x_v)} lie in \m{\cG(\omega(v))}; so \m{\chi'(x^{\alpha})}
lies in \m{\cG(\pr{\omega}{\alpha})}, and with~(iii) \m{\chi'(\pi x^{\alpha})} lies in \m{\cG(\theta)}. The
difference \m{\chi'(x^{\alpha})-\chi(x^{\alpha})}, expanded as a product of sums
\m{\chi(x_v)+(\chi'(x_v)-\chi(x_v))} less the product of the first summands, is a sum of products of
\m{|\alpha|} factors of which at least one is a difference; each such product lies in
\m{\cG(\pr{\omega}{\alpha}+\tfrac1{20})}.
\end{proof}
\leanline{\lean{aevalsub}\sep used with it: \lean{memadic} (an element lies in \m{\adic{\varpi}{\theta}} if it is $0$ or
\m{\theta\le0}), \lean{monomialmem} (a term \m{\pi x^{\alpha}} with $\pi$ in \m{\cF(\theta-\pr{\omega}{\alpha})} lies in
\m{\gauss{\cF}{\omega}{\theta}}), \lean{Xmem} (\m{x_v} lies in \m{\gauss{\adic{\varpi}{\cdot}}{\omega}{\omega(v)}})}

\begin{proof}[Proof of Theorem 1]
Throughout, $s$ runs over the letters, and we write
\[x^{+}_s=x_{\rootVar{T}{i_0}{i_1}{s}},\qquad x^{-}_s=x_{\rootVar{T}{i_1}{i_0}{s}},\]
and \m{\omega^{\pm}_s}, \m{\alpha'^{\pm}_s} for the values of $\omega$ and of an exponent~$\alpha'$ at these two
variables. For every letter the two variables differ, as \m{i_0\ne i_1}. Put
\[\mathrm{w}^{+}(\alpha')=\qT{T}\,\alpha'^{+}_{s_0}+\sum_{s\ne s_0}\PT{T}^{[s-s_0]}\alpha'^{+}_s,\qquad
\mathrm{w}^{-}(\alpha')=\qT{T}\,\alpha'^{-}_{s_0}+\sum_{s\ne s_0}\qT{T}\,\PT{T}^{[s-s_0]}\alpha'^{-}_s .\]
By~\eqref{e:excess}, \m{\qT{T}^{2}\cdot\pr{\exw{T}{s_0}{i_0}{i_1}{\cdot}}{\alpha}=\qT{T}\,\mathrm{w}^{+}(\alpha)-\mathrm{w}^{-}(\alpha)},
so the claim is \m{\qT{T}\,\mathrm{w}^{+}(\alpha)\le\mathrm{w}^{-}(\alpha)}. Put
\m{\epsilon=(\omega^{+}_{s_0}-\omega^{-}_{s_0})/\qT{T}}; then \m{\epsilon>0}. The hypothesis on~$\omega$, with
\m{\PT{T}^{[s-s_0]}\ge5} and \m{5\,\PT{T}^{[s-s_0]}\le\qT{T}} for \m{s\ne s_0}, gives for these letters
\begin{equation}\label{e:consts}
\epsilon\,\qT{T}\,\PT{T}^{[s-s_0]}\ \ge\ 2,\qquad \epsilon\,\PT{T}^{[s-s_0]}\ \le\ 0.32,\qquad
0.4\le\omega^{+}_s-\omega^{-}_s\le1.6 .
\end{equation}

\emph{Step 1: the elements of the Levi subgroup.} For $\upsilon$ in~$W$ let \m{\gamma_{\upsilon}} be the product of
the permutation matrix that exchanges \m{i_0} and~\m{i_1} with the elementary matrix
\m{1\mp\upsilon E_{i_0i_1}}, the sign being $-$ if $s_0$ is even. It lies in \m{\Lev_{\tau}(W)} because
\m{\tau(i_0)=\tau(i_1)}. Let \m{\mathrm{sg}_s=-1} if $s$ has the parity of~$s_0$ and \m{\mathrm{sg}_s=+1} otherwise.
A computation with Lemma~\ref{l:basic}(b), separately for even and odd letters, gives
\begin{equation}\label{e:action}
\polyAction{T}{\gamma_{\upsilon}}x^{+}_s=x^{-}_s+\mathrm{sg}_s\,\iota_s(\upsilon)\,x^{+}_s,\qquad
\polyAction{T}{\gamma_{\upsilon}}x^{-}_s=x^{+}_s ,
\end{equation}
and \m{\polyAction{T}{\gamma_{\upsilon}}} fixes the other variables.

\emph{Step 2: two substitutions.} Let \m{\cA=\cO[x_v][X]}, with one further variable~$X$. Let \m{\chi_0} be the
homomorphism of $\cO$-algebras from \m{\cO[x_v]} to~$\cA$ with
\[\chi_0(x^{+}_{s_0})=x^{-}_{s_0},\qquad\chi_0(x^{-}_s)=x^{-}_s\,X^{\mathrm{w}^{-}(s)}\ \text{ for all } s,\qquad
\chi_0(x_v)=x_v\ \text{ otherwise},\]
where \m{\mathrm{w}^{-}(s)=\qT{T}\,\PT{T}^{[s-s_0]}} is the coefficient of \m{\alpha'^{-}_s} in
\m{\mathrm{w}^{-}(\alpha')}; and let \m{\chi_1} agree with \m{\chi_0} except that, for \m{s\ne s_0},
\[\chi_1(x^{+}_s)=x^{+}_s-\mathrm{sg}_s\,x^{-}_s,\qquad
\chi_1(x^{-}_s)=x^{-}_s\,X^{\mathrm{w}^{-}(s)}-\mathrm{sg}_s\,x^{+}_s\,X^{\PT{T}^{[s-s_0]}}.\]
For an exponent~$\alpha'$ let \m{\alpha'^{*}} be the exponent obtained by moving the value at
\m{\rootVar{T}{i_0}{i_1}{s_0}} onto \m{\rootVar{T}{i_1}{i_0}{s_0}}; then
\m{\chi_0(x^{\alpha'})=x^{\alpha'^{*}}X^{\mathrm{w}^{-}(\alpha')}}.

On~$\cA$ consider two families of additive subgroups. An element lies in \m{\cG_{-}(\theta)}, respectively in
\m{\cG_{+}(\theta)}, if its coefficient at each \m{x^{\alpha'}X^{j}} is divisible by~$\varpi$ to the power
\begin{align*}
&\max\bigl(\lceil\theta-\pr{\omega}{\alpha'}-\epsilon\,\mathrm{w}^{-}(\alpha')+\epsilon j\rceil,0\bigr)
&&\text{for } \cG_{-}(\theta),\\
&\max\bigl(\lceil\theta-\pr{\omega}{\alpha'}+\epsilon\,\mathrm{w}^{+}(\alpha')-\epsilon j\rceil,0\bigr)
&&\text{for } \cG_{+}(\theta).
\end{align*}
Both are of the
kind of Definition~\ref{d:filt}(c), with a weight on the variables \m{x_v} and~$X$, so each has the three properties of
Definition~\ref{d:filt}(a), and each contains the image of \m{\adic{\varpi}{\theta}} at level~$\theta$. Since
\m{\epsilon\qT{T}=\omega^{+}_{s_0}-\omega^{-}_{s_0}}, moving the exponent does not change the two weights:
\begin{equation}\label{e:star}
\pr{\omega}{\alpha'^{*}}+\epsilon\,\mathrm{w}^{-}(\alpha'^{*})=\pr{\omega}{\alpha'}+\epsilon\,\mathrm{w}^{-}(\alpha'),\qquad
\pr{\omega}{\alpha'^{*}}-\epsilon\,\mathrm{w}^{+}(\alpha'^{*})=\pr{\omega}{\alpha'}-\epsilon\,\mathrm{w}^{+}(\alpha') .
\end{equation}

\emph{Step 3: the degree.} Put \m{\theta_0=e+\pr{\omega}{\alpha}}. We apply Lemma~\ref{l:subst} to $\Phi$, which
lies in \m{\gauss{\adic{\varpi}{\cdot}}{\omega}{\theta_0}} by the hypothesis for \m{\gamma=1}, with
\m{\chi=\chi_0}, \m{\chi'=\chi_1} and \m{\cG=\cG_{-}}. A term \m{x_{v'}X^{j}} lies in \m{\cG_{-}(\theta)} when
\m{\theta\le\omega(v')+\epsilon(\mathrm{w}^{-}(v')-j)}, with \m{\mathrm{w}^{-}(v')} the coefficient of~$v'$.
Condition~(i) holds with equality in each case, for \m{x^{+}_{s_0}} by the choice of~$\epsilon$. Condition~(ii)
concerns the letters \m{s\ne s_0}: for \m{x^{+}_s} it reads
\m{\omega^{+}_s+\tfrac1{20}\le\omega^{-}_s+\epsilon\,\qT{T}\PT{T}^{[s-s_0]}}, and for \m{x^{-}_s} it reads
\m{\omega^{-}_s+\tfrac1{20}\le\omega^{+}_s-\epsilon\,\PT{T}^{[s-s_0]}}; both follow from~\eqref{e:consts}. The
lemma gives
\begin{equation}\label{e:AD}
\chi_1(\Phi)\in\cG_{-}(\theta_0),\qquad\chi_1(\Phi)-\chi_0(\Phi)\in\cG_{-}(\theta_0+\tfrac1{20}).
\end{equation}

\emph{Step 4: the zeros.} Fix $\upsilon$ in~$W$, write \m{\pi=\iota_{s_0}(\upsilon)}, and let \m{\mathrm{t}_{\upsilon}}
be the automorphism \m{X\mapsto X+\pi} of~$\cA$. By hypothesis the polynomial
\m{\Phi_{\upsilon}=\polyAction{T}{\gamma_{\upsilon}^{-1}}\Phi} lies in \m{\gauss{\adic{\varpi}{\cdot}}{\omega}{\theta_0}},
and \m{\mathrm{t}_{\upsilon}(\chi_1(\Phi))=\chi_{\upsilon}(\Phi_{\upsilon})} for the homomorphism
\m{\chi_{\upsilon}=\mathrm{t}_{\upsilon}\circ\chi_1\circ\polyAction{T}{\gamma_{\upsilon}}}. We apply
Lemma~\ref{l:subst} with \m{\chi=\chi'=\chi_{\upsilon}} and \m{\cG=\cG_{+}}, and have to check that
\m{\chi_{\upsilon}(x_v)} lies in \m{\cG_{+}(\omega(v))} for every variable. A term \m{\pi'x_{v'}X^{j}} lies in
\m{\cG_{+}(\theta)} when $\varpi$ to the power
\m{\max(\lceil\theta-\omega(v')+\epsilon(\mathrm{w}^{+}(v')-j)\rceil,0)} divides~$\pi'$, with
\m{\mathrm{w}^{+}(v')} the coefficient of~$v'$ in~\m{\mathrm{w}^{+}}. By
Lemma~\ref{l:basic}(a), \m{\iota_s(\upsilon)\equiv\pi^{\PT{T}^{[s-s_0]}}} and \m{\pi^{\qT{T}}\equiv\pi}
modulo~$p$, and \m{(X+\pi)^{\PT{T}^{j}}\equiv X^{\PT{T}^{j}}+\pi^{\PT{T}^{j}}} modulo~$p$ because \m{\PT{T}} is a
power of~$p$. From \eqref{e:action} and the definition of~\m{\chi_1} one finds, with \m{P'=\PT{T}^{[s-s_0]}}:
\begin{align*}
\chi_{\upsilon}(x^{+}_{s_0})&=x^{-}_{s_0}\bigl((X+\pi)^{\qT{T}}-\pi\bigr)=x^{-}_{s_0}\bigl(X^{\qT{T}}+p\,(\dots)\bigr),\qquad
\chi_{\upsilon}(x^{-}_{s_0})=x^{-}_{s_0},\\
\chi_{\upsilon}(x^{+}_s)&=x^{-}_s\bigl((X+\pi)^{\qT{T}P'}-\iota_s(\upsilon)\bigr)-\mathrm{sg}_s\,x^{+}_s\bigl((X+\pi)^{P'}-\iota_s(\upsilon)\bigr)\\
&=x^{-}_s\bigl(X^{\qT{T}P'}+p\,(\dots)\bigr)-\mathrm{sg}_s\,x^{+}_s\bigl(X^{P'}+p\,(\dots)\bigr)\qquad(s\ne s_0),\\
\chi_{\upsilon}(x^{-}_s)&=x^{+}_s-\mathrm{sg}_s\,x^{-}_s\qquad(s\ne s_0),
\end{align*}
where the dots are polynomials in~$X$ of degree less than that of the leading power, and \m{\chi_{\upsilon}}
fixes the other variables. In the third line \m{\mathrm{sg}_s^{2}=1} was used. Now each term is checked against
the criterion. The leading terms: \m{x^{-}_{s_0}X^{\qT{T}}} is at level \m{\omega^{-}_{s_0}+\epsilon\qT{T}=\omega^{+}_{s_0}};
\m{x^{+}_sX^{P'}} is at level \m{\omega^{+}_s}; \m{x^{-}_sX^{\qT{T}P'}} is at level
\m{\omega^{-}_s+\epsilon\qT{T}P'\ge\omega^{+}_s} by~\eqref{e:consts}; \m{x^{+}_s} is at level
\m{\omega^{+}_s-\epsilon P'\ge\omega^{-}_s} by~\eqref{e:consts}, and \m{x^{-}_s} and \m{x^{-}_{s_0}} are at their own
levels. The terms with the factor \m{p=\varpi^{2}}
need the power of~$\varpi$ to be at most~$2$: for \m{x^{-}_{s_0}X^{j}} it is at most
\m{\lceil\omega^{+}_{s_0}-\omega^{-}_{s_0}\rceil\le2}; for \m{x^{-}_sX^{j}} at most
\m{\lceil\omega^{+}_s-\omega^{-}_s\rceil\le2}; for \m{x^{+}_sX^{j}} with \m{j<P'} at most
\m{\lceil\epsilon P'\rceil\le1}. So the lemma applies:
\begin{equation}\label{e:S}
\mathrm{t}_{\upsilon}(\chi_1(\Phi))\in\cG_{+}(\theta_0)\qquad\text{for every } \upsilon \text{ in } W.
\end{equation}

\emph{Step 5: the count.} Let \m{\Phi^{*}}, a polynomial in~$X$ over~$\cO$, be the coefficient of
\m{x^{\alpha^{*}}} in \m{\chi_1(\Phi)}. By \eqref{e:AD} and~\eqref{e:star} the coefficient of \m{X^{j}}
in~\m{\Phi^{*}} is divisible by~$\varpi$ to the power \m{\lceil e+\epsilon(j-\mathrm{w}^{-}(\alpha))\rceil}, which is
at least \m{e+1} for \m{j>\mathrm{w}^{-}(\alpha)}. By the second part of~\eqref{e:AD}, the coefficient of
\m{X^{\mathrm{w}^{-}(\alpha)}} in \m{\Phi^{*}} differs by a multiple of \m{\varpi^{e+1}} from the coefficient of
\m{x^{\alpha^{*}}X^{\mathrm{w}^{-}(\alpha)}} in \m{\chi_0(\Phi)}. A term \m{x^{\alpha'}} of~$\Phi$ contributes to the
latter only if \m{\alpha'^{*}=\alpha^{*}} and \m{\mathrm{w}^{-}(\alpha')=\mathrm{w}^{-}(\alpha)}; then $\alpha'$ and
$\alpha$ agree outside the two root variables of~$s_0$, have the same sum there, and the same value at
\m{\rootVar{T}{i_1}{i_0}{s_0}}, so \m{\alpha'=\alpha}. Thus that coefficient is the coefficient of
\m{x^{\alpha}} in~$\Phi$, which \m{\varpi^{e+1}} does not divide. Finally, taking the coefficient of
\m{x^{\alpha^{*}}} commutes with~\m{\mathrm{t}_{\upsilon}}, and by \eqref{e:S} and~\eqref{e:star} the coefficient of
\m{X^{j}} in \m{\Phi^{*}(X+\iota_{s_0}(\upsilon))} is divisible by~$\varpi$ to the power
\m{\lceil e+\epsilon(\mathrm{w}^{+}(\alpha)-j)\rceil\ge e+1} for \m{j<\mathrm{w}^{+}(\alpha)}.

So the image of \m{\Phi^{*}} in \m{(\cO/\varpi^{e+1}\cO)[X]} is not zero, has degree at
most~\m{\mathrm{w}^{-}(\alpha)}, and is divisible by \m{(X-\iota_{s_0}(\upsilon))^{\mathrm{w}^{+}(\alpha)}} for every
$\upsilon$ in~$W$. Lemma~\ref{l:stepanov}, for the homomorphism from $W$ to \m{\cO/\varpi^{e+1}\cO} induced
by~\m{\iota_{s_0}}, gives \m{\mathrm{w}^{+}(\alpha)\cdot\#\FF_W\le\mathrm{w}^{-}(\alpha)}, and \m{\#\FF_W=\qT{T}}.
\end{proof}
\leanline{\lean{rootlemma}\sep \lean{stepanov}\sep \lean{aevalsub}}

\sect{The volume of the cyclic interlacing cone}{s:cone}
\begin{lemma}[the determinant of an interlacing pattern]\label{l:interlace}
Let \m{k\ge0}, and let \m{\xi_0\ge\dots\ge\xi_{k-1}} and \m{\eta_0\ge\dots\ge\eta_{k-1}} be weakly decreasing
sequences of real numbers. The determinant of the \m{k\times k} matrix whose \m{(i,i')} entry is $1$ if
\m{\eta_{i'}<\xi_i} and $0$ otherwise equals $1$ if \m{\eta_i<\xi_i} for all~$i$ and \m{\xi_{i+1}\le\eta_i} for all
\m{i<k-1}, and equals $0$ otherwise.
\end{lemma}
\begin{proof}
By induction on~$k$; for \m{k=0} both sides are~$1$. Let \m{k\ge1}. If \m{\eta_0<\xi_i} for some \m{i\ge1}, then
the rows $0$ and~$i$ consist of ones, since \m{\eta_{i'}\le\eta_0<\xi_i\le\xi_0}; the determinant is~$0$, and the
condition fails because \m{\xi_1\ge\xi_i>\eta_0}. Otherwise \m{\xi_i\le\eta_0} for all \m{i\ge1}, column~$0$ has at
most its entry in row~$0$ different from zero, and the determinant is that entry, $1$ or $0$ according as
\m{\eta_0<\xi_0} or not, times the determinant of the same kind for the sequences without their first members.
The induction hypothesis evaluates the latter, and \m{\xi_1\le\eta_0} holds.
\end{proof}
\leanline{\lean{detinter}}

\begin{proposition}[the volume as a signed sum]\label{p:signed}
Let $k$ and $m$ be natural numbers and \m{L=m+1}, let $\varrho$ be a real number with \m{1<\varrho}, $\pa$ a
map from \m{\ZL} to the signs, \m{\delta=(\delta_s)_s} real numbers, and $\cC$ a subset of \m{\RR^{k\times L}} such
that, for every~$\xi$: \m{\permStep{\pa}{\varrho}{1}(\xi)\in\cC} if and only if
\m{\xi\in\pieri{\pa}{\varrho}}. For a permutation~$\sigma'$ of the indices let \m{\sigma'^{*}} be the linear
map \m{f\mapsto f\circ\sigma'} of~\m{\RR^{k}} and
\m{\mathrm{dt}(\sigma')=|\det(1-\varrho^{-L}\sigma'^{*})|}. Then
\[\vol\bigl(\cC\cap\slab{k}{\delta}\bigr)=\frac{\mathrm{dt}(\mathrm{id})}{k!}\,
\Bigl(\sum_{\sigma'}\frac{\sgn(\sigma')}{\mathrm{dt}(\sigma')}\Bigr)\,\vol\bigl(\slab{k}{\delta}\bigr),\]
the sum being over all permutations of the indices. (The formal statement is for the letters \m{0,\dots,m} with
addition modulo \m{m+1}.)
\end{proposition}
\begin{proof}
For a tuple \m{\sigma=(\sigma_s)_s} of permutations write \m{S_{\sigma}=\permStep{\pa}{\varrho}{\sigma}}, and
\m{S_1} for the identity permutations. The set \m{\slab{k}{\delta}} is bounded, and so is its intersection with
any set; all volumes below are finite.

\emph{The determinant of \m{S_{\sigma}}.} Up to the signs \m{\pa(s)}, \m{S_{\sigma}} is
\m{\varrho\cdot\mathrm{Sh}\circ(1-\mathrm{N})}, where \m{(\mathrm{Sh}\,\xi)(i,s)=\xi(i,s+1)} and
\m{(\mathrm{N}\xi)(\cdot,s)=N_s\,\xi(\cdot,s-1)} with \m{N_s=\varrho^{-1}\sigma_{s-1}^{*}}. For such a cyclic
operator,
\[\det(1-\mathrm{N})=\det(1-N_{L-1}\cdots N_1N_0).\]
Indeed, order the letters as \m{0,\dots,L-1} and let $\mathrm{El}$ be the operator with
\[(\mathrm{El}\eta')(\cdot,s)=\sum_{s'\le s}N_sN_{s-1}\cdots N_{s'+1}\,\eta'(\cdot,s');\]
then
\[(\mathrm{El}(1-\mathrm{N})\xi)(\cdot,s)=\xi(\cdot,s)-N_s\cdots N_0\,\xi(\cdot,L-1).\]
The operator $\mathrm{El}$ has determinant~$1$, and \m{\mathrm{El}(1-\mathrm{N})} is triangular with the diagonal
blocks \m{1,\dots,1} and \m{1-N_{L-1}\cdots N_0}. Hence
\begin{equation}\label{e:detS}
|\det S_{\sigma}|=\varrho^{kL}\cdot\mathrm{dt}(\sigma_0\sigma_1\cdots\sigma_{L-1}),
\end{equation}
because \m{N_{L-1}\cdots N_0=\varrho^{-L}(\sigma_{L-1}\sigma_0\cdots\sigma_{L-2})^{*}} and \m{\mathrm{dt}} takes the same
value on conjugate permutations. As \m{\varrho>1}, the right side is not zero: \m{\mathrm{dt}(\sigma')} is the
product of the numbers \m{1-\varrho^{-Ll}} over the cycles of~$\sigma'$, of lengths~$l$.

\emph{Two families of sets.} Let $\mathcal{W}$ be the set of the $\xi$ such that the column sums
\m{\sum_i(S_1\xi)(i,s)} lie in \m{[\delta_s-1,\delta_s]} for every letter. The column sums of \m{S_{\sigma}\xi}
are those of \m{S_1\xi}, the permutations only reordering the summands. For a tuple~$\sigma$ let \m{\mathcal{W}_{\sigma}} be
the set of the $\xi$ in~$\mathcal{W}$ with all coordinates of \m{S_{\sigma}\xi} positive. It is the preimage under
\m{S_{\sigma}} of the slab without its walls, which have measure zero, so by~\eqref{e:detS}
\m{\vol(\mathcal{W}_{\sigma})=\varrho^{-kL}\,\mathrm{dt}(\sigma_0\cdots\sigma_{L-1})^{-1}\,\vol(\slab{k}{\delta})}. Let
$\Upsilon$ be the set of the $\xi$ in \m{\pieri{\pa}{\varrho}} with all coordinates of \m{S_1\xi} positive.
By the hypothesis on~$\cC$, \m{\mathcal{W}\cap\Upsilon} is the preimage under~\m{S_1} of \m{\cC\cap\slab{k}{\delta}} without
the walls, so \m{\vol(\mathcal{W}\cap\Upsilon)=\varrho^{-kL}\,\mathrm{dt}(\mathrm{id})^{-1}\,\vol(\cC\cap\slab{k}{\delta})}.
For a tuple \m{\tau'=(\tau'_s)_s} of permutations let \m{\xi\circ\tau'} be the vector with coordinates
\m{\xi(\tau'_s(i),s)}; the map \m{\xi\mapsto\xi\circ\tau'} preserves volume and maps $\mathcal{W}$ to itself.

\emph{A pointwise identity.} We claim that for every $\xi$
\begin{equation}\label{e:pointwise}
\#\{\tau':\ \xi\circ\tau'\in\Upsilon\}\ =\ \sum_{\sigma}\sgn(\sigma_0\cdots\sigma_{L-1})\cdot
\bigl[\text{all coordinates of } S_{\sigma}\xi \text{ are positive}\bigr].
\end{equation}
The right side is the product over the letters~$s$ of
\[\sum_{\sigma_s}\sgn(\sigma_s)\prod_i\bigl[\pa(s)\bigl(\varrho\,\xi(i,s+1)-\xi(\sigma_s(i),s)\bigr)>0\bigr],\]
that is, of the determinants of the matrices \m{A_s} whose entry at \m{(i,i')} is
\[\bigl[\pa(s)\bigl(\varrho\,\xi(i,s+1)-\xi(i',s)\bigr)>0\bigr].\]
Reorder the rows of~\m{A_s} so that
\m{\xi(\cdot,s+1)} becomes weakly decreasing, and its columns so that \m{\xi(\cdot,s)} does. This multiplies
\m{\det A_s} by the signs of two permutations, and in the product over the circle of letters each sign occurs
twice. After the reordering Lemma~\ref{l:interlace} applies, to the transposed matrix when \m{\pa(s)=-}:
the determinant is~$1$ if the decreasing rearrangements \m{\xi^{\downarrow}} of the columns satisfy, for all~$i$,
\begin{align*}
\varrho\,\xi^{\downarrow}(i,s+1)&>\xi^{\downarrow}(i,s)\ge\varrho\,\xi^{\downarrow}(i+1,s+1)&&\text{when } \pa(s)=+,\\
\xi^{\downarrow}(i,s)&>\varrho\,\xi^{\downarrow}(i,s+1)\ge\xi^{\downarrow}(i+1,s)&&\text{when } \pa(s)=-,
\end{align*}
and it is $0$ otherwise. These are the conditions that define~$\Upsilon$, at the letter~$s$. So the right side
of~\eqref{e:pointwise} is $1$ if \m{\xi^{\downarrow}\in\Upsilon} and $0$ otherwise. On the left, the columns of
a vector in~$\Upsilon$ are strictly decreasing, by the same inequalities; so \m{\xi\circ\tau'\in\Upsilon}
forces \m{\xi\circ\tau'=\xi^{\downarrow}} and determines~$\tau'$, the entries of each column being distinct. Hence
the left side is also $1$ if \m{\xi^{\downarrow}\in\Upsilon} and $0$ otherwise.

\emph{Integration.} Integrate~\eqref{e:pointwise} over~$\mathcal{W}$. The left side gives
\m{(k!)^{L}\vol(\mathcal{W}\cap\Upsilon)}, the right side
\m{\sum_{\sigma}\sgn(\sigma_0\cdots\sigma_{L-1})\vol(\mathcal{W}_{\sigma})}.
Each permutation~$\sigma'$ is the product \m{\sigma_0\cdots\sigma_{L-1}} for \m{(k!)^{L-1}} tuples. Inserting the
two volumes and cancelling \m{\varrho^{-kL}(k!)^{L-1}} gives the formula.
\end{proof}
\leanline{\lean{signedsum}\sep \lean{isBounded} (the intersection of any set with a slab is bounded)}

\begin{proof}[Proof of Theorem 2]
By Proposition~\ref{p:signed}, \m{\vol(\cC\cap\slab{k}{\delta})} is \m{\vol(\slab{k}{\delta})} times a
number that depends only on $k$ and on \m{q=\varrho^{L}}. We evaluate this number in the case of one letter,
\m{L=1}, with
the sign~$+$, the real number~$q$ in the place of~$\varrho$, \m{\delta=1} and \m{\cC=\pieri{+}{q}}, the set of the
$\xi$ in \m{\RR^{k}} with \m{q\,\xi(i+1)\le\xi(i)} for all~$i$; the hypothesis holds because \m{S_1\xi=(q-1)\xi}
and $\cC$ is a cone. Put \m{\mathrm{f}_l=(q^{l+1}-1)/(q-1)} and let $G$ be the linear endomorphism of \m{\RR^{k}} with
\m{(G\xi)(l)=\mathrm{f}_l\,(\xi(l)-q\,\xi(l+1))}, where \m{\xi(k)=0}. It is triangular with the diagonal
entries~\m{\mathrm{f}_l}, so \m{\det G=\prod_{l=0}^{k-1}\mathrm{f}_l>0}. It preserves the sum of the coordinates,
because \m{\mathrm{f}_l-q\,\mathrm{f}_{l-1}=1}. And
\m{G\xi} has all coordinates \m{\ge0} if and only if \m{\xi(l)\ge q\,\xi(l+1)} for all~$l$, which implies
\m{\xi\ge0} by descending induction. Hence the preimage of \m{\slab{k}{1}} under~$G$ is
\m{\cC\cap\slab{k}{1}}, and \m{\vol(\cC\cap\slab{k}{1})=\vol(\slab{k}{1})/\det G}. The volume of
\m{\slab{k}{1}} is finite and positive, as the set contains the cube \m{[0,1/k]^{k}}; only here is a volume
cancelled. So the number is \m{1/\prod_{l=1}^{k}\frac{q^{l}-1}{q-1}}, and the theorem follows, also when the
slab \m{\slab{k}{\delta}} has volume zero.
\end{proof}
\leanline{\lean{prodvol}}

The proof evaluates the signed sum of Proposition~\ref{p:signed} without computing it: the identity
\m{\frac{\mathrm{dt}(\mathrm{id})}{k!}\sum_{\sigma'}\sgn(\sigma')/\mathrm{dt}(\sigma')=\prod_{l=1}^{k}\frac{q-1}{q^{l}-1}} is
obtained from the geometry of one letter.

\sect{Two dimension bounds modulo \texorpdfstring{$p$}{p}}{s:dim}
\begin{lemma}[the number of flags]\label{l:borel}
Let $\FF$ be a finite field and \m{k\ge0}. The index in \m{\GLk(\FF)} of the subgroup of the lower triangular
matrices is at most \m{\prod_{l=1}^{k}\bigl((\#\FF)^{l}-1\bigr)/(\#\FF-1)}.
\end{lemma}
\begin{proof}
The subgroup has \m{(\#\FF-1)^{k}(\#\FF)^{k(k-1)/2}} elements, the lower triangular matrices with units on the
diagonal, and Input~4 gives the order of the group; the quotient is the product.
\end{proof}
\leanline{\lean{borel}\sep \lean{cardGL}}

\begin{lemma}[the excess of a step]\label{l:exdot}
In Setting~\ref{set:T} with $\cO$ a field, let \m{i_0}, \m{i_1} be indices, $s_0$ a letter and $\xi$ a vector in
\m{\RR^{k\times L}}. Put \m{\eta=\permStep{\pa_T}{\PT{T}}{1}(\xi)}. Then
\[\pr{\exw{T}{s_0}{i_0}{i_1}{\cdot}}{\eta}=\Bigl(1-\frac1{\qT{T}}\Bigr)\cdot
\Bigl(\eta\bigl(\rootVar{T}{i_0}{i_1}{s_0}\bigr)-\bigl(\xi(i_0,s_0)-\xi(i_1,s_0)\bigr)\Bigr).\]
\end{lemma}
\begin{proof}
Let \m{\mathrm{df}_s=\xi(i_0,s)-\xi(i_1,s)}, and write \m{\eta^{+}_s}, \m{\eta^{-}_s} for the values of~$\eta$ at
\m{\rootVar{T}{i_0}{i_1}{s}} and \m{\rootVar{T}{i_1}{i_0}{s}}. For every letter, even or odd, the definition of the
step gives $\eta^{+}_s-\eta^{-}_s=\mathrm{df}_s-\PT{T}\,\mathrm{df}_{s+1}$. By the definition of the excess weight, in
which the root variable of~$s_0$ has the coefficient \m{\PT{T}^{L}/\qT{T}=1}, the excess of~$\eta$ at~$s_0$ is
\[\Bigl(1-\frac1{\qT{T}}\Bigr)\eta^{+}_{s_0}+\frac1{\qT{T}}\sum_{s}\PT{T}^{[s-s_0]}\bigl(\eta^{+}_s-\eta^{-}_s\bigr)
=\Bigl(1-\frac1{\qT{T}}\Bigr)\eta^{+}_{s_0}+\frac1{\qT{T}}\sum_{s}\PT{T}^{[s-s_0]}\bigl(\mathrm{df}_s-\PT{T}\,\mathrm{df}_{s+1}\bigr),\]
and the last sum, taken around the circle from~$s_0$, telescopes to
\[\mathrm{df}_{s_0}-\PT{T}^{L}\mathrm{df}_{s_0}=(1-\qT{T})\,\mathrm{df}_{s_0}.\qedhere\]
\end{proof}
\leanline{\lean{exdot}}

By this lemma, \m{\permStep{\pa_T}{\PT{T}}{1}(\xi)} lies in \m{\exCone{T}{i_1}{0}} if and only if $\xi$ lies
in \m{\pieri{\pa_T}{\PT{T}}}: the set \m{\cC=\exCone{T}{i_1}{0}} satisfies the hypothesis of Theorem~2 for
\m{\pa=\pa_T} and \m{\varrho=\PT{T}}, whatever \m{i_1} is.

\begin{lemma}[a point deep inside the cone]\label{l:conepoint}
In Setting~\ref{set:T} with $\cO$ a field, let \m{i_0}, \m{i_1} be indices with \m{i_0+1=i_1}. There is a vector
\m{\sigma^{\circ}} in \m{\exCone{T}{i_1}{-1}} with \m{0\le\sigma^{\circ}(v)\le\PT{T}^{3i_1}} for every
variable~$v$.
\end{lemma}
\begin{proof}
Let
\[\mathrm{g}_i=2\,\PT{T}^{2(i_0-i)}\ \text{ for } i\le i_0\qquad\text{and}\qquad\mathrm{g}_i=0\ \text{ for } i>i_0,\]
so that \m{\mathrm{g}_i=\PT{T}^{2}\mathrm{g}_{i+1}+2\,[i+1=i_1]}. Let \m{\xi(i,s)=\mathrm{g}_i} for even~$s$ and
\m{\xi(i,s)=\PT{T}\mathrm{g}_i} for odd~$s$, and
\[\sigma^{\circ}=\permStep{\pa_T}{\PT{T}}{1}(\xi).\]
Then
\m{\sigma^{\circ}(i,s)=(\PT{T}^{2}-1)\mathrm{g}_i} for even~$s$ and $0$ for odd~$s$, which lies between $0$ and
\m{\PT{T}^{3i_1}}. By Lemma~\ref{l:exdot}, the excess of \m{\sigma^{\circ}} at~$s_0$ for the indices \m{i}, \m{i+1} is
\m{(1-1/\qT{T})} times
\begin{align*}
(\PT{T}^{2}-1)\mathrm{g}_{i+1}-(\mathrm{g}_i-\mathrm{g}_{i+1})&=-2\,[i+1=i_1]&&\text{if $s_0$ is even},\\
-\PT{T}(\mathrm{g}_i-\mathrm{g}_{i+1})&\le-2\,[i+1=i_1]&&\text{if $s_0$ is odd};
\end{align*}
and \m{2(1-1/\qT{T})\ge1}.
\end{proof}
\leanline{\lean{conepoint}}

\begin{lemma}[the volume of a slice of a simplex]\label{l:simplex}
Let \m{n'\ge0} be a natural number and \m{\delta\le\delta'} real numbers. The set of the $\eta$ in
\m{\RR^{n'+1}} with all coordinates \m{\ge0} and \m{\delta-1\le\sum_i\eta_i\le\delta} has volume at most
\m{|\delta'^{\,n'}/n'!|}.
\end{lemma}
\begin{proof}
If \m{\delta'<0} the set is empty. Otherwise argue by induction on~$n'$. For \m{n'=0} the set lies in an
interval of length~$1$. For the step, the slice at a value
\m{\eta_0} of the first coordinate is empty unless \m{0\le\eta_0\le\delta\le\delta'}, and is then a set of the same
kind in one dimension less, for \m{\delta-\eta_0\le\delta'-\eta_0}; and
\m{\int_0^{\delta'}(\delta'-\eta_0)^{n'-1}/(n'-1)!\,d\eta_0=\delta'^{\,n'}/n'!}.
\end{proof}
\leanline{\lean{simplex}}

\begin{proposition}[forms obeying the excess condition, against a volume]\label{p:dimvol}
In Setting~\ref{set:T}, let the ring of~$T$ be a field~$\FF$ of characteristic~$p$, let \m{i_0}, \m{i_1} be
indices with \m{i_0+1=i_1}, $n$ a natural number and $h$ a real number. Write \m{\Lev^{\circ}} for the Levi subgroup \m{\Lev_{\twoBlock{k}{i_1}{\cdot}}(W)}, and \m{\Par^{\circ}} for the
parahoric subgroup of the labelling \m{i\mapsto i}, and let \m{\delta_s=n+k+1} for every letter~$s$. Then
\[\dim_{\FF}\exc{T}{n}{i_0}{i_1}{0}{h}\ \le\ [\Lev^{\circ}:\Lev^{\circ}\cap\Par^{\circ}]\cdot
\vol\Bigl(\exCone{T}{i_1}{h+L(k+1)}\cap\slab{k}{\delta}\Bigr).\]
\end{proposition}
\begin{proof}
Write \m{\Xi=\exc{T}{n}{i_0}{i_1}{0}{h}} and $\cC$ for the set on the right; $\cC$ is bounded and convex. The
number of the points of~$\cC$ with coordinates in \m{\frac1Q\mathbb{Z}}, divided by \m{Q^{kL}}, tends to
\m{\vol(\cC)} as the natural number~$Q$ grows, a bounded convex set having a boundary of measure zero. So it
suffices to show, for \m{Q=\qT{T}^{l}} with \m{l\ge1}, that
\begin{equation}\label{e:count}
Q^{kL}\cdot\dim_{\FF}\Xi\ \le\ [\Lev^{\circ}:\Lev^{\circ}\cap\Par^{\circ}]\cdot\#\bigl(\cC\cap\tfrac1Q\mathbb{Z}^{k\times L}\bigr).
\end{equation}
The ring of~$T$ being a field of characteristic~$p$, we have \m{p=0} and \m{\varpi=0} in~$\FF$, and
\m{\iota_{s+j}(\upsilon)=\iota_s(\upsilon)^{\PT{T}^{j}}} by Lemma~\ref{l:basic}(a); in particular
\m{\iota_s(\upsilon)^{Q}=\iota_s(\upsilon)}. Hence \m{\polyAction{T}{\gamma}} commutes with the substitution
\m{x_v\mapsto x_v^{Q}}, and for a form~$F$ the polynomial \m{F(x^{Q})} is sent by \m{\polyAction{T}{\gamma}} to
\m{(\letterRep{T}{n}{\gamma}F)(x^{Q})} (Lemma~\ref{l:basic}(c)).

\emph{Amplification.} Let $\mathcal{U}$ be the set of the functions~$\nu$ from the variables to
\m{\{0,\dots,Q-1\}}. For $\nu$ in~$\mathcal{U}$ let \m{\hat\nu} be the exponent with
\[\hat\nu(i,s)=\nu(i,s)+Q\,[i=i_0]\,\Bigl(k-\Bigl\lfloor\sum_{i'}\nu(i',s)/Q\Bigr\rfloor\Bigr);\]
it is congruent to~$\nu$ modulo~$Q$, and its degree in each letter lies in \m{[kQ,kQ+Q)}. Send a tuple
\m{(F_{\nu})_{\nu\in\mathcal{U}}} of elements of~$\Xi$ to the polynomial
\m{\Phi=\sum_{\nu}x^{\hat\nu}F_{\nu}(x^{Q})}. The exponents of the $\nu$-th summand are congruent to~$\nu$
modulo~$Q$, so the coefficient of \m{x^{Q\beta+\hat\nu}} in~$\Phi$ is \m{F_{\nu}(\beta)}, and the map is
injective. Let $\Gamma$ be a set of representatives of the right cosets \m{(\Lev^{\circ}\cap\Par^{\circ})\gamma}
in~\m{\Lev^{\circ}}. We show that the linear map sending \m{(F_{\nu})_{\nu}} to the family of the coefficients of
\m{x^{Q\eta}} in \m{\polyAction{T}{\gamma}\Phi}, for $\gamma$ in~$\Gamma$ and $\eta$ in
\m{\cC\cap\frac1Q\mathbb{Z}^{k\times L}}, is injective; comparing dimensions, this is~\eqref{e:count}.

\emph{A least term.} Let $\omega$ be the weight with \m{\omega(i,s)=i} for even~$s$ and \m{\omega(i,s)=-i} for
odd~$s$. Suppose \m{\Phi\ne0}. Among the pairs of a $\gamma$ in~$\Gamma$ and an exponent~$\alpha$ with nonzero
coefficient in \m{\polyAction{T}{\gamma}\Phi}, choose one, \m{(\gamma_0,\alpha)}, with \m{\pr{\omega}{\alpha}}
least. We show that \m{\alpha/Q} lies in~$\cC$; then the coefficient of \m{x^{\alpha}} in
\m{\polyAction{T}{\gamma_0}\Phi} is one of the coordinates of the image and is not zero.

By the commutation above,
\[\polyAction{T}{\gamma_0}\Phi=\sum_{\nu}(\polyAction{T}{\gamma_0}x^{\hat\nu})\cdot(\letterRep{T}{n}{\gamma_0}F_{\nu})(x^{Q}),\]
so \m{\alpha=Q\beta+\alpha''} for some~$\nu$, a family~$\beta$ with
\m{(\letterRep{T}{n}{\gamma_0}F_{\nu})(\beta)\ne0}, and an exponent \m{\alpha''} of
\m{\polyAction{T}{\gamma_0}x^{\hat\nu}}. The action preserves the degree in each letter, so the degree
of~\m{\alpha''} in each letter lies in \m{[kQ,kQ+Q]}. Hence \m{\alpha/Q} has coordinates \m{\ge0} and its sums over
the indices lie in \m{[n+k,n+k+1]}: it lies in the slab.

\emph{The excess at \m{(i_0,i_1)}.} As \m{F_{\nu}} lies in~$\Xi$ and \m{\gamma_0} in~\m{\Lev^{\circ}}, the
family~$\beta$ has excess at most~$h$ at every letter. The excess of \m{\alpha''} is at most its degree, at most
\m{L(kQ+Q)}, by Lemma~\ref{l:exIcc}. So the excess of \m{\alpha/Q} is at most \m{h+L(k+1)}.

\emph{The excess at the other pairs.} Let \m{i}, \m{i+1} be indices with \m{i+1\ne i_1}, and $s_0$ a letter. We
apply Theorem~1 to the indices \m{i}, \m{i+1}, the labelling \m{\twoBlock{k}{i_1}{\cdot}}, which takes the same
value at both, the weight~$\omega$, for which the differences of the hypothesis are equal to~$1$, the polynomial
\m{\polyAction{T}{\gamma_0}\Phi}, the exponent~$\alpha$ and \m{e=0}. Its coefficient at \m{x^{\alpha}} is not
zero, that is, not divisible by \m{\varpi=0}. For $\gamma$ in~\m{\Lev^{\circ}} write
\m{\gamma^{-1}\gamma_0=\gamma'\gamma_1} with \m{\gamma'} in \m{\Lev^{\circ}\cap\Par^{\circ}} and \m{\gamma_1}
in~$\Gamma$. All exponents of \m{\polyAction{T}{\gamma_1}\Phi} have weight at least \m{\pr{\omega}{\alpha}}, by the
choice of the pair. By Lemma~\ref{l:par}, for the labelling \m{i\mapsto i}, and since \m{p=0} in~$\FF$, the element
\m{\gamma'} sends a variable \m{x_{i',s}} to a combination of the variables \m{x_{l,s}} with \m{l\ge i'} if $s$
is even and with \m{l\le i'} if $s$ is odd: of variables of weight at least \m{\omega(i',s)}. So all exponents of
\m{\polyAction{T}{\gamma^{-1}}\polyAction{T}{\gamma_0}\Phi} have weight at least \m{\pr{\omega}{\alpha}}, which is
the hypothesis of Theorem~1 for \m{\varpi=0}. The theorem gives that the excess of~$\alpha$ at~$s_0$ for the
indices \m{i}, \m{i+1} is at most~$0$.
\end{proof}
\leanline{\lean{dimvol}\sep \lean{rootlemma}\sep \lean{pdvdpar}\sep \lean{exwIcc}}

\begin{proof}[Proof of Theorem 3(a)]
There is an index, so \m{k\ge1}. Of the size conditions we use \m{\kappa\le1}, \m{1\le h},
\begin{equation}\label{e:sizeA}
\frac{2(L+2)k}{h}\le\kappa\qquad\text{and}\qquad\bigl(h+L(k+1)\bigr)\,k\,\PT{T}^{3i_1}+k+2\le\frac nh .
\end{equation}
Let \m{\sigma^{\circ}} be the vector of Lemma~\ref{l:conepoint} and \m{\sigma^{\bullet}=(h+L(k+1))\,\sigma^{\circ}}.
The conditions that define the cones are linear, so the translation by~\m{\sigma^{\bullet}} maps
\m{\exCone{T}{i_1}{h+L(k+1)}\cap\slab{k}{\delta}}, with \m{\delta_s=n+k+1}, into
\m{\exCone{T}{i_1}{0}\cap\slab{k}{\delta'}} with \m{\delta'_s=n+k+1+\sum_i\sigma^{\bullet}(i,s)}. Let \m{\Lev^{\circ}} and
\m{\Par^{\circ}} be as in Proposition~\ref{p:dimvol}. Both lie in the parahoric subgroup
\m{\Par_{\twoBlock{k}{i_1}{\cdot}}(W)}, so its index times \m{[\Lev^{\circ}:\Lev^{\circ}\cap\Par^{\circ}]} is at most
\m{[\GLk(W):\Par^{\circ}]}. Now \m{\Par^{\circ}} is the preimage of the group of the lower triangular matrices
under the reduction map from \m{\GLk(W)} to \m{\GLk(\FF_W)}. So its index is the index of the lower triangular
matrices of the image in the image, which is at most their index in \m{\GLk(\FF_W)}. (The two are equal, the
reduction map being surjective: a matrix over the local ring~$W$ is invertible when its reduction is. The formal
proof uses the inequality only.) By Lemma~\ref{l:borel} and \m{\#\FF_W=\qT{T}}, that index is at most the product
\[\prod_{l=1}^{k}\frac{\qT{T}^{l}-1}{\qT{T}-1}.\]
By Theorem~2, with the remark after Lemma~\ref{l:exdot}, this product times the volume of
\m{\exCone{T}{i_1}{0}\cap\slab{k}{\delta'}} is \m{\vol(\slab{k}{\delta'})}. With Proposition~\ref{p:dimvol} we get
\[\bigl[\GLk(W):\Par_{\twoBlock{k}{i_1}{\cdot}}(W)\bigr]\cdot\dim_{\FF}\exc{T}{n}{i_0}{i_1}{0}{h}\ \le\
\vol\bigl(\slab{k}{\delta'}\bigr).\]
The slab is contained in the product over the letters of the sets of Lemma~\ref{l:simplex} in \m{\RR^{k}}. By
the second condition of~\eqref{e:sizeA} and \m{\sigma^{\circ}\le\PT{T}^{3i_1}},
\m{\delta'_s\le n+n/h\le n\exp(1/h)}. So the volume is at most
\m{\bigl((n\exp(1/h))^{k-1}/(k-1)!\bigr)^{L}}, and
\[\frac{n^{k-1}}{(k-1)!}\ \le\ \binom{n+k-1}{k-1}=\#\Mon_{k,n}.\]
Finally
\m{\exp((k-1)L/h)\le1+\kappa}: the exponent is at most \m{\kappa/2\le\tfrac12} by the first condition
of~\eqref{e:sizeA}, and \m{\exp(\theta)\le1+2\theta} for \m{0\le\theta\le\tfrac12}.
\end{proof}
\leanline{\lean{dimex}}

\begin{proposition}[a nonzero deep form has a costly monomial]\label{p:deepweight}
In Setting~\ref{set:T}, let the ring of~$T$ be a field~$\FF$ of characteristic~$p$, let $n$, $\thr$, $j$ be natural
numbers with \m{2\thr\le k} and \m{2\thr<L}, and let $F$ be a nonzero form in \m{\deepForms{T}{k}{n}{\thr}{j}}. Then
there is a family~$\beta$ with \m{F(\beta)\ne0} and
\[j\ \le\ \sum_{v=(i,s)}\beta_s(i)\,\frac{\dpi{T}{\thr}{v}}{\PT{T}^{\,e_i}},\qquad e_i=2\thr\ \text{ if } i\ge2\thr,\quad
e_i=1\ \text{ if } i<2\thr .\]
\end{proposition}
\begin{proof}
Let \m{\FF'} be the field of rational functions over~$\FF$ in the variables that are not deep, and regard
the polynomial of~$F$ as a polynomial~$\Phi$ over~\m{\FF'} in the deep variables; it is not zero. We apply
Input~3 over~\m{\FF'}, with $\cV$ the set of the deep variables and \m{\varkappa(i,s)=\PT{T}^{e_i}}. Then
\m{\prod_v\varkappa(v)=\PT{T}^{(L/2)\sum_ie_i}=\qT{T}^{\thr(k-2\thr+1)}} by Lemma~\ref{l:sumdeep}. A monomial of~$\Phi$
of cost at least~$j$ comes from a family~$\beta$ with \m{F(\beta)\ne0} whose deep part it is, and the cost is
the sum of the statement. It remains to give the points and the polynomials.

\emph{The points.} Let $\cK$ be the set of the pairs \m{(l,l')} of indices with \m{l<\thr} and either
\m{l'=l+\thr} or \m{l'\ge2\thr}; it has \m{\thr(k-2\thr+1)} elements. Let $\cS$ be the set of the functions~$\fs$
from~$\cK$ to~\m{\FF_W}, so that \m{\#\cS=\qT{T}^{\thr(k-2\thr+1)}}. For $\fs$ in~$\cS$ choose a lift
\m{\tilde{\fs}_{l,l'}} in~$W$ of each value and let \m{\gamma_{\fs}=1+\sum_{(l,l')\in\cK}\tilde{\fs}_{l,l'}E_{ll'}}, an
element of \m{\GLk(W)} with inverse \m{1-\sum\tilde{\fs}_{l,l'}E_{ll'}}, the rows and the columns that occur being
disjoint. By Lemma~\ref{l:basic}(b), \m{\polyAction{T}{\gamma_{\fs}^{-1}}} fixes the variables that are not deep and
sends a deep variable \m{x_v} to \m{x_v-\upsilon_{\fs}(v)}, where
\[\upsilon_{\fs}(l',s)=\sum_{l}\iota_s(\tilde{\fs}_{l,l'})\,x_{l,s}\ \ (s\ \text{even}),\qquad
\upsilon_{\fs}(l,s)=-\sum_{l'}\iota_s(\tilde{\fs}_{l,l'})\,x_{l',s}\ \ (s\ \text{odd}),\]
the sums being over the pairs in~$\cK$ with the given member; these are elements of~\m{\FF'}, the variables on
the right not being deep. As $F$ lies in \m{\deepForms{T}{k}{n}{\thr}{j}}, every monomial of
\m{\letterRep{T}{n}{\gamma_{\fs}}F} has at least $j$ deep variables, counted with multiplicity; applying
\m{\polyAction{T}{\gamma_{\fs}^{-1}}} and Lemma~\ref{l:basic}(c), $\Phi$ is a combination of products of at least
$j$ factors \m{x_v-\upsilon_{\fs}(v)}. So $\Phi$ lies in the $j$-th power of the ideal of the
point~\m{\upsilon_{\fs}}.

\emph{The polynomials.} Fix a letter~$s'$. For a pair \m{(l,l')} in~$\cK$ choose a natural number
\m{j_{l,l'}}: it is $0$ if \m{l'=l+\thr}; if \m{l'\ge2\thr} it is the one of \m{2l}, \m{2l+1} for which the letter
\m{s'-j_{l,l'}} is even. Let \m{v_{l,l'}} be the deep variable \m{(l',s'-j_{l,l'})} if that letter is even, and
\m{(l,s')} if \m{j_{l,l'}=0} and $s'$ is odd. In each case \m{j_{l,l'}<e}, where
\m{\PT{T}^{e}=\varkappa(v_{l,l'})}: for \m{l'\ge2\thr} we have \m{j_{l,l'}\le2\thr-1} and \m{e=2\thr}. By
Lemma~\ref{l:basic}(a) and \m{p=0},
\[\upsilon_{\fs}(v_{l,l'})^{\PT{T}^{j_{l,l'}}}=\sum_{(l_1,l'_1)\in\cK}C_{(l,l'),(l_1,l'_1)}\cdot\iota_{s'}(\tilde{\fs}_{l_1,l'_1}),\]
with a matrix~$C$ over~\m{\FF'} that does not depend on~$\fs$: its entries are, up to sign, the
\m{\PT{T}^{j_{l,l'}}}-th powers of the variables that occur in the formula for \m{\upsilon_{\fs}(v_{l,l'})}. We
claim that \m{\det C\ne0}. If $s'$ is even, $C$ is block diagonal for the second members: the block of the pairs
\m{(l,l+\thr)} is diagonal with the entries \m{x_{l,s'}}, and the block of each \m{l'\ge2\thr} is the matrix
\m{(x_{l_1,s'-2l}^{\,\PT{T}^{2l}})_{l,l_1}}, whose entries are powers of \m{\thr^{2}} distinct variables, the
letters \m{s'-2l} being distinct because \m{2\thr<L}. If $s'$ is odd, order the pairs \m{(l,l+\thr)} first: $C$ is
block triangular, with a diagonal block of entries \m{-x_{l+\thr,s'}} and blocks
\m{(x_{l_1,s'-2l-1}^{\,\PT{T}^{2l+1}})_{l,l_1}}. A matrix whose entries are powers of distinct variables has a
nonzero determinant. So each \m{\iota_{s'}(\tilde{\fs}_{l_1,l'_1})} is a combination over~\m{\FF'}, with
coefficients independent of~$\fs$, of the values at \m{\upsilon_{\fs}} of the monomials
\m{x_{v_{l,l'}}^{\PT{T}^{j_{l,l'}}}}, which cost \m{\PT{T}^{j_{l,l'}}/\varkappa(v_{l,l'})<1}.

Now let $v$ be a deep variable of the letter~$s$, and \m{\PT{T}^{e}=\varkappa(v)}. By Lemma~\ref{l:basic}(a),
\m{\upsilon_{\fs}(v)^{\varkappa(v)}} is a combination, with coefficients in~\m{\FF'} independent of~$\fs$, of the
\m{\iota_{s+e}(\tilde{\fs}_{l_1,l'_1})}; by the last paragraph for \m{s'=s+e} it is the value at \m{\upsilon_{\fs}} of
a polynomial \m{\Phi_v} whose monomials cost less than~$1$. The same paragraph, for any letter~$s'$, shows that
the point \m{\upsilon_{\fs}} determines the \m{\iota_{s'}(\tilde{\fs}_{l_1,l'_1})}, hence~$\fs$: two elements
of~$W$ with the same image under~\m{\iota_{s'}} differ by a non-unit. So the points are pairwise distinct, and
Input~3 applies.
\end{proof}
\leanline{\lean{deepweight}\sep \lean{nullst}\sep \lean{sumdeep}}

\begin{proof}[Proof of Theorem 3(b)]
From the size conditions: \m{0<\thr}, \m{2\thr\le k}, \m{1\le y}, \m{n<\PT{T}^{\lambda}} and \m{L=2\lambda}, and
\m{\PT{T}^{3\thr}\le n} from the ninth condition; so \m{3\thr<\lambda} and \m{2\thr<L}. For a family~$\beta$ let
\m{\mathrm{cost}(\beta)} be the sum of Proposition~\ref{p:deepweight}. By that proposition, a form in
\m{\deepForms{T}{k}{n}{\thr}{j}} that vanishes at all~$\beta$ with \m{j\le\mathrm{cost}(\beta)} is zero; so
\m{\dim_{\FF}\deepForms{T}{k}{n}{\thr}{j}} is at most the number of these~$\beta$. Summing over \m{j<j_0}, each
$\beta$ is counted at most \m{\mathrm{cost}(\beta)+1} times:
\[\sum_{j<j_0}\dim_{\FF}\deepForms{T}{k}{n}{\thr}{j}\ \le\ \sum_{\beta}\bigl(1+\mathrm{cost}(\beta)\bigr).\]
The cyclic shifts of the indices permute the families. Averaging over them replaces \m{\beta_s(i)} by \m{n/k},
so, by Lemma~\ref{l:sumdeep},
\[\sum_{\beta}\mathrm{cost}(\beta)=\bigl(\#\Mon_{k,n}\bigr)^{L}\cdot\frac nk\cdot\frac L2\cdot
\Bigl(\frac{k-2\thr}{\PT{T}^{2\thr}}+\frac{2\thr}{\PT{T}}\Bigr)\ \le\
\bigl(\#\Mon_{k,n}\bigr)^{L}\cdot L\,n\,\Bigl(\frac1{\PT{T}^{2\thr}}+\frac{\thr^{2}}{k\,\PT{T}}\Bigr),\]
which is at most \m{z\,(\#\Mon_{k,n})^{L}} by the tenth condition.
\end{proof}
\leanline{\lean{dimdeep}}

\sect{Conjugation-invariant means}{s:mean}
In this section the variables are the entries \m{x_{i,i'}} of one generic matrix, and for a finite set~$\cV$ of
variables \m{\Mon_{\cV}(n)} is the set of the exponents of degree~$n$ on~$\cV$.

\begin{lemma}[density of the points where a form is a unit]\label{l:zariski}
Let $W$ be a local ring, $\cO$ a ring, \m{\chi\colon W\to\cO} a ring homomorphism, $\cV$ a finite set, and $n$,
$n'$ natural numbers. Let \m{\Phi_0} be a homogeneous polynomial of degree~$n'$ in \m{W[x_{\cV}]}, with
\m{n+n'\le\#\FF_W}, and let \m{\upsilon_0} be a point of \m{W^{\cV}} at which \m{\Phi_0} takes a unit value. Let
$M_1$ be an $\cO$-submodule of \m{\cO^{\Mon_{\cV}(n)}} that contains the vector
\m{(\chi(\upsilon^{\alpha}))_{\alpha}} for every point $\upsilon$ of \m{W^{\cV}} with \m{\Phi_0(\upsilon)} a unit.
Then \m{M_1=\cO^{\Mon_{\cV}(n)}}.
\end{lemma}
\begin{proof}
Let $M_0$ be the $W$-submodule of \m{W^{\Mon_{\cV}(n)}} of the vectors whose image under~$\chi$ lies in~$M_1$. It
suffices that \m{M_0} is everything, for then $M_1$ contains the unit vectors. By Nakayama's lemma, the finitely
generated $W$-module \m{W^{\Mon_{\cV}(n)}/M_0} is zero if every $W$-linear map \m{\mathcal{L}} from it to~\m{\FF_W}
is zero. Given such a map, let $\Phi$ be the polynomial over~\m{\FF_W} with the coefficient \m{\mathcal{L}} of the
$\alpha$-th unit vector at~\m{x^{\alpha}}; it is homogeneous of degree~$n$. If $\upsilon$ is a point at which
\m{\Phi_0} is a unit, then \m{(\upsilon^{\alpha})_{\alpha}} lies in~$M_0$, so $\Phi$ vanishes at the residue
of~$\upsilon$. Hence the product of $\Phi$ with the reduction of~\m{\Phi_0} vanishes at every point of
\m{\FF_W^{\cV}}. It is homogeneous of degree \m{n+n'\le\#\FF_W}, and such a polynomial is zero, by the zero test
that follows; the reduction of \m{\Phi_0} is not zero, because of~\m{\upsilon_0}; so \m{\Phi=0} and
\m{\mathcal{L}=0}.

\emph{The zero test.} Let $\FF$ be a field and \m{\Phi'} a polynomial in \m{\FF[x_{\cV}]} that vanishes at every
point of \m{\FF^{\cV}}. First, if $\FF$ is finite and the degree of \m{\Phi'} in each variable is less
than~\m{\#\FF}, then \m{\Phi'=0}. This is proved by induction on \m{\#\cV}: write \m{\Phi'} as a polynomial in one
of the variables, with coefficients in the others; at each point of the other variables it has more zeros than
its degree, so the coefficients vanish at every point, and they are zero by the induction hypothesis. Secondly,
if $\FF$ is finite and \m{\Phi'} is homogeneous of a degree \m{l\le\#\FF}, then \m{\Phi'=0}. For \m{l<\#\FF} this is
the first statement. Let \m{l=\#\FF}. The monomials of degree~$l$ in which some variable has an exponent that
is not less than~\m{\#\FF} are the powers \m{x_v^{\#\FF}}. Replace each of them by~\m{x_v}. The values on
\m{\FF^{\cV}} do not change, as \m{\upsilon^{\#\FF}=\upsilon} for $\upsilon$ in~$\FF$. The monomials stay distinct and keep
their coefficients, as \m{\#\FF\ge2}: the new ones have degree~$1$, the others degree~\m{\#\FF}. The new
polynomial has degree less than~\m{\#\FF} in each variable, so it is zero by the first statement, and then every
coefficient of \m{\Phi'} is zero. In the proof above $\FF$ is the residue field~\m{\FF_W}. If it is infinite, the
count \m{\#\FF_W} is~$0$ by the convention, so \m{n+n'=0} and the product is a constant; it vanishes at a point,
so it is zero. The case \m{l=\#\FF} belongs to the statement of the lemma, as in the formal text; there the
polynomial \m{x_v^{\#\FF}-x_v} shows that homogeneity cannot be dropped.
\end{proof}
\leanline{\lean{zariski}}

\begin{lemma}[an integral invariant diagonal form]\label{l:diagform}
Let $\cO$ be a valuation domain of characteristic zero, $\cV$ a finite set, $p$ a prime number that is not a unit
in~$\cO$, and $e$ a natural number. There is a function~$w$ from the exponents on~$\cV$ to~$\cO$ such that
\begin{enumerate}[label=\textup{(\roman*)},leftmargin=2.3em,itemsep=1pt,topsep=2pt]
\item for every exponent~$\alpha$ with \m{|\alpha|<p^{e}}: \m{p^{\#\cV\cdot e}=\binom{|\alpha|}{\alpha}\,w(\alpha)},
where the multinomial coefficient \m{\binom{|\alpha|}{\alpha}} is \m{|\alpha|!/\prod_v\alpha(v)!};
\item for every square matrix~$G$ over~$\cO$ indexed by~$\cV$ and all exponents $\alpha$, $\alpha'$: the
coefficient of \m{x^{\alpha}} in the linear substitution of~$G$ applied to~\m{x^{\alpha'}}, times \m{w(\alpha)},
equals \m{w(\alpha')} times the coefficient of \m{x^{\alpha'}} in the linear substitution of~\m{G^{\tp}} applied
to~\m{x^{\alpha}}.
\end{enumerate}
\end{lemma}
\begin{proof}
A multinomial coefficient of degree~$l$ is a product of \m{\#\cV} binomial coefficients \m{\binom{l'}{j}} with
\m{j\le l'\le l}. Let \m{l'<p^{e}}. By Legendre's formula for the power of~$p$ in a factorial, the exponent of~$p$
in \m{\binom{l'}{j}} is
\[\sum_{r\ge1}\Bigl(\Bigl\lfloor\frac{l'}{p^{r}}\Bigr\rfloor-\Bigl\lfloor\frac{j}{p^{r}}\Bigr\rfloor
-\Bigl\lfloor\frac{l'-j}{p^{r}}\Bigr\rfloor\Bigr).\]
Each term is $0$ or~$1$, and it is $0$ when \m{p^{r}>l'}, hence for \m{r\ge e}; so the exponent is at most~$e$.
The part of \m{\binom{l'}{j}} prime to~$p$ is a unit of~$\cO$. Indeed, an integer prime to~$p$ has a multiple that
is congruent to~$1$ modulo~$p$; as $p$ is not a unit, it lies in the maximal ideal~$\fm$ of the valuation
domain~$\cO$; so that multiple, and with it the integer, does not lie in~$\fm$, and the integer is a unit. Hence
\m{\binom{l'}{j}} divides \m{p^{e}} in~$\cO$, and \m{\binom{|\alpha|}{\alpha}} divides \m{p^{\#\cV\cdot e}} when
\m{|\alpha|<p^{e}}; let \m{w(\alpha)} be the quotient, and \m{w(\alpha)=0} for the other~$\alpha$. For~(ii),
expand the product \m{\prod_{v'}(\sum_vG_{vv'}x_v)^{\alpha'(v')}}: the coefficient of \m{x^{\alpha}} is the sum,
over the matrices \m{(m_{vv'})} of natural numbers with row sums $\alpha$ and column sums~$\alpha'$, of
\m{\prod_{v'}\alpha'(v')!\cdot\prod_{v,v'}G_{vv'}^{m_{vv'}}/m_{vv'}!}. Multiplied by \m{\prod_v\alpha(v)!} this is
symmetric under transposition. So both coefficients of~(ii) vanish unless \m{|\alpha|=|\alpha'|}, and then they
differ by the factor \m{\prod\alpha'(v)!/\prod\alpha(v)!=\binom{|\alpha|}{\alpha}/\binom{|\alpha'|}{\alpha'}}, which
is \m{w(\alpha')/w(\alpha)} if the degree is less than~\m{p^{e}}; otherwise both values of~$w$ are zero.
\end{proof}
\leanline{\lean{diagform}}

\begin{proposition}[an invariant mean with prescribed values]\label{p:mean}
Let $W$ be a local ring, $\cO$ a valuation domain of characteristic zero, \m{\chi\colon W\to\cO} a ring
homomorphism, $k$ and $n_0$ natural numbers, and $\tau$ a labelling of the indices. Let $p$ be a prime number
that is not a unit in~$\cO$ and $e$ a natural number with \m{n_0<p^{e}} and \m{n_0+k<\#\FF_W}. Let $\omega$ be a
real function on the indices and $\Omega$ a set of real numbers. For an exponent~$\alpha$ on the generic
variables put \m{\mathrm{rw}(\alpha)=\sum_{i,i'}\alpha(i,i')\,\omega(i)} and
\m{\mathrm{cw}(\alpha)=\sum_{i,i'}\alpha(i,i')\,\omega(i')}. Then there is an $\cO$-linear form~$\cM$ on the
polynomials in the generic variables over~$\cO$ such that:
\begin{enumerate}[label=\textup{(\alph*)},leftmargin=2.3em,itemsep=1pt,topsep=2pt]
\item \m{\cM(\conjSubst{\chi}{\gamma}\Phi)=\cM(\Phi)} for every $\gamma$ in \m{\Lev_{\tau}(W)} and every
polynomial~$\Phi$;
\item \m{\cM(x^{\alpha})=0} for every $\alpha$ with \m{\mathrm{rw}(\alpha)\ne\mathrm{cw}(\alpha)};
\item for every $\alpha$ with \m{\mathrm{rw}(\alpha)\in\Omega}: \m{\cM(x^{\alpha})} is \m{p^{k^{2}e}} times the value
of \m{x^{\alpha}} at the identity matrix if \m{|\alpha|\le n_0}, and \m{\cM(x^{\alpha})=0} if \m{|\alpha|>n_0};
\item $\cM$ is an $\cO$-linear combination of forms \m{\cL\circ\conjSubst{\chi}{\gamma}}, with $\gamma$ in
\m{\Lev_{\tau}(W)} and $\cL$ an $\cO$-linear form that vanishes at every \m{x^{\alpha}} with
\m{\mathrm{rw}(\alpha)\notin\Omega}.
\end{enumerate}
\end{proposition}
\begin{proof}
We construct, for each \m{j\le n_0}, a form \m{\cM_j} that vanishes on the monomials of degree other than~$j$
and has the four properties for the monomials of degree~$j$; their sum is~$\cM$. Fix~$j$.

\emph{The algebra.} Let \m{V_j} be the free $\cO$-module with the basis \m{e_{\mathbf{l}}}, for the sequences
\m{\mathbf{l}=(l_1,\dots,l_j)} of indices, and put \m{|\mathbf{l}|_{\omega}=\sum_r\omega(l_r)}. Let \m{X^{\tau}} be
the matrix with the entry \m{x_{i',i}} at \m{(i,i')} if \m{\tau(i)=\tau(i')} and $0$ otherwise. An entry of its
$j$-th Kronecker power is \m{\prod_rX^{\tau}_{l_r,l'_r}}, which is zero or the monomial \m{x^{\alpha}} with
\m{\alpha=\sum_r1_{(l'_r,l_r)}}; then \m{\mathrm{rw}(\alpha)=|\mathbf{l}'|_{\omega}} and
\m{\mathrm{cw}(\alpha)=|\mathbf{l}|_{\omega}}. For $\alpha$ of degree~$j$ let \m{S_{\alpha}} be the endomorphism
of~\m{V_j} whose matrix has the entry $1$ where the Kronecker power has the entry \m{x^{\alpha}}, and $0$
elsewhere; for a vector \m{c=(c_{\alpha})} put \m{S(c)=\sum_{\alpha}c_{\alpha}S_{\alpha}}. Then for every linear
form~$\cL$, \m{S((\cL(x^{\alpha}))_{\alpha})} is $\cL$ applied to the entries of the Kronecker power. In
particular, for $\gamma$ in \m{\Lev_{\tau}(W)}, the point with \m{x_{i',i}=\gamma_{ii'}} gives a vector
\m{c_{\gamma}} with \m{S(c_{\gamma})=\chi(\gamma)^{\otimes j}}, the $j$-th Kronecker power of the matrix
\m{\chi(\gamma)}.

\emph{Density.} We claim: a submodule of the module of the vectors~$c$ that contains every $c$ with
\m{S(c)=\chi(\gamma)^{\otimes j}} for some $\gamma$ in \m{\Lev_{\tau}(W)} is everything. This is
Lemma~\ref{l:zariski} for the generic variables, the degree~$j$ and the polynomial \m{\Phi_0=\det X^{\tau}}, which
is homogeneous of degree~$k$ and equal to~$1$ at the point of the identity matrix; here
\m{j+k\le n_0+k<\#\FF_W}. Indeed, at a point~$\upsilon$ where \m{\det X^{\tau}} is a unit, the value of \m{X^{\tau}} is an
element~$\gamma$ of the Levi subgroup, and \m{S((\chi(\upsilon^{\alpha}))_{\alpha})=\chi(\gamma)^{\otimes j}}.

\emph{The module.} Let \m{V_j^{\Omega}} be the submodule of~\m{V_j} spanned by the vectors \m{S(c)e_{\mathbf{l}}}, for
all~$c$ and all \m{\mathbf{l}} with \m{|\mathbf{l}|_{\omega}\in\Omega}. It is mapped into itself by every
\m{S(c)}: by the density claim, used twice, it suffices to see that
\m{\chi(\gamma)^{\otimes j}\chi(\gamma')^{\otimes j}e_{\mathbf{l}}=\chi(\gamma\gamma')^{\otimes j}e_{\mathbf{l}}} lies
in~\m{V_j^{\Omega}}. Being finitely generated and without torsion over the valuation domain~$\cO$, \m{V_j^{\Omega}} is free.

\emph{The form.} Let $w$ be the function of Lemma~\ref{l:diagform} for the \m{k^{2}} generic variables, the
prime~$p$ and the exponent~$e$. For a polynomial~$\Phi$ with coefficients \m{\Phi_{\alpha}} let
\m{\mathrm{K}(\Phi)=(w(\alpha)\Phi_{\alpha})_{|\alpha|=j}}, and define
\[\cM_j(\Phi)=\operatorname{tr}\bigl(S(\mathrm{K}(\Phi))\,\big|\,V_j^{\Omega}\bigr),\qquad\text{so that}\quad
\cM_j(x^{\alpha})=w(\alpha)\operatorname{tr}(S_{\alpha}|V_j^{\Omega})\ \text{ for } |\alpha|=j .\]

\emph{Equivariance.} For $\gamma$ in \m{\Lev_{\tau}(W)},
\begin{equation}\label{e:equiv}
S\bigl(\mathrm{K}(\conjSubst{\chi}{\gamma}\Phi)\bigr)=\chi(\gamma)^{\otimes j}\,S(\mathrm{K}(\Phi))\,\chi(\gamma^{-1})^{\otimes j}.
\end{equation}
Indeed \m{\conjSubst{\chi}{\gamma}} is the linear substitution of a matrix~$G$ indexed by the generic variables.
By Lemma~\ref{l:diagform}(ii) the $\alpha$-th coordinate of \m{\mathrm{K}(\conjSubst{\chi}{\gamma}\Phi)} is
\m{\sum_{\alpha'}w(\alpha')\Phi_{\alpha'}} times the coefficient of \m{x^{\alpha'}} in the substitution of~\m{G^{\tp}}
applied to~\m{x^{\alpha}}. So the left side of~\eqref{e:equiv} is \m{\sum_{\alpha'}w(\alpha')\Phi_{\alpha'}} times the
matrix of the coefficients of \m{x^{\alpha'}} in the Kronecker power of the image of~\m{X^{\tau}} under the
substitution of~\m{G^{\tp}}. That substitution replaces the matrix \m{(x_{i,i'})} by
\m{\chi(\gamma^{-1})^{\tp}(x_{i,i'})\chi(\gamma)^{\tp}}, and therefore \m{X^{\tau}}, which is its transpose cut down to
the pattern of the Levi subgroup, by \m{\chi(\gamma)X^{\tau}\chi(\gamma)^{-1}}, the matrices \m{\chi(\gamma)} and
\m{\chi(\gamma)^{-1}} having that pattern. Taking coefficients gives the right side.

\emph{The four properties.} (a) By~\eqref{e:equiv} and because \m{\chi(\gamma)^{\otimes j}} and its inverse
map~\m{V_j^{\Omega}} to itself, the two traces agree. (b) \m{S_{\alpha}} sends \m{e_{\mathbf{l}'}} to zero unless
\m{|\mathbf{l}'|_{\omega}=\mathrm{rw}(\alpha)}, and then to a combination of the \m{e_{\mathbf{l}}} with
\m{|\mathbf{l}|_{\omega}=\mathrm{cw}(\alpha)}. If the two weights differ, \m{S_{\alpha}^{2}=0}, and a nilpotent
endomorphism of a free module over a domain has trace zero. (c) If \m{\mathrm{rw}(\alpha)\in\Omega}, then
\m{S_{\alpha}} maps \m{V_j} into~\m{V_j^{\Omega}}, so its trace on~\m{V_j^{\Omega}} is its trace on~\m{V_j}, the number of the
sequences \m{\mathbf{l}} with \m{\sum_r1_{(l_r,l_r)}=\alpha}. That number is \m{\binom{|\alpha|}{\alpha}} if
$\alpha$ is supported on the variables \m{x_{i,i}} and $0$ otherwise; and the value of \m{x^{\alpha}} at the
identity matrix is $1$ or $0$ in the same cases. With Lemma~\ref{l:diagform}(i), \m{j\le n_0<p^{e}}, this
is~(c). (d) Choose a basis of~\m{V_j^{\Omega}}; the trace is the sum, over the basis, of forms
\m{\Phi\mapsto\cL'(S(\mathrm{K}(\Phi))\,\mathbf{v})} with \m{\mathbf{v}} in~\m{V_j^{\Omega}} and \m{\cL'} a linear form
on~\m{V_j^{\Omega}}. By the definition of~\m{V_j^{\Omega}} and the density claim we may take
\m{\mathbf{v}=\chi(\gamma)^{\otimes j}e_{\mathbf{l}}} with
\m{|\mathbf{l}|_{\omega}\in\Omega}. By~\eqref{e:equiv} for \m{\gamma^{-1}},
\m{S(\mathrm{K}(\Phi))\chi(\gamma)^{\otimes j}e_{\mathbf{l}}=\chi(\gamma)^{\otimes j}S(\mathrm{K}(\conjSubst{\chi}{\gamma^{-1}}\Phi))e_{\mathbf{l}}},
so the form is \m{\cL\circ\conjSubst{\chi}{\gamma^{-1}}} with
\m{\cL(\Phi)=\cL'(\chi(\gamma)^{\otimes j}S(\mathrm{K}(\Phi))e_{\mathbf{l}})}, and \m{\cL(x^{\alpha})=0} unless
\m{\mathrm{rw}(\alpha)=|\mathbf{l}|_{\omega}}, which lies in~$\Omega$.
\end{proof}
\leanline{\lean{mean}\sep \lean{zariski}\sep \lean{diagform}}

The image of~$S$ is the Schur algebra of the Levi subgroup in degree~$j$, an order spanned by the Kronecker
powers of the group elements; \m{V_j^{\Omega}} is the part of the tensor space generated by the weight vectors with
weight in~$\Omega$, and \m{\cM_j} is its character, read as a function on the dual of the Schur algebra. The
function~$w$ is the price of integrality: the identification of polynomials of degree~$j$ with that dual
has the multinomial coefficients as denominators.

\sect{The Frobenius operator and the projector}{s:proj}
\begin{lemma}[the Frobenius operator under a ring homomorphism]\label{l:frobmap}
In Setting~\ref{set:T}, let $\thr$, $\lambda$, $n$ be natural numbers, $\cR$ a ring, \m{\psi=(\psi(s,j))} matrices in
\m{\Mk(\cR)}, \m{\chi\colon\cR\to\cO} a ring homomorphism, and $\gamma$ an element of
\m{\Lev_{\twoBlock{k}{\thr}{\cdot}}(W)}. Then $\chi$, applied to the entries of \m{\frob{T}{\thr}{\lambda}{n}{\psi}},
gives
\[\letterRep{T}{n}{\gamma^{-1}}\cdot\frob{T}{\thr}{\lambda}{n}{\psi'}\cdot\letterRep{T}{n}{\gamma},\qquad
\psi'(s,j)=\letterConj{T}{\gamma}{s}\bigl(\chi(\psi(s,j))\bigr).\]
\end{lemma}
\begin{proof}
Let \m{\Pi^{m}_s} and \m{\Pi'^{m}_s} be the products of Definition~\ref{d:frob}(b) for $\psi$ and~$\psi'$; then
\[\chi(\Pi^{m}_s)=\iota_s(\gamma^{-1})\,\Pi'^{m}_s\,\iota_s(\gamma).\]
For an even letter this says that
\m{\chi(\Pi^{m}_s)} is \m{\Pi'^{m}_s} conjugated by the letter matrices of~$\gamma$. For an odd letter, taking
transposed adjugates gives
\[\chi(\operatorname{adj}(\Pi^{m}_s))^{\tp}=\iota_s(\gamma)^{\tp}\,\operatorname{adj}(\Pi'^{m}_s)^{\tp}\,\iota_s(\gamma^{-1})^{\tp},\]
the determinants of \m{\iota_s(\gamma)} and \m{\iota_s(\gamma^{-1})} cancelling; and \m{\iota_s(\gamma)^{\tp}} is the
letter matrix of~\m{\gamma^{-1}} at an odd letter. The number \m{m=\lambda+\dpi{T}{\thr}{(i,s)}} depends only on the
block of~$i$, and the letter matrices of~$\gamma$ are block diagonal; so choosing the rows according to~$m$
commutes with these conjugations, and
\[\chi(B_s)=\letterMatrix{T}{\gamma^{-1}}{s}\,B'_s\,\letterMatrix{T}{\gamma}{s}.\]
Symmetric powers and Kronecker products are multiplicative.
\end{proof}
\leanline{\lean{frobmap}}

\begin{lemma}[the Frobenius operator in a filtration]\label{l:frobmem}
In Setting~\ref{set:T}, let $\thr$, $\lambda$, $n$ be natural numbers, $\cR$ a ring, \m{\psi=(\psi(s,j))} matrices
in \m{\Mk(\cR)}, and \m{\beta,\beta'} families in \m{\Fam}. Let \m{\cG} be a family of additive subgroups of~$\cR$
such that \IsRF{\cG}. Let \m{\eta_{s,j}(i)} be real numbers, for a letter, a natural number and an index, let
$\omega$ be a weight on the variables and $\epsilon$ a real number, such that
\begin{enumerate}[label=\textup{(\roman*)},leftmargin=2.3em,itemsep=1pt,topsep=2pt]
\item \m{|\eta_{s,0}(i)|\le\epsilon} for all $s$, $i$, and \m{\sum_i\eta_{s,j}(i)=0} for all $s$, $j$;
\item \m{\omega(i,s)=\eta_{s,m}(i)} if $s$ is even and \m{\omega(i,s)=-\eta_{s,m}(i)} if $s$ is odd, where
\m{m=\lambda+\dpi{T}{\thr}{(i,s)}};
\item \m{\psi(s,j)_{ii'}\in\cG\bigl(\eta_{s,j+1}(i)-\eta_{s,j}(i')-\epsilon\bigr)} for every letter~$s$, every
\m{j\le\lambda} and all indices $i$, $i'$;
\item \m{0\le\epsilon}.
\end{enumerate}
Then the \m{(\beta,\beta')} entry of \m{\frob{T}{\thr}{\lambda}{n}{\psi}} lies in
\m{\cG\bigl(\pr{\omega}{\beta}-L\,n\,k\,(\lambda+2)\,\epsilon\bigr)}.
\end{lemma}
\begin{proof}
First, \m{(\Pi^{m}_s)_{ii'}} lies in \m{\cG(\eta_{s,m}(i)-(m+1)\epsilon)} for \m{m\le\lambda+1}: for \m{m=0} the
entries are $0$ and~$1$ and \m{\eta_{s,0}(i)-\epsilon\le0}; and
\m{(\psi(s,m)\Pi^{m}_s)_{ii'}=\sum_l\psi(s,m)_{il}(\Pi^{m}_s)_{li'}}, where the levels add to
\m{\eta_{s,m+1}(i)-(m+2)\epsilon}. Next, every entry \m{(B_s)_{ii'}} lies in
\m{\cG(\omega(i,s)-k(\lambda+2)\epsilon)}. For even~$s$ this is the first step. For odd~$s$ the entry is, up to
sign, a minor of \m{\Pi^{m}_s} without the row~$i$, a sum of products of one entry from each other row; the
levels add to \m{\sum_{l\ne i}\eta_{s,m}(l)-(k-1)(m+1)\epsilon=-\eta_{s,m}(i)-(k-1)(m+1)\epsilon}, by~(i).
Finally an entry of \m{\mathrm{Sym}^{n}(B_s)} in the row~\m{\beta_s} is a sum of products of $n$ entries
of~$B_s$, of which \m{\beta_s(i)} lie in the row~$i$; and the entry of the Kronecker product is the product over
the $L$ letters.
\end{proof}
\leanline{\lean{frobmem}}

\begin{lemma}[conjugation and Frobenius modulo $p$]\label{l:cjfrob}
In Setting~\ref{set:T}, let $\gamma$ be an element of \m{\GLk(W)}, $s$ a letter, $j$ a natural number, $G$ a matrix
in \m{\Mk(\cO)} and $i$, $i'$ indices. Write \m{G^{(j)}} for the matrix of the \m{\PT{T}^{j}}-th powers of the
entries of~$G$. Then $p$ divides
\[\Bigl(\letterConj{T}{\gamma}{s}\bigl(\bigl(\letterConj{T}{\gamma^{-1}}{s-j}(G)\bigr)^{(j)}\bigr)\Bigr)_{ii'}-\bigl(G_{ii'}\bigr)^{\PT{T}^{j}} .\]
\end{lemma}
\begin{proof}
Modulo~$p$ the map \m{\pi\mapsto\pi^{\PT{T}^{j}}} of~$\cO$ is a ring homomorphism, so it commutes with products of
matrices; and it sends \m{\iota_{s-j}(\upsilon)} to \m{\iota_s(\upsilon)}, by Lemma~\ref{l:basic}(a). Hence
\m{(\iota_{s-j}(\gamma^{-1})\,G\,\iota_{s-j}(\gamma))^{(j)}\equiv\iota_s(\gamma^{-1})\,G^{(j)}\,\iota_s(\gamma)}
modulo~$p$.
\end{proof}
\leanline{\lean{cjfrob}}

\begin{proposition}[exponents of the Frobenius operator track the excess]\label{p:powdvd}
In Setting~\ref{set:T}, let \m{i_0}, \m{i_1} be indices with \m{i_0+1=i_1}, let $s_0$ be a letter, and let
$\lambda$, $n$ be natural numbers with \m{L\,n\,k\,(\lambda+2)\le\PT{T}^{\lambda}} and \m{\lambda\le L}. Let
$\gamma$ be an element of \m{\GLk(W)}, let \m{T'} be the twisting data \m{\Tmap{\chi}{T}} for the inclusion~$\chi$
of~$\cO$ in the polynomial ring in the generic variables \m{x_{s,i,i'}}, and suppose that $\psi$ is given by
\[\psi(s,j)=\letterConj{T'}{\gamma}{s}\Bigl(\bigl(\letterConj{T'}{\gamma^{-1}}{s-j}(X_{s-j})\bigr)^{(j)}\Bigr)\]
for every letter~$s$ and every natural number~$j$, in the notation of Lemma~\ref{l:cjfrob}. Let $c$ be a natural
number, \m{\beta,\beta'} families, and $\nu$ an exponent on the generic variables, with parts \m{\nu_s} on the
variables of the letter~$s$. With \m{\mathrm{rw}} and \m{\mathrm{cw}} as in Proposition~\ref{p:mean} for the function
\m{\omega_0=\idxw{T}{\lambda}{i_0}{i_1}{\cdot}} on the indices, suppose that
\m{\mathrm{rw}(\nu_s)=\mathrm{cw}(\nu_s)} for every letter \m{s\ne s_0}, and that
\[2c+1<\bigl|\pr{\exw{T}{s_0+\lambda}{i_0}{i_1}{\cdot}}{\beta}-\mathrm{rw}(\nu_{s_0})\bigr| .\]
Then \m{\varpi^{c}} divides the coefficient of \m{x^{\nu}} in the \m{(\beta,\beta')} entry of
\m{\frob{T}{i_1}{\lambda}{n}{\psi}}.
\end{proposition}
\begin{proof}
Let \m{\mathrm{sg}=\pm1} be the sign of the difference in the hypothesis, and put
\[\mathrm{a}_m(i)=\omega_0(i)\cdot\tfrac{\mathrm{sg}}{2}\cdot\PT{T}^{m}/\qT{T};\qquad\text{so}\quad
|\mathrm{a}_m(i)|\le\tfrac12\PT{T}^{-\lambda}\ \text{ for } m\le L,\qquad \sum_i\mathrm{a}_m(i)=0 .\] Give the generic
variable \m{x_{s,i,i'}} the weight \m{\mathrm{a}_L(i)} if \m{s=s_0}, and
\m{\mathrm{a}_{[s-s_0]}(i)-\mathrm{a}_{[s-s_0]}(i')} otherwise, and let \m{\cG} be the corresponding filtration
\m{\gauss{\adic{\varpi}{\cdot}}{\cdot}{\cdot}} of the polynomial ring. The weight of \m{x^{\nu}} is then
\m{\tfrac{\mathrm{sg}}{2}\mathrm{rw}(\nu_{s_0})}, the other letters contributing multiples of
\m{\mathrm{rw}(\nu_s)-\mathrm{cw}(\nu_s)=0}. We apply Lemma~\ref{l:frobmem}, for \m{\thr=i_1}, with
\[\eta_{s,j}(i)=\mathrm{a}_{j+[s-j-s_0]}(i),\qquad\omega=\tfrac{\mathrm{sg}}{2}\,\exw{T}{s_0+\lambda}{i_0}{i_1}{\cdot},\qquad
\epsilon=\tfrac12\PT{T}^{-\lambda}.\]
Condition~(i) holds by the two properties of~\m{\mathrm{a}_m}. Condition~(ii) is the definition of the excess
weight at the letter \m{s_0+\lambda}, read at the two root variables of each letter, using \m{i_0<i_1}: these
are the variables \m{(i_1,s)} and \m{(i_0,s)}, of which the deep one for the threshold~\m{i_1} is the root variable
\m{\rootVar{T}{i_0}{i_1}{s}}. For condition~(iii) let \m{j\le\lambda}. By Lemma~\ref{l:cjfrob} for~\m{T'}, the entry
\m{\psi(s,j)_{ii'}} is \m{x_{s-j,i,i'}^{\,\PT{T}^{j}}} plus $p$ times a polynomial; and it is homogeneous of degree
\m{\PT{T}^{j}} in the variables of the letter \m{s-j}. The weights of these variables are at most
\m{\PT{T}^{-\lambda}} in absolute value, so every monomial of the entry has weight within
\m{2\PT{T}^{j-\lambda}\le2} of the weight of the leading one; as \m{p=\varpi^{2}}, the entry lies in~$\cG$ at
the level of its leading monomial. That level is \m{\PT{T}^{j}} times the weight of \m{x_{s-j,i,i'}}, which is
\m{\eta_{s,j+1}(i)-\eta_{s,j}(i')} if \m{s-j\ne s_0} and \m{\eta_{s,j+1}(i)-\eta_{s,j}(i')+\mathrm{a}_j(i')} if
\m{s-j=s_0}; and \m{|\mathrm{a}_j(i')|\le\epsilon}. The lemma gives the level
\m{\tfrac{\mathrm{sg}}{2}\pr{\exw{T}{s_0+\lambda}{i_0}{i_1}{\cdot}}{\beta}-L\,n\,k\,(\lambda+2)\epsilon} for the entry,
where the last term is at most~\m{\tfrac12} by hypothesis. So the coefficient of \m{x^{\nu}} is divisible
by~$\varpi$ to a power that is at least
\m{\tfrac12\bigl|\pr{\exw{T}{s_0+\lambda}{i_0}{i_1}{\cdot}}{\beta}-\mathrm{rw}(\nu_{s_0})\bigr|-\tfrac12>c}.
\end{proof}
\leanline{\lean{powdvd}\sep \lean{frobmem}\sep \lean{cjfrob}}

\begin{lemma}[degrees in the Frobenius operator]\label{l:degle}
In Setting~\ref{set:T}, let $\thr$, $\lambda$, $n$ be natural numbers, \m{\beta,\beta'} families, and $\nu$ an
exponent on the generic variables. If the coefficient of \m{x^{\nu}} in the \m{(\beta,\beta')} entry of
\m{\frob{T}{\thr}{\lambda}{n}{\frobX{T}{\cdot}{\cdot}}} is not zero, then for every letter~$s'$ the part of~$\nu$ on
the variables of~$s'$ has degree at most \m{L\,n\,k\,(\lambda+2)\,\PT{T}^{\lambda}}.
\end{lemma}
\begin{proof}
The entries of \m{\frobX{T}{s}{j}} with \m{j\le\lambda} have degree at most~\m{\PT{T}^{\lambda}}; those of the
products \m{\Pi^{m}_s}, \m{m\le\lambda+1}, and of their adjugates have degree at most
\m{k(\lambda+1)\PT{T}^{\lambda}}; an entry of the $n$-th symmetric power has $n$ times that degree at most, and an
entry of the Kronecker product over $L$ letters $L$ times as much.
\end{proof}
\leanline{\lean{degle}\sep the formal proof is Lemma~\ref{l:frobmem} for a filtration that records the degree in the
variables of one letter}

\begin{lemma}[products of linear forms]\label{l:tensor}
Let $\cS$ be a finite set, $\cV$ a set and $\cO$ a ring. There is a map, linear in each argument, that sends a
family \m{(\cL_{\fs})_{\fs\in\cS}} of $\cO$-linear forms on \m{\cO[x_{\cV}]} to an $\cO$-linear form
\m{\bigotimes_{\fs}\cL_{\fs}} on the polynomial ring in the variables \m{x_{\fs,v}}, such that
\m{(\bigotimes_{\fs}\cL_{\fs})(\prod_{\fs}\Phi_{\fs})=\prod_{\fs}\cL_{\fs}(\Phi_{\fs})} whenever each \m{\Phi_{\fs}} is a
polynomial in the variables \m{x_{\fs,v}}, \m{v\in\cV}, read as a polynomial in \m{\cO[x_{\cV}]}.
\end{lemma}
\begin{proof}
Define the form on a monomial \m{x^{\nu}}, with parts \m{\nu_{\fs}}, as \m{\prod_{\fs}\cL_{\fs}(x^{\nu_{\fs}})}, and
extend linearly.
\end{proof}
\leanline{\lean{tensordual}}

\begin{proof}[Proof of Theorem 4]
Write \m{q=\qT{T}=p^{oL}}. The hypotheses give
\begin{equation}\label{e:nzero}
n_0+k<q\quad\text{for}\quad n_0=L\,n\,k\,(\lambda+2)\,\PT{T}^{\lambda},\qquad n<q,\qquad
L\,n\,k\,(\lambda+2)\le\PT{T}^{\lambda},
\end{equation}
because \m{(L\,n\,k\,(\lambda+2)+n+k+1)\,\PT{T}^{\lambda}\le\PT{T}^{2\lambda}=q}.

\emph{The diagonal form.} Let \m{w_0} be the function of Lemma~\ref{l:diagform} for the $k$ indices as
variables, the prime~$p$ and the exponent~\m{oL}, and put \m{w(\beta)=\prod_sw_0(\beta_s)}. As \m{n<q}, each
\m{w_0(\beta_s)} divides \m{p^{k\,oL}}; choose \m{w'(\beta)} with \m{w(\beta)\,w'(\beta)=p^{k\,oL\cdot L}}.
Property~(iii) of the theorem is Lemma~\ref{l:diagform}(ii) for each letter, since
\m{\letterRep{T}{n}{\gamma}_{\beta',\beta}=\prod_s\mathrm{Sym}^{n}(\letterMatrix{T}{\gamma}{s})_{\beta'_s,\beta_s}} and the
letter matrices of \m{\gamma^{\tp}} are the transposes of those of~$\gamma$.

\emph{The operator.} Let \m{\mathrm{Fr}} be the matrix \m{\frob{T}{i_1}{\lambda}{n}{\frobX{T}{\cdot}{\cdot}}}, with
entries in the polynomial ring in the generic variables. For a family \m{\cL=(\cL_s)_s} of $\cO$-linear forms on
the polynomials in one generic matrix let \m{Z(\cL)} be the matrix over~$\cO$ obtained by applying the form
\m{\bigotimes_s\cL_s} of Lemma~\ref{l:tensor} to the entries of~\m{\mathrm{Fr}}; it is linear in each~\m{\cL_s}.

\emph{The means.} For every letter~$s$ let \m{\cM^{+}_s} and \m{\cM^{-}_s} be the forms of
Proposition~\ref{p:mean} for the homomorphism \m{\iota_s}, the labelling \m{\twoBlock{k}{i_1}{\cdot}}, the degree
\m{n_0}, the prime~$p$ with the exponent~\m{oL}, the function \m{\omega_0=\idxw{T}{\lambda}{i_0}{i_1}{\cdot}}, and the
set~$\Omega$ of all real numbers, respectively of the real numbers greater than \m{h-(2c+1)}; the hypotheses
hold by~\eqref{e:nzero}. Put \m{\cL^{\circ}_s=\cM^{+}_s-\cM^{-}_s} and
\[Z=Z(\cL^{\circ})\cdot\operatorname{diag}(w'(\beta))_{\beta}.\]
By (a), (b), (c) of the proposition, \m{\cL^{\circ}_s} is invariant under \m{\conjSubst{\iota_s}{\gamma}} for
$\gamma$ in the Levi subgroup, and \m{\cL^{\circ}_s(x^{\alpha})=0} unless
\m{\mathrm{rw}(\alpha)=\mathrm{cw}(\alpha)\le h-(2c+1)}.

\emph{A criterion.} Let $s_0$ be a letter, $\gamma$ an element of \m{\Lev_{\twoBlock{k}{i_1}{\cdot}}(W)}, and \m{\cL},
\m{\cL'} families of forms with \m{\cL'_s=\cL_s\circ\conjSubst{\iota_s}{\gamma}} for every letter. Let \m{F_1}, \m{F_2}
be forms. Suppose that for every family~$\beta$ at which the row vector \m{F_1^{\tp}\letterRep{T}{n}{\gamma^{-1}}}
has a coordinate not divisible by~\m{\varpi^{c}}, and for every exponent~$\nu$ on the generic variables, either
\m{\prod_s\cL_s(x^{\nu_s})=0}, or the two hypotheses of Proposition~\ref{p:powdvd} hold for $s_0$, $\beta$
and~$\nu$. Then \m{\varpi^{c}} divides \m{F_1^{\tp}Z(\cL')F_2}. Indeed, \m{\bigotimes_s\cL'_s} is
\m{\bigotimes_s\cL_s} after the substitution of \m{\iota_s(\gamma^{-1})X_s\iota_s(\gamma)} for each~\m{X_s}. By
Lemma~\ref{l:frobmap}, for the twisting data~\m{T'} and the substitution as the ring homomorphism, the
substituted matrix \m{\mathrm{Fr}} is
\m{\letterRep{T}{n}{\gamma^{-1}}\cdot\frob{T}{i_1}{\lambda}{n}{\psi}\cdot\letterRep{T}{n}{\gamma}} with $\psi$ as in
Proposition~\ref{p:powdvd}; the constant matrices pass through the linear form. So
\[F_1^{\tp}Z(\cL')F_2=\sum_{\beta,\beta'}\bigl(F_1^{\tp}\letterRep{T}{n}{\gamma^{-1}}\bigr)(\beta)\cdot
\Bigl(\sum_{\nu}\mathrm{cf}_{\nu}\prod_s\cL_s(x^{\nu_s})\Bigr)\cdot\bigl(\letterRep{T}{n}{\gamma}F_2\bigr)(\beta'),\]
with \m{\mathrm{cf}_{\nu}} the coefficient of \m{x^{\nu}} in the \m{(\beta,\beta')} entry of
\m{\frob{T}{i_1}{\lambda}{n}{\psi}}, and every term is divisible by~\m{\varpi^{c}}, by the hypothesis or by
Proposition~\ref{p:powdvd}.

\emph{Property (i).} Let $F$ be a form, $\gamma$ in the Levi subgroup, and $\beta$ a family whose excess at some
letter~$s_1$ exceeds~$h$. The coordinate \m{(\letterRep{T}{n}{\gamma}ZF)(\beta)} is \m{F_1^{\tp}Z(\cL^{\circ})F_2} for
\m{F_1^{\tp}} the $\beta$-th row of \m{\letterRep{T}{n}{\gamma}} and \m{F_2=w'F}. Apply the criterion with
\m{\cL=\cL'=\cL^{\circ}} and \m{s_0=s_1-\lambda}. The row vector \m{F_1^{\tp}\letterRep{T}{n}{\gamma^{-1}}} is the
$\beta$-th unit vector. If \m{\prod_s\cL^{\circ}_s(x^{\nu_s})\ne0}, then every \m{\nu_s} has equal row and column
weights, at most \m{h-(2c+1)}, and the excess of~$\beta$ at \m{s_0+\lambda} exceeds \m{\mathrm{rw}(\nu_{s_0})} by more
than \m{2c+1}.

\emph{Property (ii).} Let \m{F'} lie in \m{\exc{T}{n}{i_0}{i_1}{\varpi^{c}\cO}{h-(4c+2)}}, let $F$ be a form, and
put $\mathrm{bil}(\cL)=(wF')^{\tp}Z(\cL)(w'F)$, with the products \m{wF'} and \m{w'F} taken coordinate by
coordinate.
The quantity of~(ii) is \m{\mathrm{bil}(\cL^{\circ})} less \m{p^{(k+k^{2})oL\cdot L}\sum_{\beta}F'(\beta)F(\beta)}.
Expanding in each letter, \m{\mathrm{bil}(\cL^{\circ})} is the sum, over the subsets~$\cS$ of the letters, of
\m{\mathrm{bil}(\cL^{\cS})} with \m{\cL^{\cS}_s=-\cM^{-}_s} for $s$ in~$\cS$ and \m{\cL^{\cS}_s=\cM^{+}_s} otherwise.

For the empty subset: by Lemma~\ref{l:degle} and~(c), \m{\cM^{+}_s} takes at every monomial that occurs in
\m{\mathrm{Fr}} the value \m{p^{k^{2}oL}} times its value at the identity; and \m{\mathrm{Fr}} with all generic
matrices equal to the identity is the identity matrix. So \m{Z(\cM^{+})=p^{k^{2}oL\cdot L}\cdot1}, and
\m{\mathrm{bil}(\cM^{+})=p^{k^{2}oL\cdot L}\sum_{\beta}w(\beta)w'(\beta)F'(\beta)F(\beta)}, which is the term
subtracted.

For a nonempty subset choose \m{s_0} in it. By~(d), \m{\cM^{-}_{s_0}} is a combination of forms
\m{\cL_0\circ\conjSubst{\iota_{s_0}}{\gamma}} with $\gamma$ in the Levi subgroup and \m{\cL_0} vanishing at the
\m{x^{\alpha}} with \m{\mathrm{rw}(\alpha)\le h-(2c+1)}. So it suffices that \m{\varpi^{c}} divides
\m{\mathrm{bil}(\cL')} when \m{\cL'_{s_0}=\cL_0\circ\conjSubst{\iota_{s_0}}{\gamma}} and \m{\cL'_s=\cL^{\cS}_s} for
\m{s\ne s_0}. Apply the criterion with \m{\cL_{s_0}=\cL_0}, \m{\cL_s=\cL'_s} otherwise, which is allowed by the
invariance~(a), with \m{F_1=wF'} and \m{F_2=w'F}. By property~(iii),
\m{((wF')^{\tp}\letterRep{T}{n}{\gamma^{-1}})(\beta)=w(\beta)\cdot(\letterRep{T}{n}{(\gamma^{-1})^{\tp}}F')(\beta)},
and \m{(\gamma^{-1})^{\tp}} lies in the Levi subgroup; so, by the choice of~\m{F'}, this coordinate is divisible
by~\m{\varpi^{c}} unless the excess of~$\beta$ is at most \m{h-(4c+2)} at every letter. For such $\beta$, and
$\nu$ with \m{\prod_s\cL_s(x^{\nu_s})\ne0}: the parts \m{\nu_s}, \m{s\ne s_0}, have equal row and column weights
by~(b), and \m{\mathrm{rw}(\nu_{s_0})>h-(2c+1)}, which exceeds the excess of~$\beta$ at \m{s_0+\lambda} by more
than \m{2c+1}.
\end{proof}
\leanline{\lean{projector}\sep \lean{mean}\sep \lean{powdvd}\sep \lean{frobmap}\sep \lean{degle}\sep \lean{tensordual}\sep
\lean{mapC} (a linear form applied to the entries commutes with multiplication by constant matrices)}

\sect{Elementary divisors}{s:div}
In this section $\cO$ is a valuation domain, with maximal ideal~$\fm$, and \m{\bar T=\Tmap{\mathrm{res}}{T}} is the
twisting data over the residue field \m{\cO/\fm} obtained from~$T$ by the residue map
\m{\mathrm{res}\colon\cO\to\cO/\fm}. For a finite set~$\cX$ a vector in \m{\cO^{\cX\times\Fam}} has the \kw{blocks}
\m{F(\fx,\cdot)}, forms in \m{\cO^{\Fam}}.

\begin{lemma}[independent columns]\label{l:cardle}
In Setting~\ref{set:T}, let $\cO$ be a valuation domain, $n$ a natural number, $\cX$ a finite set and $\cN$ a
set. Let $A$ be a matrix over~$\cO$ with rows indexed by \m{\cX\times\Fam} and columns by~$\cN$, and $A'$ a matrix
with rows indexed by~$\cN$ and columns by \m{\cX\times\Fam}. Let $\cS$ be a finite subset of~$\cN$, $\tau$ a
labelling and $\mathcal{B}$ a subset of \m{\Fam}. Suppose that \m{(A'A)_{l,l'}} is $1$ if \m{l=l'} and $0$ otherwise,
for every \m{l'} in~$\cS$ and every $l$ in~$\cN$, and that for every \m{l'} in~$\cS$ every block of the column
\m{l'} of~$A$ lies in \m{\trl{T}{n}{\tau}{\fm}{\mathcal{B}}}. Then
\[\#\cS\ \le\ \#\cX\cdot\dim_{\cO/\fm}\trl{\bar T}{n}{\tau}{0}{\mathcal{B}} .\]
\end{lemma}
\begin{proof}
Reduce the columns \m{l'\in\cS} of~$A$ modulo~$\fm$. Their blocks lie in \m{\trl{\bar T}{n}{\tau}{0}{\mathcal{B}}},
by Lemma~\ref{l:basic}(d); and they are linearly independent over \m{\cO/\fm}, because the reduction of~$A'$
recovers the coefficients of a combination of them.
\end{proof}
\leanline{\lean{cardle}}

\begin{proposition}[two matrices that pair almost perfectly]\label{p:pairing}
In Setting~\ref{set:T}, let $\cO$ be a valuation domain, $n$ a natural number, and $\cX$, $\cN$ finite sets. Let
$A$ and $B$ be matrices over~$\cO$ with rows indexed by \m{\cX\times\Fam} and columns by~$\cN$, and let $A'$ and
$C$ be matrices with rows indexed by~$\cN$ and columns by \m{\cX\times\Fam}. Let \m{\pi_1}, \m{\pi_2}, \m{\pi_0} be
elements of~$\cO$, $\tau$ a labelling and $\mathcal{B}$ a subset of \m{\Fam}. Suppose that
\begin{enumerate}[label=\textup{(\roman*)},leftmargin=2.3em,itemsep=1pt,topsep=2pt]
\item \m{A'A=1};
\item \m{\pi_0} is not a unit;
\item \m{\pi_1\pi_0} divides every entry of \m{CB-\pi_1\cdot1};
\item whenever \m{AF_1-BF_2=\pi_1\pi_0\,F} for vectors \m{F_1}, \m{F_2} in \m{\cO^{\cN}} and $F$ in
\m{\cO^{\cX\times\Fam}}, every block of~$F$ lies in \m{\trl{T}{n}{\tau}{\fm}{\mathcal{B}}}.
\end{enumerate}
Then there is a natural number~$l$ with
\[\#\cN+l\ \le\ \#\cX\cdot\dim_{\cO/\fm}\trl{\bar T}{n}{\tau}{0}{\mathcal{B}}\qquad\text{and}\qquad
\IndexDvd{\cO^{\cN}}{CA(\cO^{\cN})+\pi_2\,\cO^{\cN}}{\pi_1^{\#\cN}\pi_2^{\,l}} .\]
\end{proposition}
\begin{proof}
Write \m{\mathcal{Y}=\cX\times\Fam}. By Input~2(b), for the valuation pair \m{(\cO,\cO)}, there are invertible
matrices $U$, $U'$ over~$\cO$ and elements \m{\varepsilon_y} of~$\cO$, \m{y\in\mathcal{Y}}, with
\begin{equation}\label{e:smith}
(AA'B-B)\,A'=U\operatorname{diag}(\varepsilon_y)\,U' .
\end{equation}
Multiplying by~$A'$ on the left and using~(i) gives \m{A'U\operatorname{diag}(\varepsilon_y)=0}; so the column~$y$
of \m{A'U} is zero whenever \m{\varepsilon_y\ne0}. Let \m{\mathcal{Y}_0} be the set of the $y$ with
\m{\varepsilon_y\ne0} and \m{\varepsilon_y} dividing \m{\pi_1\pi_0}, and \m{l=\#\mathcal{Y}_0}.

\emph{The inequality.} Apply Lemma~\ref{l:cardle} to the matrix \m{(A\ U)}, with columns indexed by the disjoint
union of $\cN$ and~$\mathcal{Y}$, to the matrix with the rows of~$A'$ above those of \m{U^{-1}-U^{-1}AA'}, and to
the columns in $\cN$ and in~\m{\mathcal{Y}_0}. The product of the two matrices has the blocks $1$, \m{A'U}, $0$ and
\m{1-U^{-1}AA'U}; its columns in $\cN$ and in~\m{\mathcal{Y}_0} are unit vectors, by the remark on \m{A'U}. The
blocks of a column \m{Ae_{l'}} lie in \m{\trl{T}{n}{\tau}{\fm}{\mathcal{B}}} by~(iv) with \m{F_1=\pi_1\pi_0e_{l'}} and
\m{F_2=0}. For $y$ in~\m{\mathcal{Y}_0} write \m{\pi_1\pi_0=\varepsilon_y\pi'}; by~\eqref{e:smith}, with
\m{F_2=\pi'A'U'^{-1}e_y}, we have \m{A(A'BF_2)-BF_2=\pi_1\pi_0\,Ue_y}, and (iv) applies to the column \m{Ue_y}.

\emph{The index.} For every~$y$ one of \m{\pi_1\pi_0}, \m{\varepsilon_y} divides the other; choose \m{\varepsilon'_y}
with \m{\pi_1\pi_0\varepsilon'_y=\varepsilon_y}, or with \m{\varepsilon_y\varepsilon'_y=\pi_1\pi_0} when the former is
impossible. Put \m{C_1=CU\operatorname{diag}(\pi_1\pi_0\varepsilon'_y-\varepsilon_y)}; its column~$y$ is zero for $y$
outside~\m{\mathcal{Y}_0}. Write \m{CB-\pi_1=\pi_1\pi_0C_0} by~(iii). Using \eqref{e:smith} and~(i),
\[C_1U'A+CAA'B=\pi_1\pi_0\,CU\operatorname{diag}(\varepsilon'_y)U'A+CB
=\pi_1\bigl(1+\pi_0\,(C_0+CU\operatorname{diag}(\varepsilon'_y)U'A)\bigr).\]
The determinant of the last bracket is congruent to~$1$ modulo \m{\pi_0}, hence a unit by~(ii); so the determinant
of the left side divides \m{\pi_1^{\#\cN}}, and its image lies in \m{C_1(\cO^{\mathcal{Y}})+CA(\cO^{\cN})}. By
Input~1(a),
\m{\IndexDvd{\cO^{\cN}}{C_1(\cO^{\mathcal{Y}})+CA(\cO^{\cN})+\pi_2\cO^{\cN}}{\pi_1^{\#\cN}}}. The diagonal matrix on
\m{\cO^{\mathcal{Y}}} with the entries \m{\pi_2} at \m{\mathcal{Y}_0} and $1$ elsewhere has determinant
\m{\pi_2^{\,l}}, and $C_1$ maps its image into \m{\pi_2\cO^{\cN}}. By Input~1(a) and~(c), for the map~$C_1$,
\m{\IndexDvd{C_1(\cO^{\mathcal{Y}})+CA(\cO^{\cN})+\pi_2\cO^{\cN}}{CA(\cO^{\cN})+\pi_2\cO^{\cN}}{\pi_2^{\,l}}}, and
Input~1(b) gives the claim.
\end{proof}
\leanline{\lean{pairing}\sep \lean{cardle}\sep \lean{smith}\sep \lean{ofdetdvd}\sep \lean{idmap}\sep \lean{idtrans}}

\begin{proposition}[the depth factorisation]\label{p:depthfact}
In Setting~\ref{set:T}, let $\cO$ be a valuation domain of characteristic zero, $c$ and $n$ natural numbers,
\m{i_0}, \m{i_1} indices with \m{i_0<i_1}, and $j_0$ a natural number with \m{\dep{T}{k}{n}{i_1}{\beta}\le j_0} for
every family~$\beta$. Let $\cX$ and $\cE$ be finite sets, \m{\mathrm{pt}\colon\cE\to\cX} a map and
\m{(\gamma_{\fe})_{\fe\in\cE}} elements of \m{\GLk(W)} such that the map
\m{\fe\mapsto(\mathrm{pt}(\fe),\gamma_{\fe}^{-1}\Par_{\twoBlock{k}{i_1}{\cdot}}(W))} from $\cE$ to
\m{\cX\times\GLk(W)/\Par_{\twoBlock{k}{i_1}{\cdot}}(W)} is surjective. Let $J$ be the matrix with rows indexed by
\m{\cE\times\Fam} and columns by \m{\cX\times\Fam} whose entry at \m{((\fe,\beta),(\fx,\beta'))} is
\m{\varpi^{\dep{T}{k}{n}{i_1}{\beta}}\cdot\letterRep{T}{n}{\gamma_{\fe}}_{\beta,\beta'}} if \m{\mathrm{pt}(\fe)=\fx}, and
$0$ otherwise. Then there are a matrix~$A$ of the shape of~$J$, a square matrix~$\Delta$ indexed by
\m{\cX\times\Fam} and a matrix~$A'$ of the shape of the transpose of~$J$, all over~$\cO$, such that:
\begin{enumerate}[label=\textup{(\roman*)},leftmargin=2.3em,itemsep=1pt,topsep=2pt]
\item \m{A\Delta=J};
\item \m{\det\Delta} divides $\varpi$ to the power
\m{\#\cX\cdot\sum_{j<j_0}\dim_{\cO/\fm}\deepForms{\bar T}{k}{n}{i_1}{j}};
\item \m{A'A=1};
\item for every $F$ in \m{\cO^{\cX\times\Fam}}, every $\fe$ in~$\cE$ and every real number~$h$ with
\m{2c\le h}, the block \m{(AF)(\fe,\cdot)} lies in \m{\exc{T}{n}{i_0}{i_1}{\varpi^{c}\cO}{h}}.
\end{enumerate}
\end{proposition}
\begin{proof}
Write \m{\mathcal{Y}=\cX\times\Fam}, \m{\Par=\Par_{\twoBlock{k}{i_1}{\cdot}}(W)} and
\m{\mathrm{d}(\beta)=\dep{T}{k}{n}{i_1}{\beta}}. So
\[(JF)(\fe,\beta)=\varpi^{\mathrm{d}(\beta)}\,\bigl(\letterRep{T}{n}{\gamma_{\fe}}F(\mathrm{pt}(\fe),\cdot)\bigr)(\beta).\]

\emph{Step 1: divisibility passes to all translates.} Let $\pi$ in~$\cO$ divide every coordinate of~\m{JF}. Then
$\pi$ divides \m{\varpi^{\mathrm{d}(\beta)}(\letterRep{T}{n}{\gamma}F(\fx,\cdot))(\beta)} for every $\fx$, every
$\gamma$ in \m{\GLk(W)} and every~$\beta$. Indeed, choose $\fe$ with \m{\mathrm{pt}(\fe)=\fx} and
\m{\gamma'=\gamma\gamma_{\fe}^{-1}} in~\m{\Par}. For a letter~$s$ and indices $i$, $i'$ such that \m{(i,s)} is not
deep and \m{(i',s)} is deep, for \m{\thr=i_1}, the entry \m{(\letterMatrix{T}{\gamma'}{s})_{ii'}} is divisible
by~\m{p=\varpi^{2}} (Lemma~\ref{l:par}). So there are matrices \m{G_s} over~$\cO$ with
\m{\operatorname{diag}(\varpi^{\dpi{T}{i_1}{(i,s)}})_i\cdot\letterMatrix{T}{\gamma'}{s}=G_s\cdot\operatorname{diag}(\varpi^{\dpi{T}{i_1}{(i,s)}})_i},
and by Lemma~\ref{l:basic}(e),
\m{\operatorname{diag}(\varpi^{\mathrm{d}(\beta)})\,\letterRep{T}{n}{\gamma'}=\bigl(\bigotimes_s\mathrm{Sym}^{n}(G_s)\bigr)\operatorname{diag}(\varpi^{\mathrm{d}(\beta)})}.
Hence the vector \m{(\varpi^{\mathrm{d}(\beta)}(\letterRep{T}{n}{\gamma}F(\fx,\cdot))(\beta))_{\beta}} is the image of
the block \m{(JF)(\fe,\cdot)} under a matrix over~$\cO$.

\emph{Step 2: the factorisation.} The map \m{\mathrm{pt}} is surjective; a section of it makes \m{\mathcal{Y}} a
subset of \m{\cE\times\Fam}. Input~2(a), for \m{(\cO,\cO)} and the weights~$0$, gives an invertible matrix~$U$ and
a square matrix~$\Delta$ with \m{J=U_{\cdot,\mathcal{Y}}\Delta}. Put \m{A=U_{\cdot,\mathcal{Y}}} and
\m{A'=(U^{-1})_{\mathcal{Y},\cdot}}; then (i) and (iii) hold. By Input~2(b),
\m{\Delta=U_1\operatorname{diag}(\varepsilon_y)U_2} with $U_1$, $U_2$ invertible over~$\cO$. The column~$y$ of
\m{JU_2^{-1}=AU_1\operatorname{diag}(\varepsilon_y)} is divisible by~\m{\varepsilon_y}; so by Step~1, for all $\fx$,
$\gamma$, $\beta$,
\begin{equation}\label{e:epsdiv}
\varepsilon_y\ \text{ divides }\ \varpi^{\mathrm{d}(\beta)}\,\bigl(\letterRep{T}{n}{\gamma}\,(U_2^{-1}e_y)(\fx,\cdot)\bigr)(\beta).
\end{equation}
In particular \m{\varepsilon_y} divides \m{\varpi^{j_0}}: this element is the \m{(y,y)} entry of
\m{U_2\cdot\varpi^{j_0}\cdot U_2^{-1}}, a sum of terms each divisible by~\m{\varepsilon_y}, by~\eqref{e:epsdiv} for
\m{\gamma=1} and \m{\mathrm{d}(\beta)\le j_0}. Choose \m{\varepsilon'_y} with \m{\varepsilon_y\varepsilon'_y=\varpi^{j_0}}.

\emph{Step 3: the determinant.} Up to a unit, \m{\det\Delta=\prod_y\varepsilon_y}. For each $y$ let \m{m_y} be the
number of the \m{j<j_0} such that \m{\varepsilon_y} does not divide~\m{\varpi^{j}}. These $j$ form an initial
segment, and \m{\varepsilon_y} divides \m{\varpi^{j_0}}; so \m{\varepsilon_y} divides \m{\varpi^{m_y}}, and
\m{\sum_ym_y} is the sum over \m{j<j_0} of the number of the $y$ with \m{\varepsilon_y\nmid\varpi^{j}}. For a
fixed~$j$, apply Lemma~\ref{l:cardle} to the matrices \m{U_2^{-1}} and~\m{U_2}, to these~$y$, the constant
labelling~$0$ and the set~$\mathcal{B}$ of the $\beta$ with \m{\mathrm{d}(\beta)<j}: if some coordinate
in~\eqref{e:epsdiv} with \m{\mathrm{d}(\beta)<j} had a unit as second factor, \m{\varepsilon_y} would divide
\m{\varpi^{j}}. So the number of these $y$ is at most \m{\#\cX\cdot\dim\deepForms{\bar T}{k}{n}{i_1}{j}}. This
is~(ii).

\emph{Step 4: the excess condition.} Let $F$, $\fe$ and \m{h\ge2c} be as in~(iv), let $\gamma$ lie in
\m{\Lev_{\twoBlock{k}{i_1}{\cdot}}(W)} and let $\beta$ be a family whose excess at a letter~$s_0$ exceeds~$h$. Put
\m{F^{\circ}=U_2^{-1}\operatorname{diag}(\varepsilon'_y)U_1^{-1}F}; then \m{JF^{\circ}=\varpi^{j_0}AF}. Let
\m{F_1=F^{\circ}(\mathrm{pt}(\fe),\cdot)}. The letter matrices of~$\gamma$ are block diagonal, so
\m{\letterRep{T}{n}{\gamma}} commutes with \m{\operatorname{diag}(\varpi^{\mathrm{d}(\beta)})}, and
\[\varpi^{j_0}\cdot\bigl(\letterRep{T}{n}{\gamma}(AF)(\fe,\cdot)\bigr)(\beta)=\varpi^{\mathrm{d}(\beta)}\cdot\bigl(\letterRep{T}{n}{\gamma\gamma_{\fe}}F_1\bigr)(\beta).\]
As \m{\varpi\ne0} (Lemma~\ref{l:basic}), we have to show that \m{\varpi^{j_0+c}} divides the right side.

Let \m{\omega=\dpi{T}{i_1}{\cdot}-\tfrac12\exw{T}{s_0}{i_0}{i_1}{\cdot}}, a weight on the variables. For a
family~$\beta'$ let \m{\mathrm{ord}(\beta')} be the largest natural number, at most
\m{\max(\lceil j_0-\pr{\omega}{\beta'}\rceil,0)}, such that \m{\varpi^{\mathrm{ord}(\beta')}} divides
\m{(\letterRep{T}{n}{\gamma'}F_1)(\beta')} for every \m{\gamma'} in \m{\GLk(W)}. Choose a family~\m{\beta_0} for
which \m{\theta_0=\mathrm{ord}(\beta_0)+\pr{\omega}{\beta_0}} is least. Then, for every \m{\gamma'} and
every~$\beta'$,
\begin{equation}\label{e:least}
\varpi\ \text{ to the power }\ \max(\lceil\theta_0-\pr{\omega}{\beta'}\rceil,0)\ \text{ divides }\
(\letterRep{T}{n}{\gamma'}F_1)(\beta'),
\end{equation}
because the natural number \m{\mathrm{ord}(\beta')} is at least \m{\theta_0-\pr{\omega}{\beta'}}. We claim that
\m{\theta_0\ge j_0}. If \m{\mathrm{ord}(\beta_0)} is the bound in its definition, this is clear. Otherwise there is
\m{\gamma_0} such that \m{\varpi^{\mathrm{ord}(\beta_0)+1}} does not divide
\m{(\letterRep{T}{n}{\gamma_0}F_1)(\beta_0)}. On the one hand, Step~1 for \m{F^{\circ}} and \m{\pi=\varpi^{j_0}} shows
that \m{\varpi^{j_0}} divides \m{\varpi^{\mathrm{d}(\beta_0)}(\letterRep{T}{n}{\gamma_0}F_1)(\beta_0)}; so
\m{j_0\le\mathrm{ord}(\beta_0)+\mathrm{d}(\beta_0)}. On the other hand, Theorem~1 applies to the indices \m{i_0\ne i_1},
the letter~$s_0$, the constant labelling~$0$, the weight~$\omega$, the polynomial of the form
\m{\letterRep{T}{n}{\gamma_0}F_1}, the exponent~\m{\beta_0} and the order \m{e=\mathrm{ord}(\beta_0)}: its hypothesis
on the translates is~\eqref{e:least}, by Lemma~\ref{l:basic}(c), and the hypothesis on~$\omega$ holds because of
the two root variables of a letter exactly \m{\rootVar{T}{i_0}{i_1}{s}} is deep for \m{\thr=i_1}, so that the
difference in question is \m{1} less half a difference of two excess weights, which lies in \m{[0.4,1.6]} by
Lemma~\ref{l:exIcc}. The theorem gives \m{\pr{\exw{T}{s_0}{i_0}{i_1}{\cdot}}{\beta_0}\le0}, that is,
\m{\pr{\omega}{\beta_0}\ge\mathrm{d}(\beta_0)}; and so \m{\theta_0\ge\mathrm{ord}(\beta_0)+\mathrm{d}(\beta_0)\ge j_0}.

Finally, by~\eqref{e:least} for \m{\gamma'=\gamma\gamma_{\fe}} and~$\beta$, the coordinate
\m{(\letterRep{T}{n}{\gamma\gamma_{\fe}}F_1)(\beta)} is divisible by~$\varpi$ to a power that is at least
\m{\theta_0-\pr{\omega}{\beta}\ge j_0-\mathrm{d}(\beta)+\tfrac12\pr{\exw{T}{s_0}{i_0}{i_1}{\cdot}}{\beta}>j_0-\mathrm{d}(\beta)+c},
as the excess of~$\beta$ exceeds \m{h\ge2c}.
\end{proof}
\leanline{\lean{depthfact}\sep \lean{rootlemma}\sep \lean{pdvdpar}\sep \lean{exwIcc}\sep \lean{hermite}\sep
\lean{smith}\sep \lean{cardle}}

\begin{proof}[Proof of Theorem 5]
\emph{Notation.} Put \m{\kappa_0=(k+k^{2})\,oL\cdot L} and \m{c=2\kappa_0+2}; the seventh size condition reads
\m{6c+2\le y}. Let \m{i_1} be the index~$\thr$ and \m{i_0=\thr-1}; they exist because \m{0<\thr} and \m{2\thr\le k}.
Let \m{j_0} be a bound for the depths \m{\dep{T}{k}{n}{\thr}{\beta}}. The residue field \m{\cO/\fm} has
characteristic~$p$, as $p$ is not a unit, and \m{\bar T} satisfies the same size conditions as~$T$; also
$\varpi$ is not a unit, and \m{\varpi\ne0}. Put \m{\cX'=\cH\times\cX}, \m{\cE'=\cH\times\cE} and
\m{\mathcal{Y}=\cX'\times\Fam}; then \m{\#\mathcal{Y}=\#\cX\cdot\#\cI}, by the bijection~$\jmath$. Write
\m{\Par=\Par_{\twoBlock{k}{\thr}{\cdot}}(W)}.

\emph{Two factorisations.} The map of the hypothesis is an injective map of the finite set~$\cE$ to itself,
hence bijective. So Proposition~\ref{p:depthfact} applies, with $c$, \m{(i_0,i_1)}, \m{j_0}, to the sets $\cX'$
and~$\cE'$ with
\[\mathrm{pt}_1(\fh,\fe)=(\fh,\mathrm{src}(\fe))\ \text{ and the elements } \gamma_{\fe}^{\tp};\qquad
\mathrm{pt}_2(\fh,(\fx,\fq))=(\fh,\fx)\ \text{ and the } \tilde{\fq}^{-1}.\]
Let \m{J_1=A_1\Delta_1} and \m{J_2=A_2\Delta_2} be the factorisations, with left inverses \m{A'_1}, \m{A'_2}.

\emph{The projector.} Let $w$ and $Z$ be as in Theorem~4 for $c$, \m{(i_0,i_1)}, $\lambda$ and \m{h=y}; its
hypotheses are the fifth and sixth size conditions. On vectors indexed by \m{\cE'\times\Fam} let \m{\mathbf{u}}
act by \m{u_{\fe}} on the coordinate in~$\cH$, for each $\fe$, let \m{\mathbf{Z}} act by~$Z$ and \m{\mathbf{w}} by
\m{\operatorname{diag}(w(\beta))} on the coordinate in \m{\Fam}.

\emph{The pairing.} We apply Proposition~\ref{p:pairing} to the sets $\cE'$ (in the place of~$\cX$) and
\m{\mathcal{Y}} (in the place of~$\cN$), the matrices
\[A=A_1,\qquad B=\mathbf{u}^{-1}\mathbf{Z}\,A_2'^{\tp},\qquad A'=A'_1,\qquad C=A_2^{\tp}\,\mathbf{u}\,\mathbf{w},\]
the elements \m{\pi_1=\varpi^{2\kappa_0}=p^{\kappa_0}}, \m{\pi_2=p^{e}}, \m{\pi_0=\varpi}, the labelling
\m{\twoBlock{k}{\thr}{\cdot}} and the set~$\mathcal{B}$ of the families whose excess for \m{(i_0,i_1)} exceeds~$y$ at
some letter; so the module of the conclusion is built from \m{\exc{\bar T}{n}{i_0}{i_1}{0}{y}}. Hypotheses (i)
and~(ii) hold. For~(iii): \m{CB=A_2^{\tp}\,\mathbf{w}\mathbf{Z}\,A_2'^{\tp}}, so its \m{(y_1,y_2)} entry is the sum
over $\fe'$ in~$\cE'$ of \m{\sum_{\beta}F'(\beta)w(\beta)(ZF)(\beta)} for the blocks \m{F'=(A_2e_{y_1})(\fe',\cdot)}
and \m{F=(A_2'^{\tp}e_{y_2})(\fe',\cdot)}. By Proposition~\ref{p:depthfact}(iv), and because
\m{2c\le y-(4c+2)}, the block \m{F'} lies in
\[\exc{T}{n}{i_0}{i_1}{\varpi^{c}\cO}{y-(4c+2)}.\]
By Theorem~4(ii) the entry is congruent modulo~\m{\varpi^{c}} to \m{p^{\kappa_0}} times the \m{(y_1,y_2)} entry of
\m{A_2^{\tp}A_2'^{\tp}=1}; and \m{\pi_1\pi_0=\varpi^{c-1}}. For~(iv): if \m{A_1F_1-BF_2=\varpi^{c-1}F}, then each
block of \m{A_1F_1} lies in \m{\exc{T}{n}{i_0}{i_1}{\varpi^{c}\cO}{y}} by Proposition~\ref{p:depthfact}(iv), and
each block of \m{BF_2} is $Z$ applied to a form, so it lies there by Theorem~4(i). Hence, for $\gamma$ in the
Levi subgroup and $\beta$ in~$\mathcal{B}$, \m{\varpi^{c}} divides \m{\varpi^{c-1}} times the coordinate of the
translate of a block of~$F$, and that coordinate lies in~$\fm$. The proposition gives a natural number~$l$ with
\begin{equation}\label{e:L}
\#\mathcal{Y}+l\le\#\cE'\cdot\dim_{\cO/\fm}\exc{\bar T}{n}{i_0}{i_1}{0}{y},\qquad
\IndexDvd{\cO^{\mathcal{Y}}}{CA_1(\cO^{\mathcal{Y}})+p^{e}\cO^{\mathcal{Y}}}{p^{\kappa_0\#\mathcal{Y}+el}} .
\end{equation}

\emph{The chain.} Let \m{e_0=\#\cX'\cdot\sum_{j<j_0}\dim_{\cO/\fm}\deepForms{\bar T}{k}{n}{\thr}{j}}; the
determinants of \m{\Delta_1} and of \m{\Delta_2^{\tp}} divide~\m{\varpi^{e_0}}. By Input~1(a),
\m{\IndexDvd{\cO^{\mathcal{Y}}}{\Delta_1(\cO^{\mathcal{Y}})+p^{e}\cO^{\mathcal{Y}}}{\varpi^{e_0}}}, and the same
for~\m{\Delta_2^{\tp}}. Input~1(d), used for the endomorphism \m{CA_1} with~\eqref{e:L} and then for
\m{\Delta_2^{\tp}}, gives
\[\IndexDvd{\cO^{\mathcal{Y}}}{\Delta_2^{\tp}CA_1\Delta_1(\cO^{\mathcal{Y}})+p^{e}\cO^{\mathcal{Y}}}{p^{\,l'}},\qquad
l'=e_0+\kappa_0\,\#\mathcal{Y}+e\,l,\]
because \m{\varpi^{2e_0}=p^{e_0}}; and \m{\Delta_2^{\tp}CA_1\Delta_1=J_2^{\tp}\,\mathbf{u}\,\mathbf{w}\,J_1}.

\emph{The exponent.} By Theorem~3(b) for~\m{\bar T}, \m{e_0\le(1+z)\,\#\cX\cdot\#\cI}. By \eqref{e:L} and
Theorem~3(a) for~\m{\bar T}, with \m{h=y} and the index \m{\#\cQ} of the parahoric subgroup,
\[l\ \le\ \#\cH\cdot\#\cX\cdot\bigl(\#\cQ\cdot\dim\exc{\bar T}{n}{i_0}{i_1}{0}{y}-(\#\Mon_{k,n})^{L}\bigr)\ \le\ \kappa\cdot\#\cX\cdot\#\cI .\]
As
\m{1+\kappa_0=c/2\le y}, we get \m{l'\le(z+\kappa e+y)\cdot\#\cX\cdot\#\cI}.

\emph{The operator.} Let \m{\mathrm{r}} be the isomorphism from \m{\cO^{\mathcal{Y}}} to
\m{\cO^{\cX\times\cI}} with
\[(\mathrm{r}G)(\fx,\fii)=G\bigl((\fh_{\fii},\fx),\beta_{\fii}\bigr).\]
We claim that, for every $G$ in \m{\cO^{\mathcal{Y}}},
\begin{equation}\label{e:op}
\mathrm{r}\bigl(J_2^{\tp}\,\mathbf{u}\,\mathbf{w}\,J_1\,G\bigr)=\bigl(\mathrm{r}(\mathbf{w}G)\bigr)^{\flat},
\end{equation}
with \m{F\mapsto F^{\flat}} the operator of the statement. By the definitions of \m{J_1} and~\m{J_2}, the
coordinate of \m{J_2^{\tp}\mathbf{u}\mathbf{w}J_1G} at \m{((\fh,\fx),\beta)} is
\begin{multline*}
\sum_{\fq}\sum_{\fh',\beta',\beta''}\varpi^{2\mathrm{d}(\beta')}\,\letterRep{T}{n}{\tilde{\fq}^{-1}}_{\beta',\beta}\,
(u_{(\fx,\fq)})_{\fh,\fh'}\,w(\beta')\\
\cdot\ \letterRep{T}{n}{\gamma_{(\fx,\fq)}^{\tp}}_{\beta',\beta''}\,
G\bigl((\fh',\mathrm{src}(\fx,\fq)),\beta''\bigr),
\end{multline*}
with \m{\mathrm{d}(\beta')} the depth. By Theorem~4(iii),
\m{w(\beta')\,\letterRep{T}{n}{\gamma^{\tp}}_{\beta',\beta''}=w(\beta'')\,\letterRep{T}{n}{\gamma}_{\beta'',\beta'}}; with
\m{\varpi^{2}=p} the sum over~$\beta'$ becomes the \m{(\beta'',\beta)} entry of the matrix \m{H_{(\fx,\fq)}} of the
statement, and \eqref{e:op} follows.

\emph{Conclusion.} By hypothesis \m{Y(\operatorname{im}\Psi)\subseteq\operatorname{im}\Psi}, and
by~\eqref{e:op} the map \m{\Psi\circ\mathrm{r}} sends the image of \m{J_2^{\tp}\mathbf{u}\mathbf{w}J_1} into
\m{Y(\operatorname{im}\Psi)}. Input~1(c), for the map \m{\Psi\circ\mathrm{r}}, carries the index statement
of the chain to \m{\IndexDvd{\operatorname{im}\Psi}{Y(\operatorname{im}\Psi)+p^{e}\operatorname{im}\Psi}{p^{\,l'}}}.
\end{proof}
\leanline{\lean{hecke}\sep \lean{depthfact}\sep \lean{projector}\sep \lean{pairing}\sep \lean{dimex}\sep
\lean{dimdeep}\sep \lean{ofdetdvd}\sep \lean{idcomp}\sep \lean{idmap}}

\sect{Where the results are used}{s:use}
Theorem~5 is used once, in the unit of the columns (unit~145), for the operator attached to a marked prime of
the field of the data. There $W$ is the valuation ring of that field at the prime; the letters are the infinite
places of a certain class (the tops) that induce the prime or its complex conjugate, with the embeddings of
these places, or their composites with complex conjugation, as the maps~$\iota_s$; $\cO$ is the ring of the
complex numbers that are integral for an embedding into a $p$-adic field; $k$ is the rank of the columns,
$n$ the number that the handout calls~\m{a(D)}, and the threshold~$\thr$ the least natural number such that $p$ to
that power is at least the number of lifts of the data; and the module is that of the integral columns. The size
conditions are verified in the same unit. The result enters Theorem~A of the write-up through the index
statement (b) of that theorem: it is the price, in index, of regaining at some primes the divisibility that the
theta lifts lose there. Theorems 1 to~4 and the other results of the unit are used only inside it.

\sect{Notes on the formal text}{s:notes}
\begin{enumerate}[label=\textup{\arabic*.},leftmargin=1.6em,itemsep=5pt]
\item \emph{The three declarations without a name.} The unit has three instances: that the family
\m{\adic{\varpi}{\cdot}} has the three properties of Definition~\ref{d:filt}(a), that
\m{\gauss{\cF}{\omega}{\cdot}} has them when $\cF$ has, and that a
parahoric subgroup has finite index when the residue field is finite. The first two are stated after
Definition~\ref{d:filt}; the third is used silently wherever an index of a parahoric subgroup is counted.
\item \emph{Letters.} The formal text writes $N$ for the number of letters (here~$L$, the letter $N$ being kept
in this write-up for the level), $r$ for the threshold (here~$\thr$), $\ell$ for \m{L/2} (here~$\lambda$), $a$,
$a'$ for the pair of indices (here \m{i_0}, \m{i_1}), $i$ for the base letter (here~$s_0$), $w$ for~$\varpi$,
\m{o_0} for~$o$, $R$, $O$, $K$ for the ring of the data, $g$, $h$, $u$ for group elements (here~$\gamma$), $c$,
$x$ for weights (here~$\omega$), $P$, $v$, $D$, $\Omega$ for \m{\varrho}, \m{\pa}, \m{\delta}, $\cC$ in the
statements on cones, $\nu$, $\beta$ for the points $\xi$, $\eta$ there, $Q$, $M$, $R$, $\beta$ and $a$, $z$, $w$ for
the matrices $A$, $B$, $A'$, $C$ and the elements \m{\pi_1}, \m{\pi_2}, \m{\pi_0} of Proposition~\ref{p:pairing},
$Q$, $d$, $R$, $s$, $g$, $x$ for $A$, $\Delta$, $A'$, \m{\mathrm{pt}}, $\gamma$, \m{j_0} of
Proposition~\ref{p:depthfact}, $b$, $\theta$ for $w$, $h$ of Theorem~4, $\kappa$, $P$, $D$, $f$, $x$ for $\cM$,
$\Omega$, \m{n_0}, $e$, $\omega$ of Proposition~\ref{p:mean}, and $f$, $M$, $J$, $m$, $c$ for $\Psi$, $V$,
$\jmath$, $e$, $\fx$ of Theorem~5.
\item \emph{Arguments that a reading by program adds.} In seventeen declarations of the unit a reading of the
text by program lists as an argument a name that the declaration binds itself, because a variable of the same
name is declared earlier in the file: the group element in the letter matrices, in the action, in the matrix
on forms and in the conditions on forms; the base letter in the excess condition; the letter in the Frobenius
operator and in three lemmas on it; the matrices and the exponent in Proposition~\ref{p:depthfact}; a family in
Theorem~4; the point of~$\cX$ in the formula of Theorem~5; the element~$\varpi$ in Theorem~1, where the
statement mentions only the component of~$T$; the name under which the degree is passed on in the depth; and
the index of the cone of Definition~\ref{d:cones}(d), listed twice. This note states every declaration with the
arguments that its uses in the artifact show: the depth has five, the conditions on forms five and six, the
Frobenius operator five. Lean was not asked.
\item \emph{An argument that the reading by program misses.} In Proposition~\ref{p:signed} the formal statement
names the determinant by a local definition before the equation, and the reading by program stops there; so it
does not see that the statement depends on~$k$. Lean takes $k$ as the first, implicit, argument: the proof of
Theorem~2 applies the proposition with eight arguments, of which $k$ is the first.
\item \emph{Counting an infinite set.} The number of elements of the residue field in
Lemma~\ref{l:stepanov}, Lemma~\ref{l:zariski} and Proposition~\ref{p:mean} is, in the formal text, a natural
number that is~$0$ for an infinite field. For twisting data the field is finite, by the last condition of
Definition~\ref{d:T}.
\item \emph{Hypotheses that are not used.} Setting~\ref{set:T} asks for \m{L\ge1}. Some elementary lemmas of the
unit do not ask for it (parts of Lemma~\ref{l:basic}, Lemma~\ref{l:cjfrob}); nothing is lost, there being no
twisting data for \m{L=0}. Lemma~\ref{l:sumdeep}, Lemma~\ref{l:exdot}, Lemma~\ref{l:conepoint} and
Definition~\ref{d:cones}(d) are stated in the formal text for data over a field, which their proofs do not use.
Of the ten size conditions, Theorem~3(a) uses the third, fourth, eighth and ninth, Theorem~3(b) the first,
second, fourth, fifth, sixth, ninth and tenth, Theorem~4 the fifth and sixth; Theorem~5 uses all.
\item \emph{Arguments with default values.} In Lemma~\ref{l:subst} the second homomorphism is, by default, the
first, and hypotheses (ii) and~(iii) are then discharged by Lean; hypothesis~(iv) of Lemma~\ref{l:frobmem} and the
weights of the Hermite form have defaults as well.
\item \emph{Definitions written out.} The formal text defines the cone of Definition~\ref{d:cones}(c) through
the step of~(a); the two inequalities written there are that definition unfolded. Volumes are taken as real
numbers. In Theorem~5 the formal statement takes an enumeration of the cosets as an argument and the chosen
representative of a coset from the library; the module on which \m{\Psi} is defined is written as a module of
functions of two arguments.
\item \emph{The last line of Theorem~5.} It is the index divisibility of Input~1: for every finite family in
the image of~$\Psi$ a square matrix with determinant dividing \m{p^{l}} moves it into
\m{Y(\operatorname{im}\Psi)+p^{e}\operatorname{im}\Psi}. Nothing is assumed about torsion in~$V$, and the residue
field of~$\cO$, over which the dimensions of Theorem~3 are taken in the proof, is not the residue field
of~$W$.
\item \emph{The constant of Theorem~5.} The proof fixes \m{c=2(k+k^{2})oL\cdot L+2}, the number in the seventh
size condition; an unrelated variable of the same name is in scope in the file.
\item \emph{Assemblies.} Theorem~5 assembles Theorems 3 and~4 and Propositions \ref{p:pairing}
and~\ref{p:depthfact}; Theorem~3(a) assembles Theorem~2, Proposition~\ref{p:dimvol} and Lemmas \ref{l:borel}
to~\ref{l:simplex}. Their own content is the bookkeeping written out in their proofs.
\item \emph{Sources named by the artifact.} The headers of three files of the unit call parts of the argument
known and name related work. This note has not compared any of it; the titles and years are those of the headers.
For the root lemma: counts of zeros against degree, after Stepanov, and E.~Bombieri, \emph{Counting points on
curves over finite fields}, S\'em.\ Bourbaki 430 (1973); Ghate and Jana, \emph{Modular representations of
\m{\mathrm{GL}_2(\mathbb{F}_q)} using calculus} (2025); Imai and Koskivirta, \emph{Automorphic vector bundles on
the stack of $G$-zips} (2021). For the dimension bounds: volumes of interlacing regions, Frobenius amplification
and counting by leading terms along flags; Forrester and Zhang, \emph{Corank-1 projections and the randomised
Horn problem} (2021), after Dixon (1905) and Anderson (1991); Bombieri, as above; Okounkov,
\emph{Brunn--Minkowski inequality for multiplicities} (1996). For the projector: the Schur algebra as the dual of
the polynomials in a generic matrix, and truncation by weights; Benson and Doty, \emph{Schur--Weyl duality over
finite fields} (2009); J.~A.~Green, \emph{Polynomial representations of \m{\mathrm{GL}_n}} (1980); Donkin,
\emph{On Schur algebras and related algebras~I} (1986). The file of Input~3 names the two papers cited after the
Inputs.
\end{enumerate}

\subhead{How this note was written.} From the dossier of the unit: its \UnitMembers{} statement readings and its
structured proofs; the program check of the dossier covered \UnitSteps{} leaf steps of \UnitProofs{} proofs. A
second reader compared all readings with the Lean text and sampled every proof; the corrections of that
reading are followed here: in Proposition~\ref{p:depthfact} the root lemma is applied with the order
\m{\mathrm{ord}(\beta_0)} and the inequality \m{j_0\le\mathrm{ord}(\beta_0)+\mathrm{d}(\beta_0)} is not strict, the
numbers \m{m_y} are counts and not valuations, the first factorisation in the proof of Theorem~5 uses the
transposed elements, the divisibility in the proof of Theorem~4 rests on the condition on coefficients and not
on invariance alone, and Theorem~2 is an identity for every real \m{\varrho>1}. I read the eight files of the
unit completely, and in the background library the definitions of the conventions and the statements of the
Inputs; I did not read the proofs of the Inputs. The places where the dossier's reader was unsure were read in
the source: the seventeen declarations of note~3, the statement of Proposition~\ref{p:signed} with the place
where it is applied, and the long proofs of Theorems 1, 4 and~5 and of Propositions \ref{p:dimvol},
\ref{p:mean}, \ref{p:powdvd}, \ref{p:pairing} and~\ref{p:depthfact}, whose automation was not replayed. Each
inequality that the formal text leaves to automation was checked by hand, among them the constants $0.4$,
$1.6$, \m{\tfrac1{20}} and $5$ of Theorem~1. The proofs follow the formal ones in substance. They differ in
form in these places: Lemma~\ref{l:subst} is proved term by term instead of by induction on the polynomial;
the determinant in Proposition~\ref{p:signed} is computed by an elimination instead of a factorisation and
Sylvester's identity; Lemma~\ref{l:diagform}(ii) is proved by expanding the product instead of by an adjointness
for a pairing of polynomials; in Proposition~\ref{p:deepweight} the determinant is seen not to vanish from the
shape of the matrix instead of by a specialisation; Lemma~\ref{l:degle} is proved directly; in
Proposition~\ref{p:mean} the form is a sum over the degrees instead of the result of an induction; and
Theorem~3(b) averages over cyclic shifts instead of reflections. The diagram, the remarks after the definitions
and the paragraph after Proposition~\ref{p:mean} are mine. Lean was not run for this note.

An independent referee read the note against the Lean text. In the revision of 26~November 2026, after that
report: the threshold and the parity received the letters $\thr$ and~$\pa$ (they were $\mu$ and~$\varsigma$); three
steps were written out, namely the zero test in the proof of Lemma~\ref{l:zariski}, the index of the parahoric
subgroup in the proof of Theorem~3(a), and the units in the proof of Lemma~\ref{l:diagform}; the references for
the Inputs were sorted; and the notation list was completed. For the zero test and for the power of~$p$ in a
binomial coefficient the formal proofs call lemmas of the library Mathlib (a homogeneous polynomial of degree at
most the number of elements of an integral domain that vanishes at every point is zero; the exponent of~$p$ in
\m{\binom{l'}{j}} is at most the logarithm of~$l'$ to the base~$p$). The proofs written here are mine and were not
compared with the library's.

\head{Appendix A. Concordance with the Lean artifact}
{\fontsize{8.6}{11.6}\selectfont\renewcommand{\leanfont}{\ttfamily\fontsize{7.7}{10}\selectfont\color{leancol}}
\setlength{\LTleft}{0pt}\setlength{\LTright}{\fill}\setlength{\tabcolsep}{0pt}
\begin{longtable}{@{}p{26mm}p{66mm}p{36mm}r@{}}
\textbf{In this note} & \textbf{Lean declaration} & \textbf{File} & \textbf{Lines}\\\midrule\endhead
\ConcordanceRows
\end{longtable}}

\subhead{The five theorems and the twisting data as they are written in the artifact.}
\quoted{TwD}
\quoted{rootlemma}
\quoted{prodvol}
\quoted{dimex}
\quoted{dimdeep}
\quoted{projector}
\quoted{hecke}
In these texts {\leanfont ZMod N} is \m{\ZL}, {\leanfont Fin k} the set of the indices, {\leanfont levi W τ} and
{\leanfont parahoric W τ} the two subgroups of the conventions, {\leanfont ⊥} the zero ideal, {\leanfont weight c α}
the sum \m{\pr{\omega}{\alpha}}, and {\leanfont volume.real} the volume as a real number; note~2 of
§\ref{s:notes} gives the letters.

\head{Notation}
\begingroup\fontsize{8.6}{10.9}\selectfont\raggedright
\setlength{\columnsep}{9mm}\setlength{\parskip}{0pt}
\begin{multicols}{2}
\nt{$L$, $s'$, $s_0$, $s_1$, \m{[s-s_0]}}{the number of letters; letters, elements of \m{\ZL}; the representative in \m{\{0,\dots,L-1\}}}
\nt{$s$}{in Definition~\ref{d:T} and after it also a letter, an element of \m{\ZL} (the complex variable of Convention~0.1 does not occur in this note)}
\nt{$k$, $i$, $i'$, $l$, $l'$, $l_1$, \m{i_0}, \m{i_1}}{the number of indices; indices \m{0,\dots,k-1}; a pair of indices ($l$, $l'$ are also natural numbers)}
\nt{$v$, \m{x_v}, \m{x_{i,s}}}{a variable \m{(i,s)}; the variables of the polynomial ring}
\nt{$W$, \m{\fm_W}, \m{\FF_W}}{a local ring, its maximal ideal, its residue field}
\nt{$\cO$, $\cO'$, $\fm$, $\FF$, \m{\FF'}}{rings; the maximal ideal of a valuation domain; fields}
\nt{$T$, $p$, $o$, \m{\iota_s}, \m{\pa_T}, $\varpi$}{twisting data and their components (Definition~\ref{d:T})}
\nt{\m{\TwD{L}{W}{\cO}}, \m{\PT{T}}, \m{\qT{T}}, \m{\Tmap{\chi}{T}}, $\bar T$}{the set of twisting data; \m{p^{o}}; \m{\PT{T}^{L}}; the image under a ring homomorphism; the reduction}
\nt{$\gamma$, \m{\gamma'}, \m{\gamma_0}, \m{\gamma_1}, \m{\gamma_{\upsilon}}, \m{\gamma_{\fe}}, \m{\gamma^{\tp}}}{elements of \m{\GLk(W)}; the transpose}
\nt{\m{\letterMatrix{T}{\gamma}{s}}, \m{\polyAction{T}{\gamma}}, \m{\letterRep{T}{n}{\gamma}}}{the letter matrix, the action on polynomials, the matrix on forms (Definition~\ref{d:action})}
\nt{$\tau$, \m{\twoBlock{k}{\thr}{\cdot}}, \m{\Lev_{\tau}(W)}, \m{\Par_{\tau}(W)}, \m{\Lev^{\circ}}, \m{\Par^{\circ}}}{labellings; Levi and parahoric subgroups}
\nt{$n$, $n'$, \m{n_0}, $\thr$, $\lambda$, $j$, \m{j_0}, $e$, \m{e_0}, \m{e_i}, $c$, $m$, \m{m_y}}{natural numbers: degrees, the threshold of the depth, \m{L/2}, levels and steps, exponents, counts, an index of summation}
\nt{$r$}{in Setting~\ref{set:T} and after it also a natural number: a degree, a count or an index of summation (the rank of the data of Convention~0.3 does not occur in this note)}
\nt{$\alpha$, \m{\alpha'}, \m{\alpha''}, \m{\alpha^{*}}, \m{\alpha^{\pm}_s}, \m{\hat\nu}, \m{\nu_s}}{exponents; values at root variables; exponents on generic variables}
\nt{$\nu$}{in Setting~\ref{set:T} and after it also an exponent, a natural number (the characters of the ideles of Convention~0.6 do not occur in this note)}
\nt{$\beta$, \m{\beta'}, \m{\beta''}, \m{\beta_0}, \m{\Mon_{k,n}}, \m{\Fam}}{families; exponents of degree~$n$; the set of families}
\nt{$F$, \m{F'}, \m{F_1}, \m{F_2}, \m{F_{\nu}}, \m{F^{\circ}}, \m{F^{\flat}}}{forms and vectors of forms; the Hecke operator applied to~$F$}
\nt{$\Phi$, \m{\Phi'}, \m{\Phi_0}, \m{\Phi_v}, \m{\Phi_{\upsilon}}, \m{\Phi^{*}}, $X$}{polynomials; a further variable (also a matrix, in Input~2)}
\nt{$\omega$, \m{\omega_0}, \m{\omega^{\pm}_s}, \m{\pr{\omega}{\alpha}}}{weights; the weighted sum}
\nt{\m{\dpi{T}{\thr}{v}}, \m{\dep{T}{k}{n}{\thr}{\beta}}}{the indicator of deep variables; the depth}
\nt{\m{\rootVar{T}{i_0}{i_1}{s}}, \m{\exw{T}{s_0}{i_0}{i_1}{v}}}{the root variable; the excess weight}
\nt{\m{\trl{T}{n}{\tau}{I}{\mathcal{B}}}, \m{\exc{T}{n}{i_0}{i_1}{I}{h}}, \m{\deepForms{T}{k}{n}{\thr}{j}}}{the modules of forms of Definition~\ref{d:spaces}; $I$ an ideal, $\mathcal{B}$ a set of families, $h$ a real number}
\nt{$\cF$, $\cG$, \m{\cG_{\pm}}, $\theta$, \m{\theta'}, \m{\theta_0}, \m{\adic{\varpi}{\theta}}, \m{\gauss{\cF}{\omega}{\theta}}}{families of subgroups; real levels; the filtrations of Definition~\ref{d:filt}}
\nt{$\chi$, \m{\chi'}, \m{\chi_0}, \m{\chi_1}, \m{\chi_{\upsilon}}, $\cA$}{ring homomorphisms; an algebra}
\nt{$\upsilon$, \m{\upsilon_0}, \m{\upsilon_{\fs}}, $\pi$, \m{\pi'}, \m{\pi_0}, \m{\pi_1}, \m{\pi_2}}{elements of~$W$ and points; elements of~$\cO$}
\nt{$\epsilon$, \m{\varepsilon_y}, \m{\varepsilon'_y}}{a real number; diagonal entries of a Smith form}
\nt{$\pa$, $\varrho$, $\sigma$, \m{\sigma'}, \m{\tau'}, $\delta$, \m{\delta'}, $\cC$}{a map from the letters to the signs (for twisting data, the parity \m{\pa_T}), a real number \m{>1}, permutations, levels of a slab, a cone}
\nt{\m{\permStep{\pa}{\varrho}{\sigma}}, \m{\slab{k}{\delta}}, \m{\pieri{\pa}{\varrho}}, \m{\exCone{T}{i_1}{h}}}{the step, the slab and the cones of Definition~\ref{d:cones}}
\nt{$\xi$, $\eta$, \m{\eta'}, \m{\xi^{\downarrow}}, \m{\sigma^{\circ}}, \m{\sigma^{\bullet}}, \m{\vol}}{real vectors; a rearrangement; two special vectors; Lebesgue measure}
\nt{$q$, $Q$, \m{P'}}{\m{\varrho^{L}} or \m{\qT{T}}; a power of \m{\qT{T}}; a power of \m{\PT{T}}}
\nt{$\kappa$, $y$, $z$, \m{\kappa_0}, \m{\Sz{\mathbf{P},o,L,k,n,\thr}{\lambda,\kappa,y,z}}}{real numbers; the size conditions (Definition~\ref{d:sizes}); $\mathbf{P}$ stands for a natural number there; $y$ also an element of $\mathcal{Y}$ in the proof of Theorem~5}
\nt{\m{X_s}, \m{x_{s,i,i'}}, \m{\frobX{T}{s}{j}}, \m{\frob{T}{\thr}{\lambda}{n}{\psi}}, $\psi$, \m{\Pi^{m}_s}, \m{B_s}}{generic matrices and the Frobenius operator (Definition~\ref{d:frob})}
\nt{\m{\letterConj{T}{\gamma}{s}}, \m{\idxw{T}{\lambda}{i_0}{i_1}{i}}, \m{\conjSubst{\chi}{\gamma}}, \m{G^{(j)}}}{conjugation, index weights, the substitution by conjugation; entrywise powers}
\nt{\m{E_{ii'}}}{the matrix with the entry $1$ at \m{(i,i')} and $0$ elsewhere}
\nt{$G$, \m{G_s}, $S$, \m{S_{\alpha}}, \m{N_s}, \m{A_s}, $U$, \m{U'}, \m{U_1}, \m{U_2}, $\Delta$, $J$, \m{J_1}, \m{J_2}, $H$}{matrices}
\nt{$A$, $B$, \m{A'}, $C$, \m{C_0}, \m{C_1}, \m{A_1}, \m{A_2}, $Z$}{the matrices of §\ref{s:div}; the projector}
\nt{$w$, \m{w_0}, \m{w'}}{in Section~\ref{s:div} also the diagonal form, a matrix (a finite place of Convention~0.4 does not occur in this note)}
\nt{$\cM$, \m{\cM_j}, \m{\cM^{\pm}_s}, $\cL$, \m{\cL'}, \m{\cL_0}, \m{\cL^{\circ}}, $\Omega$}{linear forms on polynomials; a set of real numbers}
\nt{$V$, \m{V'}, \m{V_j}, \m{V_j^{\Omega}}, \m{M_0}, \m{M_1}, \m{M_2}, \m{M_3}, \m{M'_1}, \m{M'_2}, $\Psi$, $Y$}{modules and submodules; a linear map; an endomorphism}
\nt{\m{\IndexDvd{M_1}{M_2}{\pi}}}{the index of \m{M_2} in \m{M_1} divides~$\pi$ (Input~1)}
\nt{$\cX$, $\cX'$, $\cH$, $\cI$, $\cQ$, $\cN$, $\cK$, $\cS$, $\cV$, $\cR$, $\mathcal{Y}$, \m{\mathcal{Y}_0}, $\mathcal{U}$, $\mathcal{W}$, $\mathcal{J}$, $\Gamma$, $\Upsilon$, $\Xi$}{finite sets, sets of indices and of points, a ring ($\cR$), a module ($\Xi$)}
\nt{$\cE$, $\cE'$}{in Theorem~5 and Proposition~\ref{p:depthfact} also finite sets of indices or of points (the polar functions of Convention~0.8 do not occur in this note)}
\nt{$\fx$, $\fq$, \m{\tilde{\fq}}, $\fe$, $\fe'$, $\fh$, $\fii$, $\fs$, \m{y_1}, \m{y_2}}{elements of $\cX$, $\cQ$, $\cE$, $\cH$, $\cI$, $\cS$, $\mathcal{Y}$}
\nt{\m{u_{\fe}}, $\jmath$, \m{\beta_{\fii}}, \m{\fh_{\fii}}, \m{\mathbf{u}}, \m{\mathbf{Z}}, \m{\mathbf{w}}, \m{\mathbf{l}}, \m{\mathbf{v}}, \m{e_y}}{the data of Theorem~5; operators in its proof; a sequence of indices; a vector; unit vectors}
\nt{\m{\mathrm{w}^{\pm}(\alpha')}, \m{\mathrm{w}^{\pm}(v)}, \m{\mathrm{sg}_s}, \m{\mathrm{sg}}, \m{\mathrm{t}_{\upsilon}}}{in the proof of Theorem~1: two weighted degrees of an exponent and their coefficients, a sign (also in Proposition~\ref{p:powdvd}), the translation \m{X\mapsto X+\iota_{s_0}(\upsilon)}}
\nt{\m{\mathrm{dt}(\sigma')}, \m{\sigma'^{*}}, \m{S_{\sigma}}, \m{\mathrm{Sh}}, \m{\mathrm{N}}, \m{\mathrm{El}}}{in Proposition~\ref{p:signed}: \m{|\det(1-\varrho^{-L}\sigma'^{*})|}, a permutation acting on vectors, the step, the shift, the cyclic operator, the elimination}
\nt{\m{\mathrm{f}_l}, \m{\mathrm{df}_s}, \m{\mathrm{g}_i}}{the numbers \m{(q^{l+1}-1)/(q-1)}; differences of coordinates (Lemma~\ref{l:exdot}); the numbers of Lemma~\ref{l:conepoint}}
\nt{\m{\mathrm{cost}(\beta)}, $\varkappa$}{the sum of Proposition~\ref{p:deepweight}; the map of Input~3 that gives the costs}
\nt{\m{\mathrm{rw}(\alpha)}, \m{\mathrm{cw}(\alpha)}, \m{X^{\tau}}, \m{e_{\mathbf{l}}}, \m{|\mathbf{l}|_{\omega}}, \m{\mathrm{K}(\Phi)}}{in Proposition~\ref{p:mean}: the row and column weights of an exponent on the generic variables, the generic matrix of a Levi subgroup, a basis vector and its weight, the vector \m{(w(\alpha)\Phi_{\alpha})_{\alpha}}}
\nt{\m{\mathrm{a}_m(i)}, \m{\eta_{s,j}(i)}, \m{\mathrm{Fr}}, \m{Z(\cL)}, \m{\mathrm{cf}_{\nu}}, \m{\mathrm{bil}(\cL)}}{in §\ref{s:proj}: real numbers, the generic Frobenius operator, its image under a product of forms, a coefficient, the bilinear expression of property~(ii)}
\nt{\m{\mathrm{res}}, \m{\mathrm{src}}, \m{\mathrm{pt}}, \m{\mathrm{pt}_1}, \m{\mathrm{pt}_2}, \m{\mathrm{r}}}{the residue map; maps to~$\cX$ (Theorem~5, Proposition~\ref{p:depthfact}); an isomorphism in the proof of Theorem~5}
\nt{\m{\mathrm{d}(\beta)}, \m{\mathrm{ord}(\beta')}}{in §\ref{s:div}: the depth \m{\dep{T}{k}{n}{i_1}{\beta}}; the order of divisibility defined in the proof of Proposition~\ref{p:depthfact}}
\nt{$k'$, $a'$, $q'$, $r'$, \m{c_1}, \m{c_2}, \m{c_3}, \m{\phi_s}}{letters of the note on unit~145 (after Theorem~5); the letters $\mu$ and $\varsigma$ of the write-up are not used in this note (after Setting~\ref{set:T})}
\nt{$\ell$, $a$, $b$, $g$, $x$, $K$, $R$, $M$}{letters of the formal text, used only in §\ref{s:notes} and after Definition~\ref{d:sizes}}
\nt{$N$, $d$, $f$, $D$, $O$}{in §\ref{s:notes} and after Definition~\ref{d:sizes} also letters of the formal text (the meanings of Conventions~0.1 to~0.7 do not occur in this note); \m{a(D)} is a number of the handout (§\ref{s:use})}
\end{multicols}
\endgroup

\end{document}
